% 农业上市公司画像：种子、粮食种植、其他种植业、食用菌（共 20 家，按流通市值降序）
% 由 scripts/make_company_profiles.py 自动生成，请勿手工修改

\subsubsection{海南橡胶（601118）}
\label{co:601118}

\pfl{公司基本情况} 2005 年设立至今，海南橡胶是全球最大的天然橡胶种植、加工与贸易企业之一；长期处于扣非亏损状态，盈利对胶价与汇率高度敏感，并以举债和海外并购支撑规模扩张。 公司于 2011-01-07 在上交所首发上市，注册资本 42.79 亿元，总股本 42.79 亿股。 截至 2026 年 9 月 24 日收盘价 6.48 元，流通市值 277.3 亿元，近一年+24.1\%，流通市值在三级行业“其他种植业”的 6 家公司中按由大到小排第\zhnumber{1}位，近一年涨跌幅高于全样本中位数（-10.1\%）34.2 个百分点。 按 2025 年年报披露的行业构成，收入占比前两位为农业 96.8\%；其他(补充) 3.2\%。 股权结构方面，控股股东为海南省农垦投资控股集团有限公司（持股 64.35\%）。

\pfl{历史股价} 以 2021 年 1 月为起点、2026 年 9 月为终点，股价由 5.45 元变为 6.48 元，累计 +18.9\%，区间最高 7.68 元（2026 年 2 月）、最低 4.04 元（2022 年 10 月）；当前价相当于区间最高价的 84.4\%，为区间最低价的 1.60 倍。 月度收益的年化波动率 29.2\%，低于全样本中位数 37.4\%，在三级行业“其他种植业”的 6 家公司中波动率由高到低排第\zhnumber{5}位。 把股价阶段与公司事件对齐看，2022 年 10 月 至 2026 年 2 月股价上涨+90.1\%，同期（2022 年 12 月）公司将所持北京海垦商贸发展有限公司 29.34\% 股权全部转让给海垦控股集团。 从时间对应关系看，公司层面的资本运作与股权、风险事件解释了区间内多数急涨急跌，与行业指数的同步性相对有限。

\begin{center}
\begin{minipage}{\textwidth}\centering
\begin{tikzpicture}
\begin{axis}[
  width=12.6cm, height=3.9cm, scale only axis,
  xmin=2020.922, xmax=2026.888, ymin=3.76, ymax=8.22,
  xtick={2021,2022,2023,2024,2025,2026},
  yearticks,
  yticklabel style={font=\scriptsize},
  ylabel={元/股}, ylabel style={font=\scriptsize},
  grid=major, grid style={LightGray, line width=0.3pt},
  axis line style={InkGray, line width=0.5pt},
  every axis plot/.append style={line width=0.9pt},
]
\addplot[AgriGreen] table[x=x,y=y]{data/clean/profiles/px_601118.dat};
\addplot[SoftRed, dashed, line width=0.7pt] table[x=x,y=y]{data/clean/profiles/pxzz_601118.dat};
\addplot[only marks, mark=*, mark size=1.1pt, SoftRed] table[x=x,y=y]{data/clean/profiles/pxzz_601118.dat};
\end{axis}
\end{tikzpicture}

\captionof{figure}[海南橡胶（601118）股价与周期骨架]{海南橡胶（601118）月度收盘价（元/股）与 ZigZag 周期骨架（阈值 25\%，圆点为周期转折点）}
\end{minipage}
\end{center}

\pfl{过去周期分析} 以 25\% 的 ZigZag 阈值划分，样本区间内股价形成 2 个主升段与 3 个主跌段。 最大主升段出现在 2022 年 10 月 至 2026 年 2 月，40 个月+90.1\%；最大主跌段为 2026 年 2 月 至 2026 年 6 月，4 个月-35.9\%。 最近一次转折出现在 2026 年 9 月（上涨段自 2026 年 6 月起，3 个月+31.7\%），当前处于该段之后的震荡之中。 公司股价较申万农林牧渔指数提前约 43 个月见底，是板块中较早定价周期反转的公司之一。 该子行业的周期来源与猪周期不同（第 \ref{sec:grain} 章、第 \ref{sec:soft} 章），价格更多由自身供给、进口政策与全球定价决定，公司 2026 年 6 月末资产负债率 70.7\%，高于全样本中位数 52.2\%，其股价因此更多反映资金面与信用风险，而不是价格弹性本身。 综合区间形态看，该股单段涨跌幅度偏大、以趋势段为主（最大主升 +90.1\%、最大主跌 -35.9\%），年化波动率 29.2\% 与全样本中位数 37.4\% 相比波动小于样本中位数。

\pfl{盈利情况} 2026 年上半年营业收入 151.23 亿元（同比 -25.4\%），归母净利润 -1.30 亿元，亏损同比收窄 0.46 亿元。 以年报为趋势对照，2023---2025 年营业收入由 334.67 亿元变为 422.37 亿元（两年年均复合增速 +12.3\%），归母净利润由 2.97 亿元变为 -1.03 亿元。 按半年报口径，2026 年上半年的 ROE 为 -1.3\%（未年化）、销售净利率 -1.8\%、毛利率 3.4\%，ROE 在“其他种植业”6 家公司中按数值由高到低排第\zhnumber{6}位。 按同一口径，三级行业“其他种植业”6 家公司的半年 ROE 中位数为 1.70\%，该公司低于其中位数 3.00 个百分点。 从公司披露的原因看，业绩变动主要来自：2024 年，公司亏损主因贷款规模较高、外币借款利息与汇兑损失合计约 5.17 亿元；2026 年 6 月，公司亏损主要受美元汇率下跌导致汇兑损失增加、金融工具平仓损失、合同纠纷判决款等偶发性支出及递延所得税费用增加影响。 公司披露的应对举措集中在：2026 年，公司推进老挝锦森二期项目落地、柬埔寨公司设立，扩大橡胶加工产能。 综合看，公司仍处亏损区间、亏损同比收窄，半年 ROE 在三级行业内按由高到低排第\zhnumber{6}位，单位盈利能力的修复依赖行业景气与自身成本端的改善，报表弹性主要体现为价格与销量的方向性变化。

\begin{center}\small\setlength{\tabcolsep}{3.4pt}
\captionof{table}[海南橡胶（601118）关键财务指标]{海南橡胶（601118）关键财务指标（合并报表口径；两个半年报期均未年化；现金短债比 $=$ 货币资金 $\div$ 短期债务）}
\begin{adjustbox}{max width=\textwidth}
\begin{tabular}{@{}lrrrrr@{}}
\toprule
指标 & 2023 年 & 2024 年 & 2025 年 & 2025 年半年报 & 2026 年半年报 \\
\midrule
营业收入（亿元） & 334.67 & 420.04 & 422.37 & 202.72 & 151.23 \\
归母净利润（亿元） & 2.97 & 1.03 & -1.03 & -1.76 & -1.30 \\
毛利率（\%） & 2.4 & 3.8 & 2.5 & 3.8 & 3.4 \\
ROE（\%） & 3.1 & 1.1 & -1.0 & -1.8 & -1.3 \\
资产负债率（\%） & 65.7 & 68.7 & 70.3 & 69.0 & 70.7 \\
经营活动现金流净额（亿元） & 12.05 & 15.55 & 25.05 & 20.90 & 10.06 \\
货币资金（亿元） & 35.06 & 56.88 & 52.69 & 72.76 & 51.31 \\
有息负债合计（亿元） & 165.56 & 177.92 & 172.30 & 182.33 & 167.80 \\
现金短债比（倍） & 0.33 & 0.55 & 0.50 & 0.66 & 0.55 \\
\bottomrule
\end{tabular}
\end{adjustbox}
\end{center}

\begin{center}
\begin{minipage}{\textwidth}\centering
\begin{tikzpicture}
\begin{groupplot}[
  group style={group size=2 by 1, horizontal sep=0.60cm},
  width=6.85cm, height=4.60cm,
  tick label style={font=\tiny},
  title style={font=\scriptsize\heiti},
  yticklabel style={font=\tiny},
  ylabel style={font=\tiny},
  grid=major, grid style={LightGray, line width=0.3pt},
  axis line style={InkGray, line width=0.5pt},
  enlarge x limits={abs=0.8},
  legend style={font=\scriptsize, at={(0.5,-0.20)}, anchor=north,
                legend columns=2, draw=none, fill=none, column sep=6pt},
]
\nextgroupplot[
  ybar, bar width=4pt,
  ymin=-43.67, ymax=473.21,
  xtick={0,1,2,3,4,5.7},
  xticklabels={2021,2022,2023,2024,2025,2026H1},
  ylabel={亿元},
]
\addplot[fill=AgriGreen!75, draw=AgriGreen, bar shift=-2.6pt] table[x=i, y=rev]{data/clean/profiles/fin_601118.dat};
\addplot[fill=SoftRed!60, draw=SoftRed, bar shift=2.6pt] table[x=i, y=ni]{data/clean/profiles/fin_601118.dat};
\addplot[fill=AgriGreen!30, draw=AgriGreen, pattern=north east lines, bar shift=-2.6pt] table[x=x, y=rev]{data/clean/profiles/finh_601118.dat};
\addplot[fill=SoftRed!20, draw=SoftRed, pattern=north east lines, bar shift=2.6pt] table[x=x, y=ni]{data/clean/profiles/finh_601118.dat};
\legend{营业收入（年/半年）, 归母净利润（年/半年）}
\nextgroupplot[
  ymin=-7.1, ymax=77.9,
  xtick={0,1,2,3,4,5.7},
  xticklabels={2021,2022,2023,2024,2025,2026H1},
  ylabel={\%},
]
\addplot[AgriGold, mark=*, mark size=1.3pt] table[x=i, y=roe]{data/clean/profiles/fin_601118.dat};
\addplot[OilBlue, mark=square*, mark size=1.3pt] table[x=i, y=debt]{data/clean/profiles/fin_601118.dat};
\addplot[only marks, AgriGold, mark=*, mark size=2.0pt] table[x=x, y=roe]{data/clean/profiles/finh_601118.dat};
\addplot[only marks, OilBlue, mark=square*, mark size=2.0pt] table[x=x, y=debt]{data/clean/profiles/finh_601118.dat};
\legend{ROE, 资产负债率}
\end{groupplot}
\end{tikzpicture}
\begin{tikzpicture}
\begin{axis}[
  name=finbot,
  width=13.75cm, height=3.10cm, scale only axis,
  ymin=-17.79, ymax=199.27,
  xtick={0,1,2,3,4,5.7},
  xticklabels={2021,2022,2023,2024,2025,2026H1},
  tick label style={font=\tiny},
  yticklabel style={font=\tiny},
  ylabel={亿元}, ylabel style={font=\tiny},
  ybar, bar width=4pt,
  grid=major, grid style={LightGray, line width=0.3pt},
  axis line style={InkGray, line width=0.5pt},
  enlarge x limits={abs=0.8},
  legend style={font=\scriptsize, at={(0.5,-0.26)}, anchor=north,
                legend columns=2, draw=none, fill=none, column sep=10pt},
]
\addplot[fill=AgriGreen!55, draw=AgriGreen, bar shift=-3.0pt] table[x=i, y=cash]{data/clean/profiles/fin_601118.dat};
\addplot[fill=SoftRed!45, draw=SoftRed, bar shift=3.0pt] table[x=i, y=ibd]{data/clean/profiles/fin_601118.dat};
\addplot[fill=AgriGreen!25, draw=AgriGreen, pattern=north east lines, bar shift=-3.0pt] table[x=x, y=cash]{data/clean/profiles/finh_601118.dat};
\addplot[fill=SoftRed!20, draw=SoftRed, pattern=north east lines, bar shift=3.0pt] table[x=x, y=ibd]{data/clean/profiles/finh_601118.dat};
\legend{货币资金（左轴）, 有息负债（左轴）}
\end{axis}
\begin{axis}[
  at={(finbot.south east)}, anchor=south east,
  width=13.75cm, height=3.10cm, scale only axis,
  axis y line*=right, axis x line*=none, xtick=\empty,
  ymin=-6.51, ymax=32.56,
  tick label style={font=\tiny},
  yticklabel style={font=\tiny}, scaled y ticks=false,
  legend style={font=\scriptsize, at={(0.5,-0.44)}, anchor=north,
                legend columns=1, draw=none, fill=none},
]
\addplot[OilBlue, mark=*, mark size=2.2pt, line width=1.3pt] table[x=i, y=ocf]{data/clean/profiles/fin_601118.dat};
\addplot[only marks, OilBlue, mark=*, mark size=3.2pt] table[x=x, y=ocf]{data/clean/profiles/finh_601118.dat};
\legend{经营活动现金流净额（右轴，亿元，蓝点）}
\end{axis}
\end{tikzpicture}

\captionof{figure}[海南橡胶（601118）近年主要财务指标]{海南橡胶（601118）近年主要财务指标（左上：营业收入与归母净利润；右上：ROE 与资产负债率；下：货币资金、有息负债为左轴柱状，经营活动现金流净额为其右轴的蓝点折线。
2021---2025 年为年报口径，2026H1 为 2026 年半年报/半年末，与其前一年报之间留出间隔、以斜纹或空心点区分；半年报数据未年化）}
\end{minipage}
\end{center}

\pfl{财务分析} 2026 年 6 月末资产负债率 70.7\%，较 2025 年末上升 0.4 个百分点（样本中位数 52.2\%，所处三级行业内按由低到高排第\zhnumber{5}位）。 2026 年 6 月末货币资金 51.31 亿元、短期债务（短期借款与一年内到期非流动负债）93.23 亿元、长期债务 74.57 亿元，有息负债合计 167.80 亿元，现金短债比 0.55 倍——覆盖偏紧。 净有息负债 116.49 亿元，相当于其 2026 年上半年营业收入的 77\%。 流动比率 0.98、速动比率 0.70，短期偿付压力较大。 2026 年上半年经营活动现金流净额 10.06 亿元，净利润为负而经营现金流为正，差额主要来自折旧摊销与营运资本变动，说明亏损尚未完全转化为现金流出。 2026 年 6 月末在建工程 2.62 亿元，较上年末减少 0.11 亿元，在建投入在收缩。 2026 年 6 月末商誉 14.81 亿元，占归母净资产的 15.5\%，是净资产中最脆弱的一块。 综合看，现金短债比 0.55 倍，短期资金安排偏紧；资产负债率高于样本中位数 18.5 个百分点；有息负债中短期债务占比更高，期限结构仍以滚动续作为主。

\pfl{重大战略转型} 公司的战略动作可归为以下几类：并购与重组（2 项）：2022 年 12 月，公司将所持北京海垦商贸发展有限公司 29.34\% 股权全部转让给海垦控股集团；2022 年 11 月，要约收购结束后合计持有 10.86 亿股（约 68.10\%）。 公告口径上，近三年（2023 年 9 月---2026 年 9 月）公司披露重大重组与并购类公告 2 条、对外投资与转型类公告 3 条、回购与股权激励类公告 0 条，合计公告 440 条。 第一大行业仍是“农业”，收入占比 96.8\%，与 2021 年年报相比没有方向性变化。 综合判断，报告期内公司有 5 条与战略相关的公告落地（重大重组与并购 2 条、对外投资与转型 3 条），资本运作与主业调整之间存在直接关联。

\begin{center}\small\setlength{\tabcolsep}{4pt}
\captionof{table}[海南橡胶（601118）历史融资与资本运作]{海南橡胶（601118）历史融资与资本运作（据公司公告与公开报道整理；金额为披露值，含首次公开发行、再融资、债券、银行借款与重整投资等）}
\begin{adjustbox}{max width=\textwidth}
\begin{tabular}{@{}p{1.45cm}p{2.55cm}p{4.05cm}p{4.95cm}@{}}
\toprule
时间 & 方式 & 规模 & 用途与说明 \\
\midrule
2010-12 & IPO（上海证券交易所主板） & 发行 78，\allowbreak{}600 万股，\allowbreak{}发行价 5.99 元/\allowbreak{}股，\allowbreak{}募集资金总额 47.08 亿元 & 招股公告日 2010-12-16；预计募集资金 47.08 亿元 \\
2018-02 & 定向增发（非公开发行） & 发行 34，\allowbreak{}825.62 万股，\allowbreak{}发行价 5.16 元/\allowbreak{}股，\allowbreak{}募集资金总额 17.97 亿元 & 特种胶园更新种植等项目；发行对象共 7 家：华商基金、\allowbreak{}长城（天津）股权投资基金管理、\allowbreak{}太平洋资产管理、\allowbreak{}广发证券资产管理（广东）、\allowbreak{}信泰人寿保险、\allowbreak{}财通基金、\allowbreak{}国信国投基金管理（北京）；2018-02-12 完成新增股份登记，\allowbreak{}942.78 万股 \\
2025-07 & 中期票据（银行间市场） & “25 海南橡胶 MTN001（科创债）”，\allowbreak{}发行金额 3.50 亿元，\allowbreak{}票面利率 2.07\% & 偿还公司本部及下属企业有息负债；债券代码 102583120 \\
2026-04 & 2026 年度融资额度议案 & 公司及子公司拟通过金融机构借款、\allowbreak{}发行债券、\allowbreak{}银行票据等方式融资，\allowbreak{}融资余额较上年度净增加不超过 47.00 亿元 & 债务性和权益性融资；第六届董事会相关会议审议通过 \\
2026-06 & 银行授信（关联方） & 海南农垦集团财务有限公司授予公司授信总额 60.00 亿元，\allowbreak{}本期实际发生额 4.37 亿元 & 关联方授信业务；截至 2026 年 3 月末，\allowbreak{}公司合并口径共获得金融机构授信总额合计 362.74 亿元 \\
\bottomrule
\end{tabular}
\end{adjustbox}
\end{center}

\pfl{历史融投资情况} 从现金流量表看，2025 年公司取得借款收到的现金 191.94 亿元、偿还债务支付的现金 196.40 亿元、吸收权益性投资 3.51 亿元，筹资活动现金流净额 -24.82 亿元（2023 年为取得借款 150.10 亿元、偿还债务 151.41 亿元）。 2026 年上半年筹资活动现金流净额为 -8.55 亿元（取得借款 68.79 亿元、偿还债务 72.59 亿元，未年化），已转为净流出，处于净还债状态。 2025 年分配股利、利润或偿付利息支付的现金合计 7.07 亿元。 公开披露的融资记录共 5 笔，工具类型以 IPO、定向增发、中期票据、银行授信为主；金额与用途见下表。 其中规模最大的两笔为：2026 年 6 月，银行授信，海南农垦集团财务有限公司授予公司授信总额 60.00 亿元，本期实际发生额 4.37 亿元；2025 年 7 月，中期票据，“25 海南橡胶 MTN001（科创债）”，发行金额 3.50 亿元，票面利率 2.07\%。 股本方面，近三年披露股本变动 3 次；总股本自 2021 年末以来未发生变化（42.79 亿股）。 分红方面，近三年现金分红 2 次、合计每股派现 0.03 元，最近一次实施在 2024 年 12 月。 近三年“再融资”类公告 2 条，涉及定向增发。 综合看，公司融资结构上以债务滚动为主、股权融资为补充；2025 年筹资活动现金流净额 -24.82 亿元，融资节奏仍以维持存量债务周转为主，与前述经营与投资现金流的方向基本一致。

\pfl{未来潜在融资需求分析} 2026 年 6 月末货币资金 51.31 亿元，而短期债务 93.23 亿元（现金短债比 0.55 倍），2026 年上半年经营现金流净额 10.06 亿元。按“短期债务减去账面货币资金”的滚动偿付口径，静态缺口约 41.92 亿元；若再计入上半年亏损的年化额（2.61 亿元），维持一年正常经营所需的外部资金约 41.92---44.53 亿元。 按第 \ref{sec:financing-need} 节的三档划分，公司属于第一档“补血型”：2026 年上半年归母净利润 -1.30 亿元、6 月末资产负债率 70.7\%。 可行的融资路径按现实性排序为：出售非核心资产与参股股权、定向增发（需股价配合，且控股股东需放弃优先认购或同步增资）、引入国资或产业投资人（部分公司以控制权让渡为对价）。 公开披露的后续资金安排线索包括：债务与授信安排：2026 年 4 月，公司及子公司融资余额较上年度净增加不超过 47.00 亿元；资本开支安排：2026 年 6 月，特种胶园更新种植项目已累计投入募集资金 12.01 亿元。 资金需求还要叠加在建工程：期末在建工程 2.62 亿元、2026 年上半年购建固定资产等支出 1.50 亿元（未年化），这部分支出在周期底部通常会被主动放缓。 综合看，公司的外部资金需求以债务接续为主，滚动偿付口径下的静态缺口约 41.92 亿元，融资节奏将受制于银行续贷意愿与股权融资窗口。




\subsubsection{北大荒（600598）}
\label{co:600598}

\pfl{公司基本情况} 北大荒成立于 1998 年，黑龙江垦区核心农业上市公司，以国有耕地承包经营（土地发包）为主业，土地承包费收入长期占营业收入七成左右；2026 年上半年因补缴 2021—2025 年企业所得税及滞纳金 14.10 亿元而由盈转亏。 公司于 2002-03-29 在上交所首发上市，注册资本 17.78 亿元，总股本 17.78 亿股。 截至 2026 年 9 月 24 日收盘价 12.49 元，流通市值 222.0 亿元，近一年-13.1\%，流通市值在三级行业“粮食种植”的 2 家公司中按由大到小排第\zhnumber{1}位，近一年涨跌幅低于全样本中位数（-10.1\%）3.1 个百分点。 按 2026 年半年报披露的产品构成，收入占比前三位为土地承包费 91.3\%；农产品销售 5.2\%；其他商品销售或服务 2.3\%。

\pfl{历史股价} 以 2021 年 1 月为起点、2026 年 9 月为终点，股价由 16.37 元变为 12.49 元，累计 -23.7\%，区间最高 17.50 元（2026 年 1 月）、最低 10.84 元（2026 年 6 月）；当前价相当于区间最高价的 71.4\%，为区间最低价的 1.15 倍。 月度收益的年化波动率 25.1\%，低于全样本中位数 37.4\%，在三级行业“粮食种植”的 2 家公司中波动率由高到低排第\zhnumber{2}位。 把股价阶段与公司事件对齐看，2026 年 1 月 至 2026 年 6 月股价下跌-38.1\%，同期（2026 年 6 月）不属于前期会计差错，预计减少 2026 年归母净利润约 14.10 亿元系税收政策理解偏差。 整体上看，区间的几处大级别拐点都能在公司的资本运作、股权变动或风险事件上找到对应，股价波动有明显的公司层面驱动，而非单纯的行业 beta。

\begin{center}
\begin{minipage}{\textwidth}\centering
\begin{tikzpicture}
\begin{axis}[
  width=12.6cm, height=3.9cm, scale only axis,
  xmin=2020.922, xmax=2026.888, ymin=10.08, ymax=18.73,
  xtick={2021,2022,2023,2024,2025,2026},
  yearticks,
  yticklabel style={font=\scriptsize},
  ylabel={元/股}, ylabel style={font=\scriptsize},
  grid=major, grid style={LightGray, line width=0.3pt},
  axis line style={InkGray, line width=0.5pt},
  every axis plot/.append style={line width=0.9pt},
]
\addplot[AgriGreen] table[x=x,y=y]{data/clean/profiles/px_600598.dat};
\addplot[SoftRed, dashed, line width=0.7pt] table[x=x,y=y]{data/clean/profiles/pxzz_600598.dat};
\addplot[only marks, mark=*, mark size=1.1pt, SoftRed] table[x=x,y=y]{data/clean/profiles/pxzz_600598.dat};
\end{axis}
\end{tikzpicture}

\captionof{figure}[北大荒（600598）股价与周期骨架]{北大荒（600598）月度收盘价（元/股）与 ZigZag 周期骨架（阈值 25\%，圆点为周期转折点）}
\end{minipage}
\end{center}

\pfl{过去周期分析} 以 25\% 的 ZigZag 阈值划分，样本区间内股价形成 2 个主升段与 3 个主跌段。 最大主跌段为 2026 年 1 月 至 2026 年 6 月，5 个月-38.1\%；最大主升段出现在 2024 年 1 月 至 2026 年 1 月，24 个月+53.8\%。 最近一次转折出现在 2026 年 6 月（下跌段自 2026 年 1 月起，5 个月-38.1\%），当前处于该段的延续之中。 公司股价与申万农林牧渔指数几乎同步见底，个股并未走出独立行情。 土地经营型企业的收入与地租成本高度相关（第 \ref{sec:planting} 节），粮价与政策补贴是其主要变量，公司 2026 年 6 月末资产负债率 39.5\%、半年 ROE -7.1\%（未年化），在本轮下行中属于随行业同步调整的一类。 综合区间形态看，该股单段涨跌幅度偏大、以趋势段为主（最大主升 +53.8\%、最大主跌 -38.1\%），年化波动率 25.1\% 与全样本中位数 37.4\% 相比波动小于样本中位数。

\pfl{盈利情况} 2026 年上半年营业收入 21.28 亿元（同比 -29.4\%），归母净利润 -5.37 亿元，由上年同期的盈利转为亏损。 以年报为趋势对照，2023---2025 年营业收入由 50.44 亿元变为 52.29 亿元（两年年均复合增速 +1.8\%），归母净利润由 10.64 亿元变为 11.66 亿元。 按半年报口径，2026 年上半年的 ROE 为 -7.1\%（未年化）、销售净利率 -25.4\%、毛利率 54.9\%，ROE 在“粮食种植”2 家公司中按数值由高到低排第\zhnumber{2}位。 同期全样本半年 ROE 中位数 0.75\%，公司与之相差 7.84 个百分点（弱于中位数公司）。 从公司披露的原因看，业绩变动主要来自：2026 年上半年，导致销售收入同比减少 8.86 亿元；2026 年 6 月，不属于前期会计差错，预计减少 2026 年归母净利润约 14.10 亿元系税收政策理解偏差。 综合看，公司仍处亏损区间、由盈转亏，半年 ROE 在三级行业内按由高到低排第\zhnumber{2}位，单位盈利能力的修复依赖行业景气与自身成本端的改善，报表弹性主要体现为价格与销量的方向性变化。

\begin{center}\small\setlength{\tabcolsep}{3.4pt}
\captionof{table}[北大荒（600598）关键财务指标]{北大荒（600598）关键财务指标（合并报表口径；两个半年报期均未年化；现金短债比 $=$ 货币资金 $\div$ 短期债务）}
\begin{adjustbox}{max width=\textwidth}
\begin{tabular}{@{}lrrrrr@{}}
\toprule
指标 & 2023 年 & 2024 年 & 2025 年 & 2025 年半年报 & 2026 年半年报 \\
\midrule
营业收入（亿元） & 50.44 & 53.39 & 52.29 & 30.14 & 21.28 \\
归母净利润（亿元） & 10.64 & 10.87 & 11.66 & 9.84 & -5.37 \\
毛利率（\%） & 35.4 & 34.6 & 34.6 & 39.0 & 54.9 \\
ROE（\%） & 14.6 & 14.3 & 14.9 & 12.1 & -7.1 \\
资产负债率（\%） & 16.1 & 15.4 & 12.9 & 34.1 & 39.5 \\
经营活动现金流净额（亿元） & 15.15 & 14.94 & 12.40 & 36.92 & 24.86 \\
货币资金（亿元） & 23.70 & 15.86 & 2.00 & 7.80 & 19.27 \\
有息负债合计（亿元） & 0.57 & 0.73 & 0.10 & 0.16 & 0.10 \\
现金短债比（倍） & 69.01 & 24.33 & 56.25 & 112.61 & 566.76 \\
\bottomrule
\end{tabular}
\end{adjustbox}
\end{center}

\begin{center}
\begin{minipage}{\textwidth}\centering
\begin{tikzpicture}
\begin{groupplot}[
  group style={group size=2 by 1, horizontal sep=0.60cm},
  width=6.85cm, height=4.60cm,
  tick label style={font=\tiny},
  title style={font=\scriptsize\heiti},
  yticklabel style={font=\tiny},
  ylabel style={font=\tiny},
  grid=major, grid style={LightGray, line width=0.3pt},
  axis line style={InkGray, line width=0.5pt},
  enlarge x limits={abs=0.8},
  legend style={font=\scriptsize, at={(0.5,-0.20)}, anchor=north,
                legend columns=2, draw=none, fill=none, column sep=6pt},
]
\nextgroupplot[
  ybar, bar width=4pt,
  ymin=-11.24, ymax=60.44,
  xtick={0,1,2,3,4,5.7},
  xticklabels={2021,2022,2023,2024,2025,2026H1},
  ylabel={亿元},
]
\addplot[fill=AgriGreen!75, draw=AgriGreen, bar shift=-2.6pt] table[x=i, y=rev]{data/clean/profiles/fin_600598.dat};
\addplot[fill=SoftRed!60, draw=SoftRed, bar shift=2.6pt] table[x=i, y=ni]{data/clean/profiles/fin_600598.dat};
\addplot[fill=AgriGreen!30, draw=AgriGreen, pattern=north east lines, bar shift=-2.6pt] table[x=x, y=rev]{data/clean/profiles/finh_600598.dat};
\addplot[fill=SoftRed!20, draw=SoftRed, pattern=north east lines, bar shift=2.6pt] table[x=x, y=ni]{data/clean/profiles/finh_600598.dat};
\legend{营业收入（年/半年）, 归母净利润（年/半年）}
\nextgroupplot[
  ymin=-10.8, ymax=44.2,
  xtick={0,1,2,3,4,5.7},
  xticklabels={2021,2022,2023,2024,2025,2026H1},
  ylabel={\%},
]
\addplot[AgriGold, mark=*, mark size=1.3pt] table[x=i, y=roe]{data/clean/profiles/fin_600598.dat};
\addplot[OilBlue, mark=square*, mark size=1.3pt] table[x=i, y=debt]{data/clean/profiles/fin_600598.dat};
\addplot[only marks, AgriGold, mark=*, mark size=2.0pt] table[x=x, y=roe]{data/clean/profiles/finh_600598.dat};
\addplot[only marks, OilBlue, mark=square*, mark size=2.0pt] table[x=x, y=debt]{data/clean/profiles/finh_600598.dat};
\legend{ROE, 资产负债率}
\end{groupplot}
\end{tikzpicture}
\begin{tikzpicture}
\begin{axis}[
  name=finbot,
  width=13.75cm, height=3.10cm, scale only axis,
  ymin=-3.05, ymax=34.20,
  xtick={0,1,2,3,4,5.7},
  xticklabels={2021,2022,2023,2024,2025,2026H1},
  tick label style={font=\tiny},
  yticklabel style={font=\tiny},
  ylabel={亿元}, ylabel style={font=\tiny},
  ybar, bar width=4pt,
  grid=major, grid style={LightGray, line width=0.3pt},
  axis line style={InkGray, line width=0.5pt},
  enlarge x limits={abs=0.8},
  legend style={font=\scriptsize, at={(0.5,-0.26)}, anchor=north,
                legend columns=2, draw=none, fill=none, column sep=10pt},
]
\addplot[fill=AgriGreen!55, draw=AgriGreen, bar shift=-3.0pt] table[x=i, y=cash]{data/clean/profiles/fin_600598.dat};
\addplot[fill=SoftRed!45, draw=SoftRed, bar shift=3.0pt] table[x=i, y=ibd]{data/clean/profiles/fin_600598.dat};
\addplot[fill=AgriGreen!25, draw=AgriGreen, pattern=north east lines, bar shift=-3.0pt] table[x=x, y=cash]{data/clean/profiles/finh_600598.dat};
\addplot[fill=SoftRed!20, draw=SoftRed, pattern=north east lines, bar shift=3.0pt] table[x=x, y=ibd]{data/clean/profiles/finh_600598.dat};
\legend{货币资金（左轴）, 有息负债（左轴）}
\end{axis}
\begin{axis}[
  at={(finbot.south east)}, anchor=south east,
  width=13.75cm, height=3.10cm, scale only axis,
  axis y line*=right, axis x line*=none, xtick=\empty,
  ymin=-6.46, ymax=32.31,
  tick label style={font=\tiny},
  yticklabel style={font=\tiny}, scaled y ticks=false,
  legend style={font=\scriptsize, at={(0.5,-0.44)}, anchor=north,
                legend columns=1, draw=none, fill=none},
]
\addplot[OilBlue, mark=*, mark size=2.2pt, line width=1.3pt] table[x=i, y=ocf]{data/clean/profiles/fin_600598.dat};
\addplot[only marks, OilBlue, mark=*, mark size=3.2pt] table[x=x, y=ocf]{data/clean/profiles/finh_600598.dat};
\legend{经营活动现金流净额（右轴，亿元，蓝点）}
\end{axis}
\end{tikzpicture}

\captionof{figure}[北大荒（600598）近年主要财务指标]{北大荒（600598）近年主要财务指标（左上：营业收入与归母净利润；右上：ROE 与资产负债率；下：货币资金、有息负债为左轴柱状，经营活动现金流净额为其右轴的蓝点折线。
2021---2025 年为年报口径，2026H1 为 2026 年半年报/半年末，与其前一年报之间留出间隔、以斜纹或空心点区分；半年报数据未年化）}
\end{minipage}
\end{center}

\pfl{财务分析} 2026 年 6 月末资产负债率 39.5\%，较 2025 年末上升 26.6 个百分点（样本中位数 52.2\%，所处三级行业内按由低到高排第\zhnumber{1}位）。 2026 年 6 月末货币资金 19.27 亿元、短期债务（短期借款与一年内到期非流动负债）0.03 亿元、长期债务 0.06 亿元，有息负债合计 0.10 亿元，现金短债比 566.76 倍——覆盖充分。 净有息负债 -19.17 亿元，相当于其 2026 年上半年营业收入的 -90\%。 流动比率 1.46、速动比率 1.41，短期指标尚可。 2026 年上半年经营活动现金流净额 24.86 亿元，净利润为负而经营现金流为正，差额主要来自折旧摊销与营运资本变动，说明亏损尚未完全转化为现金流出。 综合看，现金短债比 566.76 倍，短期偿付压力可控；资产负债率低于样本中位数 12.7 个百分点；有息负债中长期债务占比更高，期限结构对短债滚动的压力较小。

\pfl{重大战略转型} 公司的战略动作可归为以下几类：并购与重组（1 项）：2014 年，公司剥离北大荒米业集团股权；其他动作（1 项）：2026 年 6 月，报告期末“合同负债”比上期期末增长 25，848.03\%。 公告口径上，近三年（2023 年 9 月---2026 年 9 月）公司披露重大重组与并购类公告 1 条、对外投资与转型类公告 1 条、回购与股权激励类公告 0 条，合计公告 254 条。 第一大产品“土地承包费”的收入占比由 2021 年年报的 77.0\% 变为 91.3\%（14.3 个百分点），收入结构出现再平衡。 综合判断，报告期内公司有 2 条与战略相关的公告落地（重大重组与并购 1 条、对外投资与转型 1 条），资本运作与主业调整之间存在直接关联。

\begin{center}\small\setlength{\tabcolsep}{4pt}
\captionof{table}[北大荒（600598）历史融资与资本运作]{北大荒（600598）历史融资与资本运作（据公司公告与公开报道整理；金额为披露值，含首次公开发行、再融资、债券、银行借款与重整投资等）}
\begin{adjustbox}{max width=\textwidth}
\begin{tabular}{@{}p{1.45cm}p{2.55cm}p{4.05cm}p{4.95cm}@{}}
\toprule
时间 & 方式 & 规模 & 用途与说明 \\
\midrule
2002-03 & IPO（上海证券交易所） & 发行 30，\allowbreak{}000 万股，\allowbreak{}发行价 5.38 元/\allowbreak{}股，\allowbreak{}募集资金总额 16.14 亿元 & 百万吨稻谷精深加工及十万立方米绿色稻草板生产项目、\allowbreak{}农业高新技术大规模综合应用项目、\allowbreak{}浩良河分公司油改煤化肥工程项目；经证监发行字[2002]19 号文核准；募集资金于 2002-03-19 全部划入公司指定账户 \\
2003 & 变更募投并收购资产 & 收购农垦集团所持黑龙江省北大荒米业有限公司 98.55\% 股权 & 由原投资新建改为收购控股股东资产；其中 3.50 亿元来自北大荒上市募集资金 \\
2007-12 & 可转换公司债券 & 发行数量 15.00 亿元（15.00 亿元），\allowbreak{}面值 100.00 元，\allowbreak{}按面值发行 & 经证监发行字[2007]451 号文核准；原股东享有优先配售权；截至 2010-12-31 累计转股 1.43 亿股 \\
2025-05 & 现金分红（2024 年度） & 以总股本 17.78 亿股为基数每股派 0.55 元，\allowbreak{}共派发现金红利 9.78 亿元 & 回报股东；现金红利发放已于 5 月 23 日完成 \\
2026-03 & 现金分红预案（2025 年度） & 每 10 股派发现金股利 5.50 元（含税），\allowbreak{}合计拟派发 977，\allowbreak{}723，\allowbreak{}949.95 元 & 回报股东；不送红股，\allowbreak{}不实施资本公积金转增股本 \\
\bottomrule
\end{tabular}
\end{adjustbox}
\end{center}

\pfl{历史融投资情况} 从现金流量表看，2025 年公司偿还债务支付的现金 0.48 亿元，筹资活动现金流净额 -10.32 亿元。 2026 年上半年筹资活动现金流净额为 -9.78 亿元（取得借款 — 亿元、偿还债务 0.00 亿元，未年化），已转为净流出，处于净还债状态。 2025 年分配股利、利润或偿付利息支付的现金合计 9.76 亿元。 公开披露的融资记录共 5 笔，工具类型以 IPO、可转债为主；金额与用途见下表。 其中规模最大的两笔为：2007 年 12 月，可转债，发行数量 15.00 亿元（15.00 亿元），面值 100.00 元，按面值发行；2003 年，变更募投并收购资产，收购农垦集团所持黑龙江省北大荒米业有限公司 98.55\% 股权。 股本方面，近三年披露股本变动 3 次；总股本自 2021 年末以来未发生变化（17.78 亿股）。 分红方面，近三年现金分红 4 次、合计每股派现 1.96 元，最近一次实施在 2025 年 12 月。 综合看，公司融资结构上股权与债务融资并举；2025 年筹资活动现金流净额 -10.32 亿元，融资节奏仍以维持存量债务周转为主，与前述经营与投资现金流的方向基本一致。

\pfl{未来潜在融资需求分析} 2026 年 6 月末货币资金 19.27 亿元，而短期债务 0.03 亿元（现金短债比 566.76 倍），2026 年上半年经营现金流净额 24.86 亿元。账面货币资金可以覆盖短期债务，缺口不体现在静态偿付层面。 公司 2026 年上半年归母净利润为负（-5.37 亿元），但 6 月末资产负债率 39.5\% 尚未进入第一档（亏损且负债率 $>$65\%）的区间，当前融资需求以补充流动资金、支撑低谷期经营性支出为主。 静态偿付不存在缺口，后续融资更可能服务于产能或并购安排，而不是填补流动性窟窿。




\subsubsection{苏垦农发（601952）}
\label{co:601952}

\pfl{公司基本情况} 公司前身可追溯至 2011 年，江苏省农垦集团旗下的农业全产业链上市公司，自主经营耕地约 95.5 万亩，形成种植、种子、大米、农资、食用油一体化经营。 公司于 2017-05-15 在上交所首发上市，注册资本 13.78 亿元，总股本 13.78 亿股。 截至 2026 年 9 月 24 日收盘价 9.11 元，流通市值 125.5 亿元，近一年-3.5\%，流通市值在三级行业“粮食种植”的 2 家公司中按由大到小排第\zhnumber{2}位，近一年涨跌幅高于全样本中位数（-10.1\%）6.6 个百分点。 按 2025 年年报披露的行业构成，收入占比前三位为大米加工与销售分部 36.4\%；种植业分部 35.6\%；农资服务与销售分部 21.2\%。 与 2021 年年报相比，第一大行业口径由“米业”（39.7\%）变为“大米加工与销售分部”（36.4\%）；两期披露的分类口径有调整，占比差异中同时包含分类因素。

\pfl{历史股价} 自 2021 年 1 月至 2026 年 9 月，股价由 11.62 元变为 9.11 元，累计 -21.6\%，区间最高 16.70 元（2022 年 5 月）、最低 7.49 元（2026 年 6 月）；当前价相当于区间最高价的 54.6\%，为区间最低价的 1.22 倍。 月度收益的年化波动率 26.7\%，低于全样本中位数 37.4\%，在三级行业“粮食种植”的 2 家公司中波动率由高到低排第\zhnumber{1}位。 把股价阶段与公司事件对齐看，2022 年 5 月 至 2026 年 6 月股价下跌-55.1\%，同期（2026 年 6 月）公司完成对江苏大丰华丰种业股份有限公司的投资并购。 整体上看，区间的几处大级别拐点都能在公司的资本运作、股权变动或风险事件上找到对应，股价波动有明显的公司层面驱动，而非单纯的行业 beta。

\begin{center}
\begin{minipage}{\textwidth}\centering
\begin{tikzpicture}
\begin{axis}[
  width=12.6cm, height=3.9cm, scale only axis,
  xmin=2020.922, xmax=2026.888, ymin=6.97, ymax=17.87,
  xtick={2021,2022,2023,2024,2025,2026},
  yearticks,
  yticklabel style={font=\scriptsize},
  ylabel={元/股}, ylabel style={font=\scriptsize},
  grid=major, grid style={LightGray, line width=0.3pt},
  axis line style={InkGray, line width=0.5pt},
  every axis plot/.append style={line width=0.9pt},
]
\addplot[AgriGreen] table[x=x,y=y]{data/clean/profiles/px_601952.dat};
\addplot[SoftRed, dashed, line width=0.7pt] table[x=x,y=y]{data/clean/profiles/pxzz_601952.dat};
\addplot[only marks, mark=*, mark size=1.1pt, SoftRed] table[x=x,y=y]{data/clean/profiles/pxzz_601952.dat};
\end{axis}
\end{tikzpicture}

\captionof{figure}[苏垦农发（601952）股价与周期骨架]{苏垦农发（601952）月度收盘价（元/股）与 ZigZag 周期骨架（阈值 25\%，圆点为周期转折点）}
\end{minipage}
\end{center}

\pfl{过去周期分析} 以 25\% 的 ZigZag 阈值划分，样本区间内股价形成 1 个主升段与 2 个主跌段。 最大主跌段为 2022 年 5 月 至 2026 年 6 月，49 个月-55.1\%；最大主升段出现在 2021 年 9 月 至 2022 年 5 月，8 个月+63.9\%。 最近一次转折出现在 2026 年 6 月（下跌段自 2022 年 5 月起，49 个月-55.1\%），当前处于该段的延续之中。 公司股价与申万农林牧渔指数几乎同步见底，个股并未走出独立行情。 土地经营型企业的收入与地租成本高度相关（第 \ref{sec:planting} 节），粮价与政策补贴是其主要变量，公司 2026 年上半年实现盈利、半年 ROE 2.0\%（未年化），好于全样本中位数（0.75\%），下行周期对其报表的冲击小于同业。 综合区间形态看，该股单段涨跌幅度偏大、以趋势段为主（最大主升 +63.9\%、最大主跌 -55.1\%），年化波动率 26.7\% 与全样本中位数 37.4\% 相比波动小于样本中位数。

\pfl{盈利情况} 2026 年上半年营业收入 43.32 亿元（同比 -5.6\%），归母净利润 1.44 亿元，同比 -32.5\%。 以年报为趋势对照，2023---2025 年营业收入由 121.68 亿元变为 100.71 亿元（两年年均复合增速 -9.0\%），归母净利润由 8.16 亿元变为 5.48 亿元。 按半年报口径，2026 年上半年的 ROE 为 2.0\%（未年化）、销售净利率 3.6\%、毛利率 11.3\%，ROE 在“粮食种植”2 家公司中按数值由高到低排第\zhnumber{1}位。 同期全样本半年 ROE 中位数 0.75\%，公司与之相差 1.29 个百分点（略好于中位数公司）。 综合看，公司在报告期维持盈利（高于全样本半年 ROE 中位数 0.75\%），利润的可持续性取决于粮食种植景气的延续与成本管控的执行。

\begin{center}\small\setlength{\tabcolsep}{3.4pt}
\captionof{table}[苏垦农发（601952）关键财务指标]{苏垦农发（601952）关键财务指标（合并报表口径；两个半年报期均未年化；现金短债比 $=$ 货币资金 $\div$ 短期债务）}
\begin{adjustbox}{max width=\textwidth}
\begin{tabular}{@{}lrrrrr@{}}
\toprule
指标 & 2023 年 & 2024 年 & 2025 年 & 2025 年半年报 & 2026 年半年报 \\
\midrule
营业收入（亿元） & 121.68 & 109.17 & 100.71 & 45.88 & 43.32 \\
归母净利润（亿元） & 8.16 & 7.30 & 5.48 & 2.13 & 1.44 \\
毛利率（\%） & 13.3 & 14.5 & 13.2 & 12.0 & 11.3 \\
ROE（\%） & 12.8 & 10.8 & 7.9 & 3.1 & 2.0 \\
资产负债率（\%） & 50.8 & 49.5 & 48.9 & 50.7 & 46.5 \\
经营活动现金流净额（亿元） & 23.08 & 14.68 & 22.59 & 7.22 & 1.67 \\
货币资金（亿元） & 6.26 & 6.91 & 8.49 & 13.37 & 14.56 \\
有息负债合计（亿元） & 51.18 & 48.89 & 45.57 & 48.97 & 45.25 \\
现金短债比（倍） & 1.10 & 1.16 & 1.67 & 2.49 & 3.40 \\
\bottomrule
\end{tabular}
\end{adjustbox}
\end{center}

\begin{center}
\begin{minipage}{\textwidth}\centering
\begin{tikzpicture}
\begin{groupplot}[
  group style={group size=2 by 1, horizontal sep=0.60cm},
  width=6.85cm, height=4.60cm,
  tick label style={font=\tiny},
  title style={font=\scriptsize\heiti},
  yticklabel style={font=\tiny},
  ylabel style={font=\tiny},
  grid=major, grid style={LightGray, line width=0.3pt},
  axis line style={InkGray, line width=0.5pt},
  enlarge x limits={abs=0.8},
  legend style={font=\scriptsize, at={(0.5,-0.20)}, anchor=north,
                legend columns=2, draw=none, fill=none, column sep=6pt},
]
\nextgroupplot[
  ybar, bar width=4pt,
  ymin=-12.73, ymax=142.55,
  xtick={0,1,2,3,4,5.7},
  xticklabels={2021,2022,2023,2024,2025,2026H1},
  ylabel={亿元},
]
\addplot[fill=AgriGreen!75, draw=AgriGreen, bar shift=-2.6pt] table[x=i, y=rev]{data/clean/profiles/fin_601952.dat};
\addplot[fill=SoftRed!60, draw=SoftRed, bar shift=2.6pt] table[x=i, y=ni]{data/clean/profiles/fin_601952.dat};
\addplot[fill=AgriGreen!30, draw=AgriGreen, pattern=north east lines, bar shift=-2.6pt] table[x=x, y=rev]{data/clean/profiles/finh_601952.dat};
\addplot[fill=SoftRed!20, draw=SoftRed, pattern=north east lines, bar shift=2.6pt] table[x=x, y=ni]{data/clean/profiles/finh_601952.dat};
\legend{营业收入（年/半年）, 归母净利润（年/半年）}
\nextgroupplot[
  ymin=-4.5, ymax=61.5,
  xtick={0,1,2,3,4,5.7},
  xticklabels={2021,2022,2023,2024,2025,2026H1},
  ylabel={\%},
]
\addplot[AgriGold, mark=*, mark size=1.3pt] table[x=i, y=roe]{data/clean/profiles/fin_601952.dat};
\addplot[OilBlue, mark=square*, mark size=1.3pt] table[x=i, y=debt]{data/clean/profiles/fin_601952.dat};
\addplot[only marks, AgriGold, mark=*, mark size=2.0pt] table[x=x, y=roe]{data/clean/profiles/finh_601952.dat};
\addplot[only marks, OilBlue, mark=square*, mark size=2.0pt] table[x=x, y=debt]{data/clean/profiles/finh_601952.dat};
\legend{ROE, 资产负债率}
\end{groupplot}
\end{tikzpicture}
\begin{tikzpicture}
\begin{axis}[
  name=finbot,
  width=13.75cm, height=3.10cm, scale only axis,
  ymin=-5.79, ymax=64.83,
  xtick={0,1,2,3,4,5.7},
  xticklabels={2021,2022,2023,2024,2025,2026H1},
  tick label style={font=\tiny},
  yticklabel style={font=\tiny},
  ylabel={亿元}, ylabel style={font=\tiny},
  ybar, bar width=4pt,
  grid=major, grid style={LightGray, line width=0.3pt},
  axis line style={InkGray, line width=0.5pt},
  enlarge x limits={abs=0.8},
  legend style={font=\scriptsize, at={(0.5,-0.26)}, anchor=north,
                legend columns=2, draw=none, fill=none, column sep=10pt},
]
\addplot[fill=AgriGreen!55, draw=AgriGreen, bar shift=-3.0pt] table[x=i, y=cash]{data/clean/profiles/fin_601952.dat};
\addplot[fill=SoftRed!45, draw=SoftRed, bar shift=3.0pt] table[x=i, y=ibd]{data/clean/profiles/fin_601952.dat};
\addplot[fill=AgriGreen!25, draw=AgriGreen, pattern=north east lines, bar shift=-3.0pt] table[x=x, y=cash]{data/clean/profiles/finh_601952.dat};
\addplot[fill=SoftRed!20, draw=SoftRed, pattern=north east lines, bar shift=3.0pt] table[x=x, y=ibd]{data/clean/profiles/finh_601952.dat};
\legend{货币资金（左轴）, 有息负债（左轴）}
\end{axis}
\begin{axis}[
  at={(finbot.south east)}, anchor=south east,
  width=13.75cm, height=3.10cm, scale only axis,
  axis y line*=right, axis x line*=none, xtick=\empty,
  ymin=-6.00, ymax=30.01,
  tick label style={font=\tiny},
  yticklabel style={font=\tiny}, scaled y ticks=false,
  legend style={font=\scriptsize, at={(0.5,-0.44)}, anchor=north,
                legend columns=1, draw=none, fill=none},
]
\addplot[OilBlue, mark=*, mark size=2.2pt, line width=1.3pt] table[x=i, y=ocf]{data/clean/profiles/fin_601952.dat};
\addplot[only marks, OilBlue, mark=*, mark size=3.2pt] table[x=x, y=ocf]{data/clean/profiles/finh_601952.dat};
\legend{经营活动现金流净额（右轴，亿元，蓝点）}
\end{axis}
\end{tikzpicture}

\captionof{figure}[苏垦农发（601952）近年主要财务指标]{苏垦农发（601952）近年主要财务指标（左上：营业收入与归母净利润；右上：ROE 与资产负债率；下：货币资金、有息负债为左轴柱状，经营活动现金流净额为其右轴的蓝点折线。
2021---2025 年为年报口径，2026H1 为 2026 年半年报/半年末，与其前一年报之间留出间隔、以斜纹或空心点区分；半年报数据未年化）}
\end{minipage}
\end{center}

\pfl{财务分析} 2026 年 6 月末资产负债率 46.5\%，较 2025 年末下降 2.4 个百分点（样本中位数 52.2\%，所处三级行业内按由低到高排第\zhnumber{2}位）。 2026 年 6 月末货币资金 14.56 亿元、短期债务（短期借款与一年内到期非流动负债）4.28 亿元、长期债务 40.98 亿元，有息负债合计 45.25 亿元，现金短债比 3.40 倍——覆盖充分。 净有息负债 30.70 亿元，相当于其 2026 年上半年营业收入的 71\%。 流动比率 3.47、速动比率 2.08，短期偿债能力处于正常区间。 2026 年上半年经营活动现金流净额 1.67 亿元，经营现金流与利润同向为正，内生资金可以支撑运营。 2026 年 6 月末在建工程 3.47 亿元，较上年末增加 1.12 亿元，仍有在建投入。 综合看，现金短债比 3.40 倍，短期偿付压力可控；资产负债率低于样本中位数 5.7 个百分点；有息负债中长期债务占比更高，期限结构对短债滚动的压力较小。

\pfl{重大战略转型} 近三年公司的战略动作集中在：并购与重组（1 项）：2026 年 6 月，公司完成对江苏大丰华丰种业股份有限公司的投资并购。 公告口径上，近三年（2023 年 9 月---2026 年 9 月）公司披露重大重组与并购类公告 3 条、对外投资与转型类公告 4 条、回购与股权激励类公告 0 条，合计公告 289 条。 主营结构的口径变化已经落到报表上：2021 年年报的第一大行业为“米业”（39.7\%），2025 年年报为“大米加工与销售分部”（36.4\%）；公司在这两期的分部划分方式有过调整。 综合判断，报告期内公司有 7 条与战略相关的公告落地（重大重组与并购 3 条、对外投资与转型 4 条），资本运作与主业调整之间存在直接关联。

\begin{center}\small\setlength{\tabcolsep}{4pt}
\captionof{table}[苏垦农发（601952）历史融资与资本运作]{苏垦农发（601952）历史融资与资本运作（据公司公告与公开报道整理；金额为披露值，含首次公开发行、再融资、债券、银行借款与重整投资等）}
\begin{adjustbox}{max width=\textwidth}
\begin{tabular}{@{}p{1.45cm}p{2.55cm}p{4.05cm}p{4.95cm}@{}}
\toprule
时间 & 方式 & 规模 & 用途与说明 \\
\midrule
2017-05 & IPO（上海证券交易所主板） & 公开发行 26，\allowbreak{}000 万股，\allowbreak{}发行价 9.32 元/\allowbreak{}股，\allowbreak{}募集资金总额 24.23 亿元 & 百万亩农田改造建设项目（16.68 亿元）、\allowbreak{}大华种业集团改扩建项目（1.32 亿元）、\allowbreak{}农业科学研究院建设项目、\allowbreak{}农业信息化建设项目、\allowbreak{}补充流动资金（3.00 亿元）；；募集资金 2017-05-09 到账 \\
2022-04 & 变更部分募投项目 & 使用首发募投项目部分现金管理收益及利息 2.46 亿元投资建设“苏垦米业集团改扩建项目”（含江心沙、\allowbreak{}淮海公司、\allowbreak{}宝应湖三个子项目） & 大米加工产能扩建；2024 年 8 月因政策收严、\allowbreak{}淮海项目用地手续耗时较长 \\
2026-04-29 & 子公司增资扩股引进战略投资者（关联交易） & 控股子公司江苏种业完成 10.00 亿元外部资金引入及战略投资者引进；控股股东江苏省农垦集团投资 4.50 亿元认缴新增注册资本 2.90 亿元 & 优化江苏种业股权结构、\allowbreak{}充实种业资本实力；2026-06-17 完成工商变更登记；增资后江苏种业注册资本 20.00 亿元 \\
2026-08 & 中期现金分红 & 以 2026 年半年末总股本 13.78 亿股为基数每 10 股派现金红利 0.40 元 & 回报股东；属 2025 年度股东会授权董事会决策的权限范围，\allowbreak{}无需提交股东会审议 \\
\bottomrule
\end{tabular}
\end{adjustbox}
\end{center}

\pfl{历史融投资情况} 从现金流量表看，2025 年公司取得借款收到的现金 1.05 亿元、偿还债务支付的现金 2.15 亿元，筹资活动现金流净额 -13.17 亿元（2023 年为取得借款 3.68 亿元、偿还债务 7.84 亿元）。 2026 年上半年筹资活动现金流净额为 6.32 亿元（取得借款 0.10 亿元、偿还债务 1.16 亿元，未年化），仍处于净融入状态。 2025 年分配股利、利润或偿付利息支付的现金合计 4.17 亿元。 公开披露的融资记录共 4 笔，工具类型以 IPO、战略投资为主；金额与用途见下表。 其中规模最大的两笔为：2026 年 4 月，战略投资，控股子公司江苏种业完成 10.00 亿元外部资金引入及战略投资者引进，控股股东江苏省农垦集团投资 4.50 亿元认缴新增注册资本 2.90 亿元；2017 年 5 月，IPO，公开发行 26，000 万股，发行价 9.32 元/股，募集资金总额 24.23 亿元。 股本方面，近三年披露股本变动 3 次；总股本自 2021 年末以来未发生变化（13.78 亿股）。 分红方面，近三年现金分红 7 次、合计每股派现 1.12 元，最近一次实施在 2026 年 6 月。 综合看，公司融资结构上股权与债务融资并举；2025 年筹资活动现金流净额 -13.17 亿元，融资节奏仍以维持存量债务周转为主，与前述经营与投资现金流的方向基本一致。

\pfl{未来潜在融资需求分析} 2026 年 6 月末货币资金 14.56 亿元，而短期债务 4.28 亿元（现金短债比 3.40 倍），2026 年上半年经营现金流净额 1.67 亿元。账面货币资金可以覆盖短期债务，缺口不体现在静态偿付层面。 公司 2026 年 6 月末资产负债率 46.5\%、2026 年上半年 ROE 2.0\%（未年化），资产负债表尚有余量，融资需求不是“补血型”的刚性需求，而是配套在建项目投入的主动性安排。 资金需求还要叠加在建工程：期末在建工程 3.47 亿元、2026 年上半年购建固定资产等支出 2.02 亿元（未年化），这部分支出在周期底部通常会被主动放缓。 以现有货币资金与经营现金流衡量，公司不存在刚性补血的紧迫性，融资动作更多是主动型的。




\subsubsection{隆平高科（000998）}
\label{co:000998}

\pfl{公司基本情况} 隆平高科成立于 1999 年，由袁隆平院士发起设立、中信集团为实际控制人的综合性跨国种业集团，全球种业企业前八强，杂交水稻种子国内领先；因国内玉米种业高库存与巴西业务大幅承压，2026 年上半年亏损扩大，商誉规模较大。 公司于 2000-12-11 在深交所主板首发上市，注册资本 14.69 亿元，总股本 13.17 亿股。 截至 2026 年 9 月 24 日收盘价 8.68 元，流通市值 114.1 亿元，近一年-11.2\%，流通市值在三级行业“种子”的 8 家公司中按由大到小排第\zhnumber{1}位，近一年涨跌幅低于全样本中位数（-10.1\%）1.1 个百分点。 按 2026 年半年报披露的行业构成，收入几乎全部来自农业 100.0\%。

\pfl{历史股价} 自 2021 年 1 月至 2026 年 9 月，股价由 17.98 元变为 8.68 元，累计 -51.7\%，区间最高 24.04 元（2021 年 11 月）、最低 7.54 元（2026 年 6 月）；当前价相当于区间最高价的 36.1\%，为区间最低价的 1.15 倍。 月度收益的年化波动率 28.9\%，低于全样本中位数 37.4\%，在三级行业“种子”的 8 家公司中波动率由高到低排第\zhnumber{8}位。 把股价阶段与公司事件对齐看，2021 年 11 月 至 2026 年 6 月股价下跌-68.6\%，同期（2025 年 5 月）发行后中信集团合计持股比例由 17.36\% 提升至不超过 25.93\%。 从时间对应关系看，公司层面的资本运作与股权、风险事件解释了区间内多数急涨急跌，与行业指数的同步性相对有限。

\begin{center}
\begin{minipage}{\textwidth}\centering
\begin{tikzpicture}
\begin{axis}[
  width=12.6cm, height=3.9cm, scale only axis,
  xmin=2020.922, xmax=2026.888, ymin=7.01, ymax=25.72,
  xtick={2021,2022,2023,2024,2025,2026},
  yearticks,
  yticklabel style={font=\scriptsize},
  ylabel={元/股}, ylabel style={font=\scriptsize},
  grid=major, grid style={LightGray, line width=0.3pt},
  axis line style={InkGray, line width=0.5pt},
  every axis plot/.append style={line width=0.9pt},
]
\addplot[AgriGreen] table[x=x,y=y]{data/clean/profiles/px_000998.dat};
\addplot[SoftRed, dashed, line width=0.7pt] table[x=x,y=y]{data/clean/profiles/pxzz_000998.dat};
\addplot[only marks, mark=*, mark size=1.1pt, SoftRed] table[x=x,y=y]{data/clean/profiles/pxzz_000998.dat};
\end{axis}
\end{tikzpicture}

\captionof{figure}[隆平高科（000998）股价与周期骨架]{隆平高科（000998）月度收盘价（元/股）与 ZigZag 周期骨架（阈值 25\%，圆点为周期转折点）}
\end{minipage}
\end{center}

\pfl{过去周期分析} 以 25\% 的 ZigZag 阈值划分，样本区间内股价形成 1 个主升段与 2 个主跌段。 最大主跌段为 2021 年 11 月 至 2026 年 6 月，55 个月-68.6\%；最大主升段出现在 2021 年 6 月 至 2021 年 11 月，5 个月+49.9\%。 最近一次转折出现在 2026 年 6 月（下跌段自 2021 年 11 月起，55 个月-68.6\%），当前处于该段的延续之中。 公司股价与申万农林牧渔指数几乎同步见底，个股并未走出独立行情。 种子行业的价格由制种成本、粮价与政策（转基因商业化、种业振兴）共同决定（第 \ref{sec:planting} 节），与生猪价格的联动有限，公司 2026 年 6 月末资产负债率 59.1\%、半年 ROE -4.1\%（未年化），在本轮下行中属于随行业同步调整的一类。 综合区间形态看，该股单段涨跌幅度偏大、以趋势段为主（最大主升 +49.9\%、最大主跌 -68.6\%），年化波动率 28.9\% 与全样本中位数 37.4\% 相比波动小于样本中位数。

\pfl{盈利情况} 2026 年上半年营业收入 17.10 亿元（同比 -21.1\%），归母净利润 -2.73 亿元，亏损同比扩大 1.09 亿元。 以年报为趋势对照，2023---2025 年营业收入由 92.23 亿元变为 84.77 亿元（两年年均复合增速 -4.1\%），归母净利润由 2.00 亿元变为 1.66 亿元。 按半年报口径，2026 年上半年的 ROE 为 -4.1\%（未年化）、销售净利率 -22.4\%、毛利率 39.1\%，ROE 在“种子”8 家公司中按数值由高到低排第\zhnumber{8}位。 按同一口径，三级行业“种子”8 家公司的半年 ROE 中位数为 0.17\%，该公司低于其中位数 4.26 个百分点。 公司披露的应对举措集中在：2024 年 9 月，公司商誉账面价值 42.07 亿元。 综合看，公司仍处亏损区间、亏损同比扩大，半年 ROE 在三级行业内按由高到低排第\zhnumber{8}位，单位盈利能力的修复依赖行业景气与自身成本端的改善，报表弹性主要体现为价格与销量的方向性变化。

\begin{center}\small\setlength{\tabcolsep}{3.4pt}
\captionof{table}[隆平高科（000998）关键财务指标]{隆平高科（000998）关键财务指标（合并报表口径；两个半年报期均未年化；现金短债比 $=$ 货币资金 $\div$ 短期债务）}
\begin{adjustbox}{max width=\textwidth}
\begin{tabular}{@{}lrrrrr@{}}
\toprule
指标 & 2023 年 & 2024 年 & 2025 年 & 2025 年半年报 & 2026 年半年报 \\
\midrule
营业收入（亿元） & 92.23 & 85.66 & 84.77 & 21.66 & 17.10 \\
归母净利润（亿元） & 2.00 & 1.14 & 1.66 & -1.64 & -2.73 \\
毛利率（\%） & 39.8 & 36.6 & 39.3 & 36.3 & 39.1 \\
ROE（\%） & 3.2 & 2.2 & 2.9 & -3.1 & -4.1 \\
资产负债率（\%） & 63.5 & 66.5 & 59.5 & 60.0 & 59.1 \\
经营活动现金流净额（亿元） & 7.19 & 4.89 & 7.52 & -9.38 & -8.39 \\
货币资金（亿元） & 39.24 & 27.61 & 29.18 & 16.64 & 24.25 \\
有息负债合计（亿元） & 113.34 & 108.68 & 99.37 & 101.46 & 107.27 \\
现金短债比（倍） & 0.55 & 0.44 & 0.52 & 0.23 & 0.37 \\
\bottomrule
\end{tabular}
\end{adjustbox}
\end{center}

\begin{center}
\begin{minipage}{\textwidth}\centering
\begin{tikzpicture}
\begin{groupplot}[
  group style={group size=2 by 1, horizontal sep=0.60cm},
  width=6.85cm, height=4.60cm,
  tick label style={font=\tiny},
  title style={font=\scriptsize\heiti},
  yticklabel style={font=\tiny},
  ylabel style={font=\tiny},
  grid=major, grid style={LightGray, line width=0.3pt},
  axis line style={InkGray, line width=0.5pt},
  enlarge x limits={abs=0.8},
  legend style={font=\scriptsize, at={(0.5,-0.20)}, anchor=north,
                legend columns=2, draw=none, fill=none, column sep=6pt},
]
\nextgroupplot[
  ybar, bar width=4pt,
  ymin=-18.38, ymax=104.30,
  xtick={0,1,2,3,4,5.7},
  xticklabels={2021,2022,2023,2024,2025,2026H1},
  ylabel={亿元},
]
\addplot[fill=AgriGreen!75, draw=AgriGreen, bar shift=-2.6pt] table[x=i, y=rev]{data/clean/profiles/fin_000998.dat};
\addplot[fill=SoftRed!60, draw=SoftRed, bar shift=2.6pt] table[x=i, y=ni]{data/clean/profiles/fin_000998.dat};
\addplot[fill=AgriGreen!30, draw=AgriGreen, pattern=north east lines, bar shift=-2.6pt] table[x=x, y=rev]{data/clean/profiles/finh_000998.dat};
\addplot[fill=SoftRed!20, draw=SoftRed, pattern=north east lines, bar shift=2.6pt] table[x=x, y=ni]{data/clean/profiles/finh_000998.dat};
\legend{营业收入（年/半年）, 归母净利润（年/半年）}
\nextgroupplot[
  ymin=-20.0, ymax=74.5,
  xtick={0,1,2,3,4,5.7},
  xticklabels={2021,2022,2023,2024,2025,2026H1},
  ylabel={\%},
]
\addplot[AgriGold, mark=*, mark size=1.3pt] table[x=i, y=roe]{data/clean/profiles/fin_000998.dat};
\addplot[OilBlue, mark=square*, mark size=1.3pt] table[x=i, y=debt]{data/clean/profiles/fin_000998.dat};
\addplot[only marks, AgriGold, mark=*, mark size=2.0pt] table[x=x, y=roe]{data/clean/profiles/finh_000998.dat};
\addplot[only marks, OilBlue, mark=square*, mark size=2.0pt] table[x=x, y=debt]{data/clean/profiles/finh_000998.dat};
\legend{ROE, 资产负债率}
\end{groupplot}
\end{tikzpicture}
\begin{tikzpicture}
\begin{axis}[
  name=finbot,
  width=13.75cm, height=3.10cm, scale only axis,
  ymin=-11.33, ymax=126.94,
  xtick={0,1,2,3,4,5.7},
  xticklabels={2021,2022,2023,2024,2025,2026H1},
  tick label style={font=\tiny},
  yticklabel style={font=\tiny},
  ylabel={亿元}, ylabel style={font=\tiny},
  ybar, bar width=4pt,
  grid=major, grid style={LightGray, line width=0.3pt},
  axis line style={InkGray, line width=0.5pt},
  enlarge x limits={abs=0.8},
  legend style={font=\scriptsize, at={(0.5,-0.26)}, anchor=north,
                legend columns=2, draw=none, fill=none, column sep=10pt},
]
\addplot[fill=AgriGreen!55, draw=AgriGreen, bar shift=-3.0pt] table[x=i, y=cash]{data/clean/profiles/fin_000998.dat};
\addplot[fill=SoftRed!45, draw=SoftRed, bar shift=3.0pt] table[x=i, y=ibd]{data/clean/profiles/fin_000998.dat};
\addplot[fill=AgriGreen!25, draw=AgriGreen, pattern=north east lines, bar shift=-3.0pt] table[x=x, y=cash]{data/clean/profiles/finh_000998.dat};
\addplot[fill=SoftRed!20, draw=SoftRed, pattern=north east lines, bar shift=3.0pt] table[x=x, y=ibd]{data/clean/profiles/finh_000998.dat};
\legend{货币资金（左轴）, 有息负债（左轴）}
\end{axis}
\begin{axis}[
  at={(finbot.south east)}, anchor=south east,
  width=13.75cm, height=3.10cm, scale only axis,
  axis y line*=right, axis x line*=none, xtick=\empty,
  ymin=-14.30, ymax=21.15,
  tick label style={font=\tiny},
  yticklabel style={font=\tiny}, scaled y ticks=false,
  legend style={font=\scriptsize, at={(0.5,-0.44)}, anchor=north,
                legend columns=1, draw=none, fill=none},
]
\addplot[OilBlue, mark=*, mark size=2.2pt, line width=1.3pt] table[x=i, y=ocf]{data/clean/profiles/fin_000998.dat};
\addplot[only marks, OilBlue, mark=*, mark size=3.2pt] table[x=x, y=ocf]{data/clean/profiles/finh_000998.dat};
\legend{经营活动现金流净额（右轴，亿元，蓝点）}
\end{axis}
\end{tikzpicture}

\captionof{figure}[隆平高科（000998）近年主要财务指标]{隆平高科（000998）近年主要财务指标（左上：营业收入与归母净利润；右上：ROE 与资产负债率；下：货币资金、有息负债为左轴柱状，经营活动现金流净额为其右轴的蓝点折线。
2021---2025 年为年报口径，2026H1 为 2026 年半年报/半年末，与其前一年报之间留出间隔、以斜纹或空心点区分；半年报数据未年化）}
\end{minipage}
\end{center}

\pfl{财务分析} 2026 年 6 月末资产负债率 59.1\%，较 2025 年末下降 0.4 个百分点（样本中位数 52.2\%，所处三级行业内按由低到高排第\zhnumber{7}位）。 2026 年 6 月末货币资金 24.25 亿元、短期债务（短期借款与一年内到期非流动负债）64.83 亿元、长期债务 42.44 亿元，有息负债合计 107.27 亿元，现金短债比 0.37 倍——覆盖偏紧。 净有息负债 83.02 亿元，相当于其 2026 年上半年营业收入的 486\%。 流动比率 1.06、速动比率 0.60，短期流动性无明显缺口。 2026 年上半年经营活动现金流净额 -8.39 亿元，经营现金流与利润同时为负，现金消耗较快。 2026 年 6 月末在建工程 1.39 亿元，较上年末增加 0.41 亿元，仍有在建投入。 2026 年 6 月末商誉 43.51 亿元，占归母净资产的 42.2\%，是净资产中最脆弱的一块。 综合看，账面货币资金对有息负债与短期债务的覆盖不足，滚动偿付对新增融资的依赖度较高；资产负债率高于样本中位数 6.9 个百分点；有息负债中短期债务占比更高，期限结构仍以滚动续作为主。

\pfl{重大战略转型} 公司的战略动作可归为以下几类：其他动作（3 项）：2024 年 8 月，公司在机构调研中表示正积极应对转基因玉米市场日益增长的机遇；2019 年，成立隆平生物；并购与重组（2 项）：2025 年 3 月，一方面大力处置和剥离非核心资产（含与主业无关的类金融业务）；2017 年，先后收购湖北惠民科技、湖南金稻米业、三瑞农科、河北巡天农业以及巴西陶氏益农特定玉米种子业务。 公告口径上，近三年（2023 年 9 月---2026 年 9 月）公司披露重大重组与并购类公告 4 条、对外投资与转型类公告 0 条、回购与股权激励类公告 4 条，合计公告 398 条。 第一大行业仍是“农业”，收入占比 100.0\%，与 2021 年年报相比没有方向性变化。 综合判断，报告期内公司有 4 条与战略相关的公告落地（重大重组与并购 4 条、对外投资与转型 0 条），资本运作与主业调整之间存在直接关联。

\begin{center}\small\setlength{\tabcolsep}{4pt}
\captionof{table}[隆平高科（000998）历史融资与资本运作]{隆平高科（000998）历史融资与资本运作（据公司公告与公开报道整理；金额为披露值，含首次公开发行、再融资、债券、银行借款与重整投资等）}
\begin{adjustbox}{max width=\textwidth}
\begin{tabular}{@{}p{1.45cm}p{2.55cm}p{4.05cm}p{4.95cm}@{}}
\toprule
时间 & 方式 & 规模 & 用途与说明 \\
\midrule
2016-01 & 定向增发（非公开发行） & 向中信兴业、\allowbreak{}中信建设、\allowbreak{}信农投资、\allowbreak{}现代种业基金和汇添富－优势企业定增计划 5 号资产管理计划合计非公开发行 2.60 亿股 & 发行完成后中信系合计持股 18.72\%，\allowbreak{}中信集团成为实际控制人 \\
2018-03 & 发行股份购买资产 & 向王义波、\allowbreak{}彭泽斌等 45 名自然人发行股份购买联创种业 90\% 股权 & 注入优质玉米种子资产；2018 年 8 月通过证监会审批 \\
2023-08 & 股权受让（现金收购） & 斥资 8.00 亿元受让隆平农业发展股份有限公司 7.14\% 股份 & 获取巴西玉米种子业务，\allowbreak{}扩大营收规模；隆平发展 2023 年实现营收 92.23 亿元 \\
2025-05 & 定向增发（向特定对象发行） & 发行价格 7.87 元/\allowbreak{}股，\allowbreak{}数量不超过 1.52 亿股；实际募集资金人民币 10.35 亿元 & 扣除发行费用后全部用于偿还银行贷款及补充流动资金；发行对象为公司控股股东中信农业科技股份有限公司，\allowbreak{}北京中信信远向中信农业增资 36.00 亿元 \\
\bottomrule
\end{tabular}
\end{adjustbox}
\end{center}

\pfl{历史融投资情况} 从现金流量表看，2025 年公司取得借款收到的现金 102.29 亿元、偿还债务支付的现金 110.19 亿元、吸收权益性投资 22.01 亿元，筹资活动现金流净额 -0.24 亿元（2023 年为取得借款 102.48 亿元、偿还债务 65.61 亿元）。 2026 年上半年筹资活动现金流净额为 3.44 亿元（取得借款 46.34 亿元、偿还债务 38.23 亿元，未年化），仍处于净融入状态。 2025 年分配股利、利润或偿付利息支付的现金合计 6.75 亿元。 公开披露的融资记录共 4 笔，工具类型以定向增发、发行股份购买资产为主；金额与用途见下表。 其中规模最大的两笔为：2025 年 5 月，定向增发，发行价格 7.87 元/股，数量不超过 1.52 亿股，实际募集资金人民币 10.35 亿元；2023 年 8 月，股权受让，斥资 8.00 亿元受让隆平农业发展股份有限公司 7.14\% 股份。 股本方面，近三年披露股本变动 3 次；总股本自 2021 年末以来未发生变化（13.17 亿股）。 分红方面，近三年现金分红 3 次、合计每股派现 0.13 元，最近一次实施在 2025 年 12 月。 近三年“再融资”类公告 47 条，涉及定向增发、募集资金使用。 综合看，公司融资结构上股权与债务融资并举；2025 年筹资活动现金流净额 -0.24 亿元，融资节奏仍以维持存量债务周转为主，与前述经营与投资现金流的方向基本一致。

\pfl{未来潜在融资需求分析} 2026 年 6 月末货币资金 24.25 亿元，而短期债务 64.83 亿元（现金短债比 0.37 倍），2026 年上半年经营现金流净额 -8.39 亿元。按“短期债务减去账面货币资金”的滚动偿付口径，静态缺口约 40.58 亿元；若再计入上半年亏损的年化额（5.45 亿元），维持一年正常经营所需的外部资金约 40.58---46.03 亿元。 公司 2026 年上半年归母净利润为负（-2.73 亿元），但 6 月末资产负债率 59.1\% 尚未进入第一档（亏损且负债率 $>$65\%）的区间，当前融资需求以补充流动资金、支撑低谷期经营性支出为主。 公开披露的后续资金安排线索包括：股权融资线索：2024 年 8 月，公司拟发行不超过 1.52 亿股、募集资金不超过 12.00 亿元，由控股股东中信农业全额认购。 资金需求还要叠加在建工程：期末在建工程 1.39 亿元、2026 年上半年购建固定资产等支出 1.53 亿元（未年化），这部分支出在周期底部通常会被主动放缓。 综合看，公司的外部资金需求以营运资金周转与在建项目投入为主，滚动偿付口径下的静态缺口约 40.58 亿元，融资节奏将受制于银行续贷意愿与股权融资窗口。




\subsubsection{登海种业（002041）}
\label{co:002041}

\pfl{公司基本情况} 登海种业成立于 2000 年，国内玉米杂交种龙头之一，以“登海”“先玉”系列玉米种为核心，玉米种收入占营业收入约九成，控股股东莱州市农业科学院、李登海为实际控制人；已审定多个转基因玉米品种，2026 年上半年业绩同比明显回升。 公司于 2005-04-18 在深交所主板首发上市，注册资本 8.80 亿元，总股本 8.80 亿股。 截至 2026 年 9 月 24 日收盘价 9.94 元，流通市值 87.5 亿元，近一年+7.2\%，流通市值在三级行业“种子”的 8 家公司中按由大到小排第\zhnumber{2}位，近一年涨跌幅高于全样本中位数（-10.1\%）17.3 个百分点。 按 2026 年半年报披露的行业构成，收入几乎全部来自农业 100.0\%。

\pfl{历史股价} 自 2021 年 1 月至 2026 年 9 月，股价由 18.52 元变为 9.94 元，累计 -46.3\%，区间最高 26.42 元（2021 年 11 月）、最低 7.87 元（2026 年 6 月）；当前价相当于区间最高价的 37.6\%，为区间最低价的 1.26 倍。 月度收益的年化波动率 33.5\%，低于全样本中位数 37.4\%，在三级行业“种子”的 8 家公司中波动率由高到低排第\zhnumber{7}位。

\begin{center}
\begin{minipage}{\textwidth}\centering
\begin{tikzpicture}
\begin{axis}[
  width=12.6cm, height=3.9cm, scale only axis,
  xmin=2020.922, xmax=2026.888, ymin=7.32, ymax=28.27,
  xtick={2021,2022,2023,2024,2025,2026},
  yearticks,
  yticklabel style={font=\scriptsize},
  ylabel={元/股}, ylabel style={font=\scriptsize},
  grid=major, grid style={LightGray, line width=0.3pt},
  axis line style={InkGray, line width=0.5pt},
  every axis plot/.append style={line width=0.9pt},
]
\addplot[AgriGreen] table[x=x,y=y]{data/clean/profiles/px_002041.dat};
\addplot[SoftRed, dashed, line width=0.7pt] table[x=x,y=y]{data/clean/profiles/pxzz_002041.dat};
\addplot[only marks, mark=*, mark size=1.1pt, SoftRed] table[x=x,y=y]{data/clean/profiles/pxzz_002041.dat};
\end{axis}
\end{tikzpicture}

\captionof{figure}[登海种业（002041）股价与周期骨架]{登海种业（002041）月度收盘价（元/股）与 ZigZag 周期骨架（阈值 25\%，圆点为周期转折点）}
\end{minipage}
\end{center}

\pfl{过去周期分析} 以 25\% 的 ZigZag 阈值划分，样本区间内股价形成 3 个主升段与 3 个主跌段。 最大主跌段为 2021 年 11 月 至 2024 年 6 月，31 个月-69.4\%；最大主升段出现在 2021 年 6 月 至 2021 年 11 月，5 个月+74.9\%。 最近一次转折出现在 2026 年 9 月（上涨段自 2026 年 6 月起，3 个月+26.3\%），当前处于该段之后的震荡之中。 公司股价与申万农林牧渔指数几乎同步见底，个股并未走出独立行情。 种子行业的价格由制种成本、粮价与政策（转基因商业化、种业振兴）共同决定（第 \ref{sec:planting} 节），与生猪价格的联动有限，公司 2026 年上半年实现盈利、半年 ROE 1.3\%（未年化），好于全样本中位数（0.75\%），下行周期对其报表的冲击小于同业。 综合区间形态看，该股单段涨跌幅度偏大、以趋势段为主（最大主升 +74.9\%、最大主跌 -69.4\%），年化波动率 33.5\% 与全样本中位数 37.4\% 相比波动小于样本中位数。

\pfl{盈利情况} 2026 年上半年营业收入 4.69 亿元（同比 +27.0\%），归母净利润 0.50 亿元，同比 +41.0\%。 以年报为趋势对照，2023---2025 年营业收入由 15.52 亿元变为 11.02 亿元（两年年均复合增速 -15.7\%），归母净利润由 2.56 亿元变为 0.92 亿元。 按半年报口径，2026 年上半年的 ROE 为 1.3\%（未年化）、销售净利率 8.9\%、毛利率 30.4\%，ROE 在“种子”8 家公司中按数值由高到低排第\zhnumber{4}位。 按同一口径，三级行业“种子”8 家公司的半年 ROE 中位数为 0.17\%，该公司高于其中位数 1.17 个百分点。 从公司披露的原因看，业绩变动主要来自：2025 年上半年，当期扣非净利润亏损 433.00 万元；2026 年上半年，主要原因是本期借款借贷净额较上年同期减少 2，555.00 万元及上期分红影响。 综合看，公司在报告期维持盈利（高于全样本半年 ROE 中位数 0.75\%），利润的可持续性取决于种子景气的延续与成本管控的执行。

\begin{center}\small\setlength{\tabcolsep}{3.4pt}
\captionof{table}[登海种业（002041）关键财务指标]{登海种业（002041）关键财务指标（合并报表口径；两个半年报期均未年化；现金短债比 $=$ 货币资金 $\div$ 短期债务）}
\begin{adjustbox}{max width=\textwidth}
\begin{tabular}{@{}lrrrrr@{}}
\toprule
指标 & 2023 年 & 2024 年 & 2025 年 & 2025 年半年报 & 2026 年半年报 \\
\midrule
营业收入（亿元） & 15.52 & 12.46 & 11.02 & 3.69 & 4.69 \\
归母净利润（亿元） & 2.56 & 0.57 & 0.92 & 0.35 & 0.50 \\
毛利率（\%） & 28.8 & 28.1 & 29.6 & 27.7 & 30.4 \\
ROE（\%） & 7.4 & 1.6 & 2.5 & 1.0 & 1.3 \\
资产负债率（\%） & 19.1 & 19.1 & 18.7 & 16.9 & 15.8 \\
经营活动现金流净额（亿元） & 1.79 & 1.28 & 2.46 & -0.63 & -0.46 \\
货币资金（亿元） & 13.17 & 15.36 & 23.49 & 17.07 & 21.43 \\
有息负债合计（亿元） & 0.14 & 0.27 & 0.39 & 0.39 & 0.26 \\
现金短债比（倍） & 167.83 & 74.75 & 74.45 & 54.10 & 112.76 \\
\bottomrule
\end{tabular}
\end{adjustbox}
\end{center}

\begin{center}
\begin{minipage}{\textwidth}\centering
\begin{tikzpicture}
\begin{groupplot}[
  group style={group size=2 by 1, horizontal sep=0.60cm},
  width=6.85cm, height=4.60cm,
  tick label style={font=\tiny},
  title style={font=\scriptsize\heiti},
  yticklabel style={font=\tiny},
  ylabel style={font=\tiny},
  grid=major, grid style={LightGray, line width=0.3pt},
  axis line style={InkGray, line width=0.5pt},
  enlarge x limits={abs=0.8},
  legend style={font=\scriptsize, at={(0.5,-0.20)}, anchor=north,
                legend columns=2, draw=none, fill=none, column sep=6pt},
]
\nextgroupplot[
  ybar, bar width=4pt,
  ymin=-1.55, ymax=17.38,
  xtick={0,1,2,3,4,5.7},
  xticklabels={2021,2022,2023,2024,2025,2026H1},
  ylabel={亿元},
]
\addplot[fill=AgriGreen!75, draw=AgriGreen, bar shift=-2.6pt] table[x=i, y=rev]{data/clean/profiles/fin_002041.dat};
\addplot[fill=SoftRed!60, draw=SoftRed, bar shift=2.6pt] table[x=i, y=ni]{data/clean/profiles/fin_002041.dat};
\addplot[fill=AgriGreen!30, draw=AgriGreen, pattern=north east lines, bar shift=-2.6pt] table[x=x, y=rev]{data/clean/profiles/finh_002041.dat};
\addplot[fill=SoftRed!20, draw=SoftRed, pattern=north east lines, bar shift=2.6pt] table[x=x, y=ni]{data/clean/profiles/finh_002041.dat};
\legend{营业收入（年/半年）, 归母净利润（年/半年）}
\nextgroupplot[
  ymin=-1.8, ymax=24.7,
  xtick={0,1,2,3,4,5.7},
  xticklabels={2021,2022,2023,2024,2025,2026H1},
  ylabel={\%},
]
\addplot[AgriGold, mark=*, mark size=1.3pt] table[x=i, y=roe]{data/clean/profiles/fin_002041.dat};
\addplot[OilBlue, mark=square*, mark size=1.3pt] table[x=i, y=debt]{data/clean/profiles/fin_002041.dat};
\addplot[only marks, AgriGold, mark=*, mark size=2.0pt] table[x=x, y=roe]{data/clean/profiles/finh_002041.dat};
\addplot[only marks, OilBlue, mark=square*, mark size=2.0pt] table[x=x, y=debt]{data/clean/profiles/finh_002041.dat};
\legend{ROE, 资产负债率}
\end{groupplot}
\end{tikzpicture}
\begin{tikzpicture}
\begin{axis}[
  name=finbot,
  width=13.75cm, height=3.10cm, scale only axis,
  ymin=-2.35, ymax=26.31,
  xtick={0,1,2,3,4,5.7},
  xticklabels={2021,2022,2023,2024,2025,2026H1},
  tick label style={font=\tiny},
  yticklabel style={font=\tiny},
  ylabel={亿元}, ylabel style={font=\tiny},
  ybar, bar width=4pt,
  grid=major, grid style={LightGray, line width=0.3pt},
  axis line style={InkGray, line width=0.5pt},
  enlarge x limits={abs=0.8},
  legend style={font=\scriptsize, at={(0.5,-0.26)}, anchor=north,
                legend columns=2, draw=none, fill=none, column sep=10pt},
]
\addplot[fill=AgriGreen!55, draw=AgriGreen, bar shift=-3.0pt] table[x=i, y=cash]{data/clean/profiles/fin_002041.dat};
\addplot[fill=SoftRed!45, draw=SoftRed, bar shift=3.0pt] table[x=i, y=ibd]{data/clean/profiles/fin_002041.dat};
\addplot[fill=AgriGreen!25, draw=AgriGreen, pattern=north east lines, bar shift=-3.0pt] table[x=x, y=cash]{data/clean/profiles/finh_002041.dat};
\addplot[fill=SoftRed!20, draw=SoftRed, pattern=north east lines, bar shift=3.0pt] table[x=x, y=ibd]{data/clean/profiles/finh_002041.dat};
\legend{货币资金（左轴）, 有息负债（左轴）}
\end{axis}
\begin{axis}[
  at={(finbot.south east)}, anchor=south east,
  width=13.75cm, height=3.10cm, scale only axis,
  axis y line*=right, axis x line*=none, xtick=\empty,
  ymin=-1.53, ymax=4.89,
  tick label style={font=\tiny},
  yticklabel style={font=\tiny}, scaled y ticks=false,
  legend style={font=\scriptsize, at={(0.5,-0.44)}, anchor=north,
                legend columns=1, draw=none, fill=none},
]
\addplot[OilBlue, mark=*, mark size=2.2pt, line width=1.3pt] table[x=i, y=ocf]{data/clean/profiles/fin_002041.dat};
\addplot[only marks, OilBlue, mark=*, mark size=3.2pt] table[x=x, y=ocf]{data/clean/profiles/finh_002041.dat};
\legend{经营活动现金流净额（右轴，亿元，蓝点）}
\end{axis}
\end{tikzpicture}

\captionof{figure}[登海种业（002041）近年主要财务指标]{登海种业（002041）近年主要财务指标（左上：营业收入与归母净利润；右上：ROE 与资产负债率；下：货币资金、有息负债为左轴柱状，经营活动现金流净额为其右轴的蓝点折线。
2021---2025 年为年报口径，2026H1 为 2026 年半年报/半年末，与其前一年报之间留出间隔、以斜纹或空心点区分；半年报数据未年化）}
\end{minipage}
\end{center}

\pfl{财务分析} 2026 年 6 月末资产负债率 15.8\%，较 2025 年末下降 2.9 个百分点（样本中位数 52.2\%，所处三级行业内按由低到高排第\zhnumber{2}位）。 2026 年 6 月末货币资金 21.43 亿元、短期债务（短期借款与一年内到期非流动负债）0.19 亿元、长期债务 0.07 亿元，有息负债合计 0.26 亿元，现金短债比 112.76 倍——覆盖充分。 净有息负债 -21.16 亿元，相当于其 2026 年上半年营业收入的 -451\%。 流动比率 6.34、速动比率 5.53，短期偿债能力处于正常区间。 2026 年上半年经营活动现金流净额 -0.46 亿元，净利润为正但经营现金流为负，利润的现金含量偏低。 2026 年 6 月末在建工程 0.52 亿元，较上年末增加 0.05 亿元，仍有在建投入。 综合看，现金短债比 112.76 倍，短期偿付压力可控；资产负债率低于样本中位数 36.4 个百分点；有息负债中短期债务占比更高，期限结构仍以滚动续作为主。

\pfl{重大战略转型} 公告口径上，近三年（2023 年 9 月---2026 年 9 月）公司披露重大重组与并购类公告 0 条、对外投资与转型类公告 1 条、回购与股权激励类公告 0 条，合计公告 151 条。 第一大行业仍是“农业”，收入占比 100.0\%，与 2021 年年报相比没有方向性变化。 综合判断，报告期内公司有 1 条与战略相关的公告落地（重大重组与并购 0 条、对外投资与转型 1 条），资本运作与主业调整之间存在直接关联。

\begin{center}\small\setlength{\tabcolsep}{4pt}
\captionof{table}[登海种业（002041）历史融资与资本运作]{登海种业（002041）历史融资与资本运作（据公司公告与公开报道整理；金额为披露值，含首次公开发行、再融资、债券、银行借款与重整投资等）}
\begin{adjustbox}{max width=\textwidth}
\begin{tabular}{@{}p{1.45cm}p{2.55cm}p{4.05cm}p{4.95cm}@{}}
\toprule
时间 & 方式 & 规模 & 用途与说明 \\
\midrule
2005 年以来 & 现金分红（累计） & A 股上市后累计派现 5.84 亿元；近三年累计派现 5,\allowbreak{}720.00 万元 & 回报股东；2026 年上半年不派发现金红利、\allowbreak{}不送红股、\allowbreak{}不以公积金转增股本 \\
2005—2026 & IPO 后的股权／债权再融资 & --- & 本次检索未发现公司 IPO 之后的定增、\allowbreak{}可转债或公司债记录；2015 年年度报告披露“公司报告期无募集资金使用情况” \\
2005-04 & IPO（深圳证券交易所中小企业板） & 发行 2，\allowbreak{}200 万股，\allowbreak{}发行价 16.70 元/\allowbreak{}股 & 发行费用明细：承销费用 1，\allowbreak{}上网发行费用 97.36 万元；验资确认截至 2005-04-12 募集资金总额 3.67 亿元 \\
\bottomrule
\end{tabular}
\end{adjustbox}
\end{center}

\pfl{历史融投资情况} 从现金流量表看，2025 年公司取得借款收到的现金 0.50 亿元、偿还债务支付的现金 0.37 亿元，筹资活动现金流净额 -0.16 亿元（2023 年为取得借款 0.13 亿元、偿还债务 0.07 亿元）。 2026 年上半年筹资活动现金流净额为 -0.14 亿元（取得借款 0.19 亿元、偿还债务 0.32 亿元，未年化），已转为净流出，处于净还债状态。 2025 年分配股利、利润或偿付利息支付的现金合计 0.27 亿元。 公开披露的融资记录共 3 笔，工具类型以 IPO 为主；金额与用途见下表。 其中规模最大的两笔为：2005 年 4 月，IPO，发行 2，200 万股，发行价 16.70 元/股；2005 年，现金分红，A 股上市后累计派现 5.84 亿元，近三年累计派现 5,720.00 万元。 股本方面，近三年披露股本变动 3 次；总股本自 2021 年末以来未发生变化（8.80 亿股）。 分红方面，近三年现金分红 2 次、合计每股派现 0.06 元，最近一次实施在 2024 年 12 月。 综合看，公司融资结构上股权与债务融资并举；2025 年筹资活动现金流净额 -0.16 亿元，融资节奏仍以维持存量债务周转为主，与前述经营与投资现金流的方向基本一致。

\pfl{未来潜在融资需求分析} 2026 年 6 月末货币资金 21.43 亿元，而短期债务 0.19 亿元（现金短债比 112.76 倍），2026 年上半年经营现金流净额 -0.46 亿元。账面货币资金可以覆盖短期债务，缺口不体现在静态偿付层面。 公司 2026 年 6 月末资产负债率 15.8\%、2026 年上半年 ROE 1.3\%（未年化），账面杠杆不高、盈利能力一般，短期内没有必须依靠外部融资才能维持的经营缺口；后续融资更可能指向技改、扩产或产业并购。 公开披露的后续资金安排线索包括：债务与授信安排：2026 年 6 月，本期借款借贷净额较上年同期减少 2，555.00 万元。 资金需求还要叠加在建工程：期末在建工程 0.52 亿元、2026 年上半年购建固定资产等支出 0.39 亿元（未年化），这部分支出在周期底部通常会被主动放缓。 综合看，公司在静态偿付层面不存在缺口，外部融资更多是配合扩张节奏与并购整合的主动性安排，而非偿债压力驱动。




\subsubsection{亚盛集团（600108）}
\label{co:600108}

\pfl{公司基本情况} 公司前身可追溯至 1995 年，甘肃省农垦集团控股的现代农业企业集团，种植基地主要分布在甘肃、内蒙，以牧草、马铃薯、制种玉米等规模化种植为基础并延伸加工与农资服务；2026 年上半年盈利微降，经营现金流净流出扩大，资金缺口主要靠借款与中期票据填补。 公司于 1997-08-18 在上交所首发上市，注册资本 19.47 亿元，总股本 19.47 亿股。 截至 2026 年 9 月 24 日收盘价 4.13 元，流通市值 80.4 亿元，近一年+40.5\%，流通市值在三级行业“其他种植业”的 6 家公司中按由大到小排第\zhnumber{2}位，近一年涨跌幅高于全样本中位数（-10.1\%）50.6 个百分点。 按 2026 年半年报披露的行业构成，收入占比前三位为农业 57.1\%；商贸 36.1\%；工业 6.0\%。 股权结构方面，控股股东为甘肃省农垦集团有限责任公司。

\pfl{历史股价} 本报告样本区间（2021 年 1 月 至 2026 年 9 月）内，股价由 2.94 元变为 4.13 元，累计 +40.5\%，区间最高 4.55 元（2026 年 3 月）、最低 2.37 元（2024 年 6 月）；当前价相当于区间最高价的 90.8\%，为区间最低价的 1.74 倍。 月度收益的年化波动率 28.2\%，低于全样本中位数 37.4\%，在三级行业“其他种植业”的 6 家公司中波动率由高到低排第\zhnumber{6}位。

\begin{center}
\begin{minipage}{\textwidth}\centering
\begin{tikzpicture}
\begin{axis}[
  width=12.6cm, height=3.9cm, scale only axis,
  xmin=2020.922, xmax=2026.888, ymin=2.20, ymax=4.87,
  xtick={2021,2022,2023,2024,2025,2026},
  yearticks,
  yticklabel style={font=\scriptsize},
  ylabel={元/股}, ylabel style={font=\scriptsize},
  grid=major, grid style={LightGray, line width=0.3pt},
  axis line style={InkGray, line width=0.5pt},
  every axis plot/.append style={line width=0.9pt},
]
\addplot[AgriGreen] table[x=x,y=y]{data/clean/profiles/px_600108.dat};
\addplot[SoftRed, dashed, line width=0.7pt] table[x=x,y=y]{data/clean/profiles/pxzz_600108.dat};
\addplot[only marks, mark=*, mark size=1.1pt, SoftRed] table[x=x,y=y]{data/clean/profiles/pxzz_600108.dat};
\end{axis}
\end{tikzpicture}

\captionof{figure}[亚盛集团（600108）股价与周期骨架]{亚盛集团（600108）月度收盘价（元/股）与 ZigZag 周期骨架（阈值 25\%，圆点为周期转折点）}
\end{minipage}
\end{center}

\pfl{过去周期分析} 以 25\% 的 ZigZag 阈值划分，样本区间内股价形成 3 个主升段与 3 个主跌段。 最大主升段出现在 2024 年 6 月 至 2026 年 3 月，21 个月+92.0\%；最大主跌段为 2022 年 5 月 至 2024 年 6 月，25 个月-39.7\%。 最近一次转折出现在 2026 年 9 月（上涨段自 2026 年 6 月起，3 个月+39.5\%），当前处于该段之后的震荡之中。 公司股价较申万农林牧渔指数提前约 23 个月见底，是板块中较早定价周期反转的公司之一。 该子行业的周期来源与猪周期不同（第 \ref{sec:grain} 章、第 \ref{sec:soft} 章），价格更多由自身供给、进口政策与全球定价决定，公司 2026 年 6 月末资产负债率 58.9\%、半年 ROE 0.6\%（未年化），在本轮下行中属于随行业同步调整的一类。 综合区间形态看，该股单段涨跌幅度偏大、以趋势段为主（最大主升 +92.0\%、最大主跌 -39.7\%），年化波动率 28.2\% 与全样本中位数 37.4\% 相比波动小于样本中位数。

\pfl{盈利情况} 2026 年上半年营业收入 14.72 亿元（同比 -2.6\%），归母净利润 0.25 亿元，同比 -11.0\%。 以年报为趋势对照，2023---2025 年营业收入由 40.05 亿元变为 42.42 亿元（两年年均复合增速 +2.9\%），归母净利润由 1.04 亿元变为 0.57 亿元。 按半年报口径，2026 年上半年的 ROE 为 0.6\%（未年化）、销售净利率 1.7\%、毛利率 17.6\%，ROE 在“其他种植业”6 家公司中按数值由高到低排第\zhnumber{4}位。 按同一口径，三级行业“其他种植业”6 家公司的半年 ROE 中位数为 1.70\%，该公司低于其中位数 1.12 个百分点。 公司给出的应对方向是：公司增资后该公司注册资本将变更为 1.00 亿元，以提升马铃薯加工能力与产能规模效益。 综合看，公司在报告期维持盈利（低于全样本半年 ROE 中位数 0.75\%），利润的可持续性取决于其他种植业景气的延续与成本管控的执行。

\begin{center}\small\setlength{\tabcolsep}{3.4pt}
\captionof{table}[亚盛集团（600108）关键财务指标]{亚盛集团（600108）关键财务指标（合并报表口径；两个半年报期均未年化；现金短债比 $=$ 货币资金 $\div$ 短期债务）}
\begin{adjustbox}{max width=\textwidth}
\begin{tabular}{@{}lrrrrr@{}}
\toprule
指标 & 2023 年 & 2024 年 & 2025 年 & 2025 年半年报 & 2026 年半年报 \\
\midrule
营业收入（亿元） & 40.05 & 41.92 & 42.42 & 15.12 & 14.72 \\
归母净利润（亿元） & 1.04 & 0.96 & 0.57 & 0.28 & 0.25 \\
毛利率（\%） & 17.8 & 17.5 & 15.5 & 17.7 & 17.6 \\
ROE（\%） & 2.5 & 2.3 & 1.3 & 0.7 & 0.6 \\
资产负债率（\%） & 53.5 & 55.3 & 56.3 & 55.8 & 58.9 \\
经营活动现金流净额（亿元） & 1.60 & 1.24 & 1.50 & -1.03 & -1.67 \\
货币资金（亿元） & 6.64 & 4.88 & 3.40 & 3.11 & 5.31 \\
有息负债合计（亿元） & 35.45 & 38.78 & 41.09 & 40.65 & 46.72 \\
现金短债比（倍） & 0.54 & 0.35 & 0.12 & 0.20 & 0.20 \\
\bottomrule
\end{tabular}
\end{adjustbox}
\end{center}

\begin{center}
\begin{minipage}{\textwidth}\centering
\begin{tikzpicture}
\begin{groupplot}[
  group style={group size=2 by 1, horizontal sep=0.60cm},
  width=6.85cm, height=4.60cm,
  tick label style={font=\tiny},
  title style={font=\scriptsize\heiti},
  yticklabel style={font=\tiny},
  ylabel style={font=\tiny},
  grid=major, grid style={LightGray, line width=0.3pt},
  axis line style={InkGray, line width=0.5pt},
  enlarge x limits={abs=0.8},
  legend style={font=\scriptsize, at={(0.5,-0.20)}, anchor=north,
                legend columns=2, draw=none, fill=none, column sep=6pt},
]
\nextgroupplot[
  ybar, bar width=4pt,
  ymin=-4.24, ymax=47.51,
  xtick={0,1,2,3,4,5.7},
  xticklabels={2021,2022,2023,2024,2025,2026H1},
  ylabel={亿元},
]
\addplot[fill=AgriGreen!75, draw=AgriGreen, bar shift=-2.6pt] table[x=i, y=rev]{data/clean/profiles/fin_600108.dat};
\addplot[fill=SoftRed!60, draw=SoftRed, bar shift=2.6pt] table[x=i, y=ni]{data/clean/profiles/fin_600108.dat};
\addplot[fill=AgriGreen!30, draw=AgriGreen, pattern=north east lines, bar shift=-2.6pt] table[x=x, y=rev]{data/clean/profiles/finh_600108.dat};
\addplot[fill=SoftRed!20, draw=SoftRed, pattern=north east lines, bar shift=2.6pt] table[x=x, y=ni]{data/clean/profiles/finh_600108.dat};
\legend{营业收入（年/半年）, 归母净利润（年/半年）}
\nextgroupplot[
  ymin=-4.7, ymax=64.8,
  xtick={0,1,2,3,4,5.7},
  xticklabels={2021,2022,2023,2024,2025,2026H1},
  ylabel={\%},
]
\addplot[AgriGold, mark=*, mark size=1.3pt] table[x=i, y=roe]{data/clean/profiles/fin_600108.dat};
\addplot[OilBlue, mark=square*, mark size=1.3pt] table[x=i, y=debt]{data/clean/profiles/fin_600108.dat};
\addplot[only marks, AgriGold, mark=*, mark size=2.0pt] table[x=x, y=roe]{data/clean/profiles/finh_600108.dat};
\addplot[only marks, OilBlue, mark=square*, mark size=2.0pt] table[x=x, y=debt]{data/clean/profiles/finh_600108.dat};
\legend{ROE, 资产负债率}
\end{groupplot}
\end{tikzpicture}
\begin{tikzpicture}
\begin{axis}[
  name=finbot,
  width=13.75cm, height=3.10cm, scale only axis,
  ymin=-4.67, ymax=52.33,
  xtick={0,1,2,3,4,5.7},
  xticklabels={2021,2022,2023,2024,2025,2026H1},
  tick label style={font=\tiny},
  yticklabel style={font=\tiny},
  ylabel={亿元}, ylabel style={font=\tiny},
  ybar, bar width=4pt,
  grid=major, grid style={LightGray, line width=0.3pt},
  axis line style={InkGray, line width=0.5pt},
  enlarge x limits={abs=0.8},
  legend style={font=\scriptsize, at={(0.5,-0.26)}, anchor=north,
                legend columns=2, draw=none, fill=none, column sep=10pt},
]
\addplot[fill=AgriGreen!55, draw=AgriGreen, bar shift=-3.0pt] table[x=i, y=cash]{data/clean/profiles/fin_600108.dat};
\addplot[fill=SoftRed!45, draw=SoftRed, bar shift=3.0pt] table[x=i, y=ibd]{data/clean/profiles/fin_600108.dat};
\addplot[fill=AgriGreen!25, draw=AgriGreen, pattern=north east lines, bar shift=-3.0pt] table[x=x, y=cash]{data/clean/profiles/finh_600108.dat};
\addplot[fill=SoftRed!20, draw=SoftRed, pattern=north east lines, bar shift=3.0pt] table[x=x, y=ibd]{data/clean/profiles/finh_600108.dat};
\legend{货币资金（左轴）, 有息负债（左轴）}
\end{axis}
\begin{axis}[
  at={(finbot.south east)}, anchor=south east,
  width=13.75cm, height=3.10cm, scale only axis,
  axis y line*=right, axis x line*=none, xtick=\empty,
  ymin=-2.78, ymax=3.86,
  tick label style={font=\tiny},
  yticklabel style={font=\tiny}, scaled y ticks=false,
  legend style={font=\scriptsize, at={(0.5,-0.44)}, anchor=north,
                legend columns=1, draw=none, fill=none},
]
\addplot[OilBlue, mark=*, mark size=2.2pt, line width=1.3pt] table[x=i, y=ocf]{data/clean/profiles/fin_600108.dat};
\addplot[only marks, OilBlue, mark=*, mark size=3.2pt] table[x=x, y=ocf]{data/clean/profiles/finh_600108.dat};
\legend{经营活动现金流净额（右轴，亿元，蓝点）}
\end{axis}
\end{tikzpicture}

\captionof{figure}[亚盛集团（600108）近年主要财务指标]{亚盛集团（600108）近年主要财务指标（左上：营业收入与归母净利润；右上：ROE 与资产负债率；下：货币资金、有息负债为左轴柱状，经营活动现金流净额为其右轴的蓝点折线。
2021---2025 年为年报口径，2026H1 为 2026 年半年报/半年末，与其前一年报之间留出间隔、以斜纹或空心点区分；半年报数据未年化）}
\end{minipage}
\end{center}

\pfl{财务分析} 2026 年 6 月末资产负债率 58.9\%，较 2025 年末上升 2.6 个百分点（样本中位数 52.2\%，所处三级行业内按由低到高排第\zhnumber{3}位）。 2026 年 6 月末货币资金 5.31 亿元、短期债务（短期借款与一年内到期非流动负债）27.04 亿元、长期债务 19.68 亿元，有息负债合计 46.72 亿元，现金短债比 0.20 倍——缺口明显。 净有息负债 41.41 亿元，相当于其 2026 年上半年营业收入的 281\%。 流动比率 1.47、速动比率 0.90，短期流动性无明显缺口。 2026 年上半年经营活动现金流净额 -1.67 亿元，净利润为正但经营现金流为负，利润的现金含量偏低。 2026 年 6 月末在建工程 4.70 亿元，较上年末增加 1.00 亿元，仍有在建投入。 综合看，账面货币资金对有息负债与短期债务的覆盖不足，滚动偿付对新增融资的依赖度较高；资产负债率高于样本中位数 6.7 个百分点；有息负债中短期债务占比更高，期限结构仍以滚动续作为主。 公司披露的重整与债务违约风险为：2025 年 6 月，公司对甘肃证券有限公司股权投资成本 2，000.00 万元、甘肃普华甜菊糖有限公司股权投资成本 464.25 万元。

\pfl{重大战略转型} 从公开披露看，公司的战略推进集中在：产能与项目（2 项）：2001 年 11 月，甘肃省财政厅拨付公司国债转贷资金 1，655.00 万元；并购与重组（1 项）：2009 年 3 月，2009-03-09 签署购买资产协议。 公告口径上，近三年（2023 年 9 月---2026 年 9 月）公司披露重大重组与并购类公告 2 条、对外投资与转型类公告 6 条、回购与股权激励类公告 0 条，合计公告 299 条。 第一大行业“农业”的收入占比由 2021 年年报的 81.2\% 变为 57.1\%（24.1 个百分点），收入结构出现再平衡。 综合判断，报告期内公司有 8 条与战略相关的公告落地（重大重组与并购 2 条、对外投资与转型 6 条），资本运作与主业调整之间存在直接关联。

\begin{center}\small\setlength{\tabcolsep}{4pt}
\captionof{table}[亚盛集团（600108）历史融资与资本运作]{亚盛集团（600108）历史融资与资本运作（据公司公告与公开报道整理；金额为披露值，含首次公开发行、再融资、债券、银行借款与重整投资等）}
\begin{adjustbox}{max width=\textwidth}
\begin{tabular}{@{}p{1.45cm}p{2.55cm}p{4.05cm}p{4.95cm}@{}}
\toprule
时间 & 方式 & 规模 & 用途与说明 \\
\midrule
1997-07 & IPO（上海证券交易所） & 公开发行股票 7,\allowbreak{}000 万股 & 以全额预交款、\allowbreak{}比例配售、\allowbreak{}余款即退方式发行；1997-08-18 在上交所挂牌 \\
2009-12 & 发行股份购买资产（定向增发） & 向甘肃省农垦集团有限责任公司非公开发行人民币普通股（A 股）2.96 亿股；同花顺资本运作口径实际募集资金净额 11.57 亿元 & 认购甘肃农垦全资拥有的七家农场全部农业类资产及十宗土地，\allowbreak{}总股本由 14.41 亿股增至 17.37 亿股 \\
2012-05 & 增发 A 股 & 同花顺资本运作口径实际募集资金净额 10.91 亿元（发行起始日 2012-05-08） & 同花顺“资本运作”页披露的历次增发记录 \\
2023-11 & 中期票据（银行间债券市场） & 甘肃亚盛实业（集团）股份有限公司 2023 年度第一期中期票据（乡村振兴）8.00 亿元 & 简称“23 亚盛实业 MTN001（乡村振兴）”，\allowbreak{}200.00 万元 \\
2026-04-30 & 现金分红（2025 年度） & 以 2025-12-31 总股本 19.47 亿股为基数 & 回报股东；2025 年度合并报表归属于母公司股东的净利润 5，\allowbreak{}693.36 万元；截至 2025-12-31 母公司报表期末未分配利润 14.48 亿元；该预案尚需 2025 年年度股东会审议通过 \\
2026-06-30 & 计息负债扩张 & 应付债券期末余额 11.10 亿元 & 补充经营资金缺口；公司说明系报告期内借款增加；一年内到期的应付债券 8.19 亿元 \\
\bottomrule
\end{tabular}
\end{adjustbox}
\end{center}

\pfl{历史融投资情况} 从现金流量表看，2025 年公司取得借款收到的现金 19.74 亿元、偿还债务支付的现金 17.39 亿元，筹资活动现金流净额 0.49 亿元（2023 年为取得借款 26.66 亿元、偿还债务 32.02 亿元）。 2026 年上半年筹资活动现金流净额为 4.90 亿元（取得借款 14.51 亿元、偿还债务 15.78 亿元，未年化），仍处于净融入状态。 2025 年分配股利、利润或偿付利息支付的现金合计 1.42 亿元。 公开披露的融资记录共 6 笔，工具类型以 IPO、发行股份购买资产、定向增发、中期票据为主；金额与用途见下表。 其中规模最大的两笔为：2026 年 6 月，计息负债扩张，应付债券期末余额 11.10 亿元；2023 年 11 月，中期票据，甘肃亚盛实业（集团）股份有限公司 2023 年度第一期中期票据（乡村振兴）8.00 亿元。 股本方面，近三年披露股本变动 3 次；总股本自 2021 年末以来未发生变化（19.47 亿股）。 分红方面，近三年现金分红 4 次、合计每股派现 0.02 元，最近一次实施在 2025 年 12 月。 综合看，公司融资结构上股权与债务融资并举；2025 年筹资活动现金流净额 0.49 亿元，年度内仍为净融入，与前述经营与投资现金流的方向基本一致。

\pfl{未来潜在融资需求分析} 2026 年 6 月末货币资金 5.31 亿元，而短期债务 27.04 亿元（现金短债比 0.20 倍），2026 年上半年经营现金流净额 -1.67 亿元。按“短期债务减去账面货币资金”的滚动偿付口径，静态缺口约 21.73 亿元；上半年盈利或接近盈亏平衡，该缺口不随经营亏损进一步扩大。 公司 2026 年 6 月末资产负债率 58.9\%、2026 年上半年 ROE 0.6\%（未年化），现金短债比 0.20 倍、短债覆盖不足，融资需求首先用于置换与滚动存量债务、补充营运资金，属于“补血型”而非扩张型。 公开披露的后续资金安排线索包括：债务与授信安排：2026 年 6 月，筹资活动现金流量净额 4.90 亿元（同比增长 704.58\%）、应付债券余额增加 174.85\% 公司对外部债务融资依赖上升；2026 年 9 月，公司为子公司提供的最高额信用担保额度于 该日到期。 资金需求还要叠加在建工程：期末在建工程 4.70 亿元、2026 年上半年购建固定资产等支出 1.28 亿元（未年化），这部分支出在周期底部通常会被主动放缓。 综合看，公司的外部资金需求以营运资金周转与在建项目投入为主，滚动偿付口径下的静态缺口约 21.73 亿元，能否落地取决于授信续作与发行窗口的配合。




\subsubsection{诺普信（002215）}
\label{co:002215}

\pfl{公司基本情况} 诺普信成立于 1999 年，由农药制剂龙头转型的作物科学公司，以蓝莓为核心的特色生鲜消费业务营收占比已过半，形成“农药制剂＋特色生鲜消费”双轮驱动；实际控制人卢柏强股权质押规模较大，正推进 14.5 亿元定增投向蓝莓基地与研发中心。 公司于 2008-02-18 在深交所主板首发上市，注册资本 10.04 亿元，总股本 10.05 亿股。 截至 2026 年 9 月 24 日收盘价 9.62 元，流通市值 77.7 亿元，近一年-26.4\%，流通市值在三级行业“其他种植业”的 6 家公司中按由大到小排第\zhnumber{3}位，近一年涨跌幅低于全样本中位数（-10.1\%）16.3 个百分点。 按 2026 年半年报披露的行业构成，收入占比前三位为特色生鲜消费业务 52.2\%；农药制剂业务 38.5\%；田田圈业务 9.3\%。 与 2021 年年报相比，第一大行业口径由“主体业务”（59.4\%）变为“特色生鲜消费业务”（52.2\%）；两期披露的分类口径有调整，占比差异中同时包含分类因素。

\pfl{历史股价} 以 2021 年 1 月为起点、2026 年 9 月为终点，股价由 5.30 元变为 9.62 元，累计 +81.5\%，区间最高 12.91 元（2025 年 9 月）、最低 5.15 元（2021 年 7 月）；当前价相当于区间最高价的 74.5\%，为区间最低价的 1.87 倍。 月度收益的年化波动率 33.8\%，低于全样本中位数 37.4\%，在三级行业“其他种植业”的 6 家公司中波动率由高到低排第\zhnumber{3}位。

\begin{center}
\begin{minipage}{\textwidth}\centering
\begin{tikzpicture}
\begin{axis}[
  width=12.6cm, height=3.9cm, scale only axis,
  xmin=2020.922, xmax=2026.888, ymin=4.79, ymax=13.81,
  xtick={2021,2022,2023,2024,2025,2026},
  yearticks,
  yticklabel style={font=\scriptsize},
  ylabel={元/股}, ylabel style={font=\scriptsize},
  grid=major, grid style={LightGray, line width=0.3pt},
  axis line style={InkGray, line width=0.5pt},
  every axis plot/.append style={line width=0.9pt},
]
\addplot[AgriGreen] table[x=x,y=y]{data/clean/profiles/px_002215.dat};
\addplot[SoftRed, dashed, line width=0.7pt] table[x=x,y=y]{data/clean/profiles/pxzz_002215.dat};
\addplot[only marks, mark=*, mark size=1.1pt, SoftRed] table[x=x,y=y]{data/clean/profiles/pxzz_002215.dat};
\end{axis}
\end{tikzpicture}

\captionof{figure}[诺普信（002215）股价与周期骨架]{诺普信（002215）月度收盘价（元/股）与 ZigZag 周期骨架（阈值 25\%，圆点为周期转折点）}
\end{minipage}
\end{center}

\pfl{过去周期分析} 以 25\% 的 ZigZag 阈值划分，样本区间内股价形成 2 个主升段与 3 个主跌段。 最大主升段出现在 2024 年 7 月 至 2025 年 9 月，14 个月+86.3\%；最大主跌段为 2025 年 9 月 至 2026 年 6 月，9 个月-36.3\%。 最近一次转折出现在 2026 年 6 月（下跌段自 2025 年 9 月起，9 个月-36.3\%），当前处于该段的延续之中。 公司股价较申万农林牧渔指数提前约 58 个月见底，是板块中较早定价周期反转的公司之一。 该子行业的周期来源与猪周期不同（第 \ref{sec:grain} 章、第 \ref{sec:soft} 章），价格更多由自身供给、进口政策与全球定价决定，公司 2026 年上半年实现盈利、半年 ROE 16.8\%（未年化），好于全样本中位数（0.75\%），下行周期对其报表的冲击小于同业。 综合区间形态看，该股单段涨跌幅度偏大、以趋势段为主（最大主升 +86.3\%、最大主跌 -36.3\%），年化波动率 33.8\% 与全样本中位数 37.4\% 相比波动与样本中位数接近。

\pfl{盈利情况} 2026 年上半年营业收入 39.64 亿元（同比 +7.8\%），归母净利润 7.65 亿元，同比 +18.0\%。 以年报为趋势对照，2023---2025 年营业收入由 41.20 亿元变为 57.83 亿元（两年年均复合增速 +18.5\%），归母净利润由 2.36 亿元变为 6.50 亿元。 按半年报口径，2026 年上半年的 ROE 为 16.8\%（未年化）、销售净利率 19.2\%、毛利率 36.8\%，ROE 在“其他种植业”6 家公司中按数值由高到低排第\zhnumber{2}位。 同期全样本半年 ROE 中位数 0.75\%，公司与之相差 16.00 个百分点（略好于中位数公司）。 从公司披露的原因看，业绩变动主要来自：2026 年 5 月，公司回应蓝莓行业价格回落：属规模化发展过程中的正常调整。 综合看，公司在报告期维持盈利（高于全样本半年 ROE 中位数 0.75\%），利润的可持续性取决于其他种植业景气的延续与成本管控的执行。

\begin{center}\small\setlength{\tabcolsep}{3.4pt}
\captionof{table}[诺普信（002215）关键财务指标]{诺普信（002215）关键财务指标（合并报表口径；两个半年报期均未年化；现金短债比 $=$ 货币资金 $\div$ 短期债务）}
\begin{adjustbox}{max width=\textwidth}
\begin{tabular}{@{}lrrrrr@{}}
\toprule
指标 & 2023 年 & 2024 年 & 2025 年 & 2025 年半年报 & 2026 年半年报 \\
\midrule
营业收入（亿元） & 41.20 & 52.88 & 57.83 & 36.79 & 39.64 \\
归母净利润（亿元） & 2.36 & 5.85 & 6.50 & 6.48 & 7.65 \\
毛利率（\%） & 29.5 & 37.7 & 37.4 & 37.7 & 36.8 \\
ROE（\%） & 6.4 & 15.4 & 16.4 & 15.8 & 16.8 \\
资产负债率（\%） & 61.1 & 63.7 & 63.3 & 63.7 & 60.5 \\
经营活动现金流净额（亿元） & 6.40 & 13.06 & 14.75 & 2.54 & 0.92 \\
货币资金（亿元） & 13.22 & 14.38 & 11.79 & 15.11 & 12.33 \\
有息负债合计（亿元） & 45.93 & 51.22 & 52.39 & 60.17 & 58.48 \\
现金短债比（倍） & 0.37 & 0.40 & 0.30 & 0.33 & 0.36 \\
\bottomrule
\end{tabular}
\end{adjustbox}
\end{center}

\begin{center}
\begin{minipage}{\textwidth}\centering
\begin{tikzpicture}
\begin{groupplot}[
  group style={group size=2 by 1, horizontal sep=0.60cm},
  width=6.85cm, height=4.60cm,
  tick label style={font=\tiny},
  title style={font=\scriptsize\heiti},
  yticklabel style={font=\tiny},
  ylabel style={font=\tiny},
  grid=major, grid style={LightGray, line width=0.3pt},
  axis line style={InkGray, line width=0.5pt},
  enlarge x limits={abs=0.8},
  legend style={font=\scriptsize, at={(0.5,-0.20)}, anchor=north,
                legend columns=2, draw=none, fill=none, column sep=6pt},
]
\nextgroupplot[
  ybar, bar width=4pt,
  ymin=-5.78, ymax=64.77,
  xtick={0,1,2,3,4,5.7},
  xticklabels={2021,2022,2023,2024,2025,2026H1},
  ylabel={亿元},
]
\addplot[fill=AgriGreen!75, draw=AgriGreen, bar shift=-2.6pt] table[x=i, y=rev]{data/clean/profiles/fin_002215.dat};
\addplot[fill=SoftRed!60, draw=SoftRed, bar shift=2.6pt] table[x=i, y=ni]{data/clean/profiles/fin_002215.dat};
\addplot[fill=AgriGreen!30, draw=AgriGreen, pattern=north east lines, bar shift=-2.6pt] table[x=x, y=rev]{data/clean/profiles/finh_002215.dat};
\addplot[fill=SoftRed!20, draw=SoftRed, pattern=north east lines, bar shift=2.6pt] table[x=x, y=ni]{data/clean/profiles/finh_002215.dat};
\legend{营业收入（年/半年）, 归母净利润（年/半年）}
\nextgroupplot[
  ymin=-5.1, ymax=70.0,
  xtick={0,1,2,3,4,5.7},
  xticklabels={2021,2022,2023,2024,2025,2026H1},
  ylabel={\%},
]
\addplot[AgriGold, mark=*, mark size=1.3pt] table[x=i, y=roe]{data/clean/profiles/fin_002215.dat};
\addplot[OilBlue, mark=square*, mark size=1.3pt] table[x=i, y=debt]{data/clean/profiles/fin_002215.dat};
\addplot[only marks, AgriGold, mark=*, mark size=2.0pt] table[x=x, y=roe]{data/clean/profiles/finh_002215.dat};
\addplot[only marks, OilBlue, mark=square*, mark size=2.0pt] table[x=x, y=debt]{data/clean/profiles/finh_002215.dat};
\legend{ROE, 资产负债率}
\end{groupplot}
\end{tikzpicture}
\begin{tikzpicture}
\begin{axis}[
  name=finbot,
  width=13.75cm, height=3.10cm, scale only axis,
  ymin=-5.85, ymax=65.50,
  xtick={0,1,2,3,4,5.7},
  xticklabels={2021,2022,2023,2024,2025,2026H1},
  tick label style={font=\tiny},
  yticklabel style={font=\tiny},
  ylabel={亿元}, ylabel style={font=\tiny},
  ybar, bar width=4pt,
  grid=major, grid style={LightGray, line width=0.3pt},
  axis line style={InkGray, line width=0.5pt},
  enlarge x limits={abs=0.8},
  legend style={font=\scriptsize, at={(0.5,-0.26)}, anchor=north,
                legend columns=2, draw=none, fill=none, column sep=10pt},
]
\addplot[fill=AgriGreen!55, draw=AgriGreen, bar shift=-3.0pt] table[x=i, y=cash]{data/clean/profiles/fin_002215.dat};
\addplot[fill=SoftRed!45, draw=SoftRed, bar shift=3.0pt] table[x=i, y=ibd]{data/clean/profiles/fin_002215.dat};
\addplot[fill=AgriGreen!25, draw=AgriGreen, pattern=north east lines, bar shift=-3.0pt] table[x=x, y=cash]{data/clean/profiles/finh_002215.dat};
\addplot[fill=SoftRed!20, draw=SoftRed, pattern=north east lines, bar shift=3.0pt] table[x=x, y=ibd]{data/clean/profiles/finh_002215.dat};
\legend{货币资金（左轴）, 有息负债（左轴）}
\end{axis}
\begin{axis}[
  at={(finbot.south east)}, anchor=south east,
  width=13.75cm, height=3.10cm, scale only axis,
  axis y line*=right, axis x line*=none, xtick=\empty,
  ymin=-3.84, ymax=19.18,
  tick label style={font=\tiny},
  yticklabel style={font=\tiny}, scaled y ticks=false,
  legend style={font=\scriptsize, at={(0.5,-0.44)}, anchor=north,
                legend columns=1, draw=none, fill=none},
]
\addplot[OilBlue, mark=*, mark size=2.2pt, line width=1.3pt] table[x=i, y=ocf]{data/clean/profiles/fin_002215.dat};
\addplot[only marks, OilBlue, mark=*, mark size=3.2pt] table[x=x, y=ocf]{data/clean/profiles/finh_002215.dat};
\legend{经营活动现金流净额（右轴，亿元，蓝点）}
\end{axis}
\end{tikzpicture}

\captionof{figure}[诺普信（002215）近年主要财务指标]{诺普信（002215）近年主要财务指标（左上：营业收入与归母净利润；右上：ROE 与资产负债率；下：货币资金、有息负债为左轴柱状，经营活动现金流净额为其右轴的蓝点折线。
2021---2025 年为年报口径，2026H1 为 2026 年半年报/半年末，与其前一年报之间留出间隔、以斜纹或空心点区分；半年报数据未年化）}
\end{minipage}
\end{center}

\pfl{财务分析} 2026 年 6 月末资产负债率 60.5\%，较 2025 年末下降 2.8 个百分点（样本中位数 52.2\%，所处三级行业内按由低到高排第\zhnumber{4}位）。 2026 年 6 月末货币资金 12.33 亿元、短期债务（短期借款与一年内到期非流动负债）33.99 亿元、长期债务 24.49 亿元，有息负债合计 58.48 亿元，现金短债比 0.36 倍——覆盖偏紧。 净有息负债 46.16 亿元，相当于其 2026 年上半年营业收入的 116\%。 流动比率 1.03、速动比率 0.81，短期偿债能力处于正常区间。 2026 年上半年经营活动现金流净额 0.92 亿元，经营现金流与利润同向为正，内生资金可以支撑运营。 2026 年 6 月末在建工程 1.74 亿元，较上年末增加 1.11 亿元，仍有在建投入。 综合看，账面货币资金对有息负债与短期债务的覆盖不足，滚动偿付对新增融资的依赖度较高；资产负债率高于样本中位数 8.3 个百分点；有息负债中短期债务占比更高，期限结构仍以滚动续作为主。 公司披露的监管与合规风险为：2026 年 9 月，公司下属控股公司涉及 2，000.00 万元侵害植物新品种权诉讼。

\pfl{重大战略转型} 从公开披露看，公司的战略推进集中在：融资与资本运作（1 项）：2025 年 11 月，公司启动定增是基于对蓝莓产业长期发展的信心以及把握时间窗口进行战略扩张的内在需求。 公告口径上，近三年（2023 年 9 月---2026 年 9 月）公司披露重大重组与并购类公告 0 条、对外投资与转型类公告 2 条、回购与股权激励类公告 50 条，合计公告 374 条。 主营结构的口径变化已经落到报表上：2021 年年报的第一大行业为“主体业务”（59.4\%），2026 年半年报为“特色生鲜消费业务”（52.2\%）；公司在这两期的分部划分方式有过调整。 综合判断，报告期内公司有 2 条与战略相关的公告落地（重大重组与并购 0 条、对外投资与转型 2 条），资本运作与主业调整之间存在直接关联。

\begin{center}\small\setlength{\tabcolsep}{4pt}
\captionof{table}[诺普信（002215）历史融资与资本运作]{诺普信（002215）历史融资与资本运作（据公司公告与公开报道整理；金额为披露值，含首次公开发行、再融资、债券、银行借款与重整投资等）}
\begin{adjustbox}{max width=\textwidth}
\begin{tabular}{@{}p{1.45cm}p{2.55cm}p{4.05cm}p{4.95cm}@{}}
\toprule
时间 & 方式 & 规模 & 用途与说明 \\
\midrule
2008-02 & IPO（深圳证券交易所中小企业板） & 发行 3，\allowbreak{}000 万股（网下配售 600 万股、\allowbreak{}网上定价发行 2，\allowbreak{}400 万股） & 发行前总股本 9,\allowbreak{}000.00 万股，\allowbreak{}发行后增加 3,\allowbreak{}000.00 万股 \\
2025 & 债务融资规模 & 2025 年筹资活动现金流入小计 33.44 亿元（上年 26.74 亿元） & 日常经营及蓝莓产业投入；公司层面借款规模较大，\allowbreak{}与实控人质押融资为不同口径 \\
2025-11-26 & 定向增发（向特定对象发行）预案 & 募集资金总额不超过 14.50 亿元；拟发行股份数量不超过本次发行前总股本的 30\%（即不超过 3.02 亿股）；发行对象为不超过 35 名特定投资者 & 蓝莓基地新增扩建项目拟用募集资金 11.00 亿元（投资总额 12.21 亿元）、\allowbreak{}小浆果国际研发中心建设项目 1.50 亿元（投资总额 2.74 亿元）、\allowbreak{}补充流动资金 2.00 亿元；已经公司第七届董事会第九次会议（临时）审议通过 \\
2026-04 & 现金分红（2025 年度） & 以 10.05 亿股为基数，\allowbreak{}向全体股东每 10 股派发现金红利 3.00 元（含税） & 回报股东；公司该利润分配预案 \\
\bottomrule
\end{tabular}
\end{adjustbox}
\end{center}

\pfl{历史融投资情况} 从现金流量表看，2025 年公司取得借款收到的现金 65.86 亿元、偿还债务支付的现金 66.78 亿元、吸收权益性投资 0.02 亿元，筹资活动现金流净额 -4.26 亿元（2023 年为取得借款 69.94 亿元、偿还债务 59.81 亿元）。 2026 年上半年筹资活动现金流净额为 -0.20 亿元（取得借款 44.76 亿元、偿还债务 39.85 亿元，未年化），已转为净流出，处于净还债状态。 2025 年分配股利、利润或偿付利息支付的现金合计 4.29 亿元。 公开披露的融资记录共 4 笔，工具类型以 IPO、债务融资规模、定向增发为主；金额与用途见下表。 其中规模最大的两笔为：2025 年 11 月，定向增发，募集资金总额不超过 14.50 亿元，拟发行股份数量不超过本次发行前总股本的 30\%（即不超过 3.02 亿股），发行对象为不超过 35 名特定投资者。 股本方面，近三年披露股本变动 8 次（激励股份解禁 2 次、股份回购 2 次、增发新股上市 1 次、股权激励 1 次）；总股本由 2021 年末的 9.87 亿股变为 10.05 亿股（+1.81\%）。 分红方面，近三年现金分红 4 次、合计每股派现 0.95 元，最近一次实施在 2025 年 12 月。 近三年“再融资”类公告 7 条，涉及定向增发、募集资金使用。 综合看，公司融资结构上以债务滚动为主、股权融资为补充；2025 年筹资活动现金流净额 -4.26 亿元，融资节奏仍以维持存量债务周转为主，与前述经营与投资现金流的方向基本一致。

\pfl{未来潜在融资需求分析} 2026 年 6 月末货币资金 12.33 亿元，而短期债务 33.99 亿元（现金短债比 0.36 倍），2026 年上半年经营现金流净额 0.92 亿元。按“短期债务减去账面货币资金”的滚动偿付口径，静态缺口约 21.66 亿元；上半年盈利或接近盈亏平衡，该缺口不随经营亏损进一步扩大。 按第 \ref{sec:financing-need} 节的三档划分，公司属于第三档“扩张型”：2026 年上半年 ROE 16.8\%（未年化，折合约 33.5\%）、6 月末资产负债率 60.5\%，资产负债表仍具备加杠杆空间；融资用途以产能扩张、品类并购与海外布局为主，单笔规模通常在 5---20 亿元量级，会被市场定价为成长而非补血。 公开披露的后续资金安排线索包括：股权融资线索：2025 年度，公司向特定对象发行 A 股股票预案；2025 年 11 月，公司在投资者关系活动中说明：本次定增为预案。 资金需求还要叠加在建工程：期末在建工程 1.74 亿元、2026 年上半年购建固定资产等支出 3.71 亿元（未年化），这部分支出在周期底部通常会被主动放缓。 综合看，公司的外部资金需求以营运资金周转与在建项目投入为主，滚动偿付口径下的静态缺口约 21.66 亿元，融资的时点与规模更多取决于信用条件与资本市场窗口。




\subsubsection{农发种业（600313）}
\label{co:600313}

\pfl{公司基本情况} 公司前身可追溯至 1999 年，中国农业发展集团有限公司（中国农发集团）旗下的种业与农资综合企业，主营小麦、玉米、水稻等种子及农药化肥贸易；2026 年 1 月完成向控股股东定增 4.02 亿元，2025 年度实现自 2017 年以来首次现金分红，但母公司累计未分配利润长期为负。 公司于 2001-01-19 在上交所首发上市，注册资本 11.61 亿元，总股本 10.82 亿股。 截至 2026 年 9 月 24 日收盘价 6.80 元，流通市值 73.6 亿元，近一年+7.6\%，流通市值在三级行业“种子”的 8 家公司中按由大到小排第\zhnumber{3}位，近一年涨跌幅高于全样本中位数（-10.1\%）17.7 个百分点。 股权结构方面，实际控制人为中国农业发展集团有限公司。

\pfl{历史股价} 以 2021 年 1 月为起点、2026 年 9 月为终点，股价由 5.09 元变为 6.80 元，累计 +33.6\%，区间最高 15.84 元（2022 年 5 月）、最低 4.41 元（2021 年 6 月）；当前价相当于区间最高价的 42.9\%，为区间最低价的 1.54 倍。 月度收益的年化波动率 49.7\%，高于全样本中位数 37.4\%，在三级行业“种子”的 8 家公司中波动率由高到低排第\zhnumber{2}位。

\begin{center}
\begin{minipage}{\textwidth}\centering
\begin{tikzpicture}
\begin{axis}[
  width=12.6cm, height=3.9cm, scale only axis,
  xmin=2020.922, xmax=2026.888, ymin=4.10, ymax=16.95,
  xtick={2021,2022,2023,2024,2025,2026},
  yearticks,
  yticklabel style={font=\scriptsize},
  ylabel={元/股}, ylabel style={font=\scriptsize},
  grid=major, grid style={LightGray, line width=0.3pt},
  axis line style={InkGray, line width=0.5pt},
  every axis plot/.append style={line width=0.9pt},
]
\addplot[AgriGreen] table[x=x,y=y]{data/clean/profiles/px_600313.dat};
\addplot[SoftRed, dashed, line width=0.7pt] table[x=x,y=y]{data/clean/profiles/pxzz_600313.dat};
\addplot[only marks, mark=*, mark size=1.1pt, SoftRed] table[x=x,y=y]{data/clean/profiles/pxzz_600313.dat};
\end{axis}
\end{tikzpicture}

\captionof{figure}[农发种业（600313）股价与周期骨架]{农发种业（600313）月度收盘价（元/股）与 ZigZag 周期骨架（阈值 25\%，圆点为周期转折点）}
\end{minipage}
\end{center}

\pfl{过去周期分析} 以 25\% 的 ZigZag 阈值划分，样本区间内股价形成 3 个主升段与 3 个主跌段。 最大主升段出现在 2021 年 6 月 至 2022 年 5 月，11 个月+259.2\%；最大主跌段为 2022 年 5 月 至 2024 年 8 月，27 个月-66.0\%。 最近一次转折出现在 2026 年 9 月（上涨段自 2026 年 6 月起，3 个月+29.8\%），当前处于该段之后的震荡之中。 公司股价较申万农林牧渔指数提前约 59 个月见底，是板块中较早定价周期反转的公司之一。 种子行业的价格由制种成本、粮价与政策（转基因商业化、种业振兴）共同决定（第 \ref{sec:planting} 节），与生猪价格的联动有限，公司 2026 年上半年实现盈利、半年 ROE 2.2\%（未年化），好于全样本中位数（0.75\%），下行周期对其报表的冲击小于同业。 综合区间形态看，该股单段涨跌幅度偏大、以趋势段为主（最大主升 +259.2\%、最大主跌 -66.0\%），年化波动率 49.7\% 与全样本中位数 37.4\% 相比波动明显大于样本中位数。

\pfl{盈利情况} 2026 年上半年营业收入 38.73 亿元（同比 +32.8\%），归母净利润 0.53 亿元，同比 -16.6\%。 以年报为趋势对照，2023---2025 年营业收入由 67.61 亿元变为 66.54 亿元（两年年均复合增速 -0.8\%），归母净利润由 1.27 亿元变为 0.84 亿元。 按半年报口径，2026 年上半年的 ROE 为 2.2\%（未年化）、销售净利率 1.7\%、毛利率 5.5\%，ROE 在“种子”8 家公司中按数值由高到低排第\zhnumber{3}位。 按同一口径，三级行业“种子”8 家公司的半年 ROE 中位数为 0.17\%，该公司高于其中位数 2.06 个百分点。 从公司披露的原因看，业绩变动主要来自：2026 年上半年，营业总收入 38.73 亿元（同比 +32.84\%）；2024 年，归母净利润同比减少：受国际市场影响农药下游企业出口订单锐减、行业去库存导致国内市场需求走低、农药主要产品价格下跌。 综合看，公司在报告期维持盈利（高于全样本半年 ROE 中位数 0.75\%），利润的可持续性取决于种子景气的延续与成本管控的执行。

\begin{center}\small\setlength{\tabcolsep}{3.4pt}
\captionof{table}[农发种业（600313）关键财务指标]{农发种业（600313）关键财务指标（合并报表口径；两个半年报期均未年化；现金短债比 $=$ 货币资金 $\div$ 短期债务）}
\begin{adjustbox}{max width=\textwidth}
\begin{tabular}{@{}lrrrrr@{}}
\toprule
指标 & 2023 年 & 2024 年 & 2025 年 & 2025 年半年报 & 2026 年半年报 \\
\midrule
营业收入（亿元） & 67.61 & 55.63 & 66.54 & 29.16 & 38.73 \\
归母净利润（亿元） & 1.27 & 0.47 & 0.84 & 0.63 & 0.53 \\
毛利率（\%） & 8.4 & 7.6 & 7.5 & 7.0 & 5.5 \\
ROE（\%） & 7.1 & 2.5 & 4.3 & 3.2 & 2.2 \\
资产负债率（\%） & 40.5 & 44.4 & 43.3 & 39.5 & 37.2 \\
经营活动现金流净额（亿元） & 6.06 & -0.77 & 3.02 & -1.02 & 0.79 \\
货币资金（亿元） & 12.06 & 9.99 & 15.09 & 9.65 & 8.73 \\
有息负债合计（亿元） & 3.75 & 3.87 & 5.05 & 5.97 & 4.77 \\
现金短债比（倍） & 4.57 & 3.06 & 3.83 & 1.98 & 2.44 \\
\bottomrule
\end{tabular}
\end{adjustbox}
\end{center}

\begin{center}
\begin{minipage}{\textwidth}\centering
\begin{tikzpicture}
\begin{groupplot}[
  group style={group size=2 by 1, horizontal sep=0.60cm},
  width=6.85cm, height=4.60cm,
  tick label style={font=\tiny},
  title style={font=\scriptsize\heiti},
  yticklabel style={font=\tiny},
  ylabel style={font=\tiny},
  grid=major, grid style={LightGray, line width=0.3pt},
  axis line style={InkGray, line width=0.5pt},
  enlarge x limits={abs=0.8},
  legend style={font=\scriptsize, at={(0.5,-0.20)}, anchor=north,
                legend columns=2, draw=none, fill=none, column sep=6pt},
]
\nextgroupplot[
  ybar, bar width=4pt,
  ymin=-6.76, ymax=75.72,
  xtick={0,1,2,3,4,5.7},
  xticklabels={2021,2022,2023,2024,2025,2026H1},
  ylabel={亿元},
]
\addplot[fill=AgriGreen!75, draw=AgriGreen, bar shift=-2.6pt] table[x=i, y=rev]{data/clean/profiles/fin_600313.dat};
\addplot[fill=SoftRed!60, draw=SoftRed, bar shift=2.6pt] table[x=i, y=ni]{data/clean/profiles/fin_600313.dat};
\addplot[fill=AgriGreen!30, draw=AgriGreen, pattern=north east lines, bar shift=-2.6pt] table[x=x, y=rev]{data/clean/profiles/finh_600313.dat};
\addplot[fill=SoftRed!20, draw=SoftRed, pattern=north east lines, bar shift=2.6pt] table[x=x, y=ni]{data/clean/profiles/finh_600313.dat};
\legend{营业收入（年/半年）, 归母净利润（年/半年）}
\nextgroupplot[
  ymin=-3.6, ymax=48.9,
  xtick={0,1,2,3,4,5.7},
  xticklabels={2021,2022,2023,2024,2025,2026H1},
  ylabel={\%},
]
\addplot[AgriGold, mark=*, mark size=1.3pt] table[x=i, y=roe]{data/clean/profiles/fin_600313.dat};
\addplot[OilBlue, mark=square*, mark size=1.3pt] table[x=i, y=debt]{data/clean/profiles/fin_600313.dat};
\addplot[only marks, AgriGold, mark=*, mark size=2.0pt] table[x=x, y=roe]{data/clean/profiles/finh_600313.dat};
\addplot[only marks, OilBlue, mark=square*, mark size=2.0pt] table[x=x, y=debt]{data/clean/profiles/finh_600313.dat};
\legend{ROE, 资产负债率}
\end{groupplot}
\end{tikzpicture}
\begin{tikzpicture}
\begin{axis}[
  name=finbot,
  width=13.75cm, height=3.10cm, scale only axis,
  ymin=-1.51, ymax=16.91,
  xtick={0,1,2,3,4,5.7},
  xticklabels={2021,2022,2023,2024,2025,2026H1},
  tick label style={font=\tiny},
  yticklabel style={font=\tiny},
  ylabel={亿元}, ylabel style={font=\tiny},
  ybar, bar width=4pt,
  grid=major, grid style={LightGray, line width=0.3pt},
  axis line style={InkGray, line width=0.5pt},
  enlarge x limits={abs=0.8},
  legend style={font=\scriptsize, at={(0.5,-0.26)}, anchor=north,
                legend columns=2, draw=none, fill=none, column sep=10pt},
]
\addplot[fill=AgriGreen!55, draw=AgriGreen, bar shift=-3.0pt] table[x=i, y=cash]{data/clean/profiles/fin_600313.dat};
\addplot[fill=SoftRed!45, draw=SoftRed, bar shift=3.0pt] table[x=i, y=ibd]{data/clean/profiles/fin_600313.dat};
\addplot[fill=AgriGreen!25, draw=AgriGreen, pattern=north east lines, bar shift=-3.0pt] table[x=x, y=cash]{data/clean/profiles/finh_600313.dat};
\addplot[fill=SoftRed!20, draw=SoftRed, pattern=north east lines, bar shift=3.0pt] table[x=x, y=ibd]{data/clean/profiles/finh_600313.dat};
\legend{货币资金（左轴）, 有息负债（左轴）}
\end{axis}
\begin{axis}[
  at={(finbot.south east)}, anchor=south east,
  width=13.75cm, height=3.10cm, scale only axis,
  axis y line*=right, axis x line*=none, xtick=\empty,
  ymin=-2.54, ymax=8.10,
  tick label style={font=\tiny},
  yticklabel style={font=\tiny}, scaled y ticks=false,
  legend style={font=\scriptsize, at={(0.5,-0.44)}, anchor=north,
                legend columns=1, draw=none, fill=none},
]
\addplot[OilBlue, mark=*, mark size=2.2pt, line width=1.3pt] table[x=i, y=ocf]{data/clean/profiles/fin_600313.dat};
\addplot[only marks, OilBlue, mark=*, mark size=3.2pt] table[x=x, y=ocf]{data/clean/profiles/finh_600313.dat};
\legend{经营活动现金流净额（右轴，亿元，蓝点）}
\end{axis}
\end{tikzpicture}

\captionof{figure}[农发种业（600313）近年主要财务指标]{农发种业（600313）近年主要财务指标（左上：营业收入与归母净利润；右上：ROE 与资产负债率；下：货币资金、有息负债为左轴柱状，经营活动现金流净额为其右轴的蓝点折线。
2021---2025 年为年报口径，2026H1 为 2026 年半年报/半年末，与其前一年报之间留出间隔、以斜纹或空心点区分；半年报数据未年化）}
\end{minipage}
\end{center}

\pfl{财务分析} 2026 年 6 月末资产负债率 37.2\%，较 2025 年末下降 6.1 个百分点（样本中位数 52.2\%，所处三级行业内按由低到高排第\zhnumber{4}位）。 2026 年 6 月末货币资金 8.73 亿元、短期债务（短期借款与一年内到期非流动负债）3.58 亿元、长期债务 1.19 亿元，有息负债合计 4.77 亿元，现金短债比 2.44 倍——覆盖充分。 净有息负债 -3.96 亿元，相当于其 2026 年上半年营业收入的 -10\%。 流动比率 2.24、速动比率 1.69，短期流动性无明显缺口。 2026 年上半年经营活动现金流净额 0.79 亿元，经营现金流与利润同向为正，内生资金可以支撑运营。 2026 年 6 月末在建工程 1.36 亿元，较上年末减少 0.04 亿元，在建投入在收缩。 2026 年 6 月末商誉 1.36 亿元，占归母净资产的 3.9\%，是净资产中最脆弱的一块。 综合看，现金短债比 2.44 倍，短期偿付压力可控；资产负债率低于样本中位数 15.0 个百分点；有息负债中短期债务占比更高，期限结构仍以滚动续作为主。

\pfl{重大战略转型} 从公开披露看，公司的战略推进集中在：其他动作（2 项）：2026 年 8 月，公告向全资子公司增资，以现金出资、资金来源为自有资金向全资子公司增资；2022 年 8 月，第七届董事会第十九次会议审议通过投资设立控股子公司。 公告口径上，近三年（2023 年 9 月---2026 年 9 月）公司披露重大重组与并购类公告 3 条、对外投资与转型类公告 5 条、回购与股权激励类公告 0 条，合计公告 291 条。 综合判断，报告期内公司有 8 条与战略相关的公告落地（重大重组与并购 3 条、对外投资与转型 5 条），资本运作与主业调整之间存在直接关联。

\begin{center}\small\setlength{\tabcolsep}{4pt}
\captionof{table}[农发种业（600313）历史融资与资本运作]{农发种业（600313）历史融资与资本运作（据公司公告与公开报道整理；金额为披露值，含首次公开发行、再融资、债券、银行借款与重整投资等）}
\begin{adjustbox}{max width=\textwidth}
\begin{tabular}{@{}p{1.45cm}p{2.55cm}p{4.05cm}p{4.95cm}@{}}
\toprule
时间 & 方式 & 规模 & 用途与说明 \\
\midrule
2000-12 & IPO（上海证券交易所主板） & 向社会公开发行人民币普通股 8，\allowbreak{}000 万股，\allowbreak{}每股发行价格 6.30 元，\allowbreak{}发行后总股本 25 & 经中国证监会证监发行字〔2000〕178 号文批准发行 \\
2025 & 公积金弥补亏损 & 使用盈余公积弥补亏损 1,\allowbreak{}563.18 万元 & 为提升投资者回报能力 \\
2025-11 至 2026-01 & 定向增发（向特定对象发行 A 股股票并在主板上市） & 发行 7，\allowbreak{}917.53 万股，\allowbreak{}发行价格 5.14 元/\allowbreak{}股，\allowbreak{}募集资金总额 4.07 亿元 & 将中国农发集团以委托贷款方式拨付公司的国拨资金转为对公司的直接投资（股权投资），\allowbreak{}同时补充流动资金；发行对象为公司实际控制人中国农发集团（拟认购 3.07 亿元）及其一致行动人华农资产（拟认购 1.00 亿元）；定价基准日为第七届董事会第五十次会议决议公告日 \\
2026-06 & 现金分红（2025 年度） & 以股权登记日总股本 11.61 亿股为基数 & 回报股东 \\
\bottomrule
\end{tabular}
\end{adjustbox}
\end{center}

\pfl{历史融投资情况} 从现金流量表看，2025 年公司取得借款收到的现金 6.49 亿元、偿还债务支付的现金 6.15 亿元，筹资活动现金流净额 -0.01 亿元（2023 年为取得借款 3.15 亿元、偿还债务 2.49 亿元）。 2026 年上半年筹资活动现金流净额为 0.54 亿元（取得借款 2.69 亿元、偿还债务 2.74 亿元，未年化），仍处于净融入状态。 2025 年分配股利、利润或偿付利息支付的现金合计 0.29 亿元。 公开披露的融资记录共 4 笔，工具类型以 IPO、定向增发为主；金额与用途见下表。 股本方面，近三年披露股本变动 4 次（限售股份上市 1 次）；总股本自 2021 年末以来未发生变化（10.82 亿股）。 分红方面，近三年现金分红 1 次、合计每股派现 0.01 元，最近一次实施在 2026 年 6 月。 近三年“再融资”类公告 32 条，涉及定向增发、发行股份购买资产、募集资金使用。 综合看，公司融资结构上股权与债务融资并举；2025 年筹资活动现金流净额 -0.01 亿元，融资节奏仍以维持存量债务周转为主，与前述经营与投资现金流的方向基本一致。

\pfl{未来潜在融资需求分析} 2026 年 6 月末货币资金 8.73 亿元，而短期债务 3.58 亿元（现金短债比 2.44 倍），2026 年上半年经营现金流净额 0.79 亿元。账面货币资金可以覆盖短期债务，缺口不体现在静态偿付层面。 公司 2026 年 6 月末资产负债率 37.2\%、2026 年上半年 ROE 2.2\%（未年化），资产负债表尚有余量，融资需求不是“补血型”的刚性需求，而是配套在建项目投入的主动性安排。 公开披露的后续资金安排线索包括：股权融资线索：2026 年 1 月，公司完成向中国农发集团及华农资产定增。 资金需求还要叠加在建工程：期末在建工程 1.36 亿元、2026 年上半年购建固定资产等支出 0.71 亿元（未年化），这部分支出在周期底部通常会被主动放缓。 静态偿付不存在缺口，后续融资更可能服务于产能或并购安排，而不是填补流动性窟窿。




\subsubsection{宏辉果蔬（603336）}
\label{co:603336}

\pfl{公司基本情况} 宏辉果蔬成立于 1992 年，汕头起家的果蔬采后加工与冷链配送服务商，2016 年上交所主板上市；2025 年控制权由创始人黄俊辉转让给苏州国资背景的申泽瑞泰，2026 年以 7 亿元收购施美药业 41.128\% 股权跨界医药，形成果蔬+医药双主业。 公司于 2016-11-24 在上交所首发上市，注册资本 6.09 亿元，总股本 5.70 亿股。 截至 2026 年 9 月 24 日收盘价 10.30 元，流通市值 62.7 亿元，近一年+11.1\%，流通市值在三级行业“其他种植业”的 6 家公司中按由大到小排第\zhnumber{4}位，近一年涨跌幅高于全样本中位数（-10.1\%）21.2 个百分点。 按 2025 年年报披露的行业构成，收入占比前三位为果蔬服务业 72.1\%；供应链管理服务业 22.6\%；农副食品加工业 4.9\%。

\pfl{历史股价} 自 2021 年 1 月至 2026 年 9 月，股价由 10.11 元变为 10.30 元，累计 +1.9\%，区间最高 11.75 元（2026 年 4 月）、最低 3.01 元（2024 年 6 月）；当前价相当于区间最高价的 87.7\%，为区间最低价的 3.42 倍。 月度收益的年化波动率 40.8\%，高于全样本中位数 37.4\%，在三级行业“其他种植业”的 6 家公司中波动率由高到低排第\zhnumber{2}位。 把股价阶段与公司事件对齐看，2024 年 6 月 至 2026 年 4 月股价上涨+290.4\%，同期（2026 年 4 月）第六届董事会第六次会议审议通过收购江西施美药业股份有限公司部分股份，以自有及自筹资金 7.00 亿元受让自然人江鸿持有的施美药业 41.128\% 股份。 整体上看，区间的几处大级别拐点都能在公司的资本运作、股权变动或风险事件上找到对应，股价波动有明显的公司层面驱动，而非单纯的行业 beta。

\begin{center}
\begin{minipage}{\textwidth}\centering
\begin{tikzpicture}
\begin{axis}[
  width=12.6cm, height=3.9cm, scale only axis,
  xmin=2020.922, xmax=2026.888, ymin=2.80, ymax=12.57,
  xtick={2021,2022,2023,2024,2025,2026},
  yearticks,
  yticklabel style={font=\scriptsize},
  ylabel={元/股}, ylabel style={font=\scriptsize},
  grid=major, grid style={LightGray, line width=0.3pt},
  axis line style={InkGray, line width=0.5pt},
  every axis plot/.append style={line width=0.9pt},
]
\addplot[AgriGreen] table[x=x,y=y]{data/clean/profiles/px_603336.dat};
\addplot[SoftRed, dashed, line width=0.7pt] table[x=x,y=y]{data/clean/profiles/pxzz_603336.dat};
\addplot[only marks, mark=*, mark size=1.1pt, SoftRed] table[x=x,y=y]{data/clean/profiles/pxzz_603336.dat};
\end{axis}
\end{tikzpicture}

\captionof{figure}[宏辉果蔬（603336）股价与周期骨架]{宏辉果蔬（603336）月度收盘价（元/股）与 ZigZag 周期骨架（阈值 25\%，圆点为周期转折点）}
\end{minipage}
\end{center}

\pfl{过去周期分析} 以 25\% 的 ZigZag 阈值划分，样本区间内股价形成 3 个主升段与 3 个主跌段。 最大主升段出现在 2024 年 6 月 至 2026 年 4 月，22 个月+290.4\%；最大主跌段为 2021 年 10 月 至 2024 年 6 月，32 个月-62.7\%。 最近一次转折出现在 2026 年 9 月（上涨段自 2026 年 6 月起，3 个月+30.1\%），当前处于该段之后的震荡之中。 公司股价较申万农林牧渔指数提前约 23 个月见底，是板块中较早定价周期反转的公司之一。 该子行业的周期来源与猪周期不同（第 \ref{sec:grain} 章、第 \ref{sec:soft} 章），价格更多由自身供给、进口政策与全球定价决定，公司 2026 年 6 月末资产负债率 36.3\%、半年 ROE 0.4\%（未年化），在本轮下行中属于随行业同步调整的一类。 综合区间形态看，该股单段涨跌幅度偏大、以趋势段为主（最大主升 +290.4\%、最大主跌 -62.7\%），年化波动率 40.8\% 与全样本中位数 37.4\% 相比波动与样本中位数接近。

\pfl{盈利情况} 2026 年上半年营业收入 4.18 亿元（同比 -11.1\%），归母净利润 0.05 亿元，同比 -25.4\%。 以年报为趋势对照，2023---2025 年营业收入由 10.86 亿元变为 10.52 亿元（两年年均复合增速 -1.6\%），归母净利润由 0.24 亿元变为 0.18 亿元。 按半年报口径，2026 年上半年的 ROE 为 0.4\%（未年化）、销售净利率 5.4\%、毛利率 16.2\%，ROE 在“其他种植业”6 家公司中按数值由高到低排第\zhnumber{5}位。 按同一口径，三级行业“其他种植业”6 家公司的半年 ROE 中位数为 1.70\%，该公司低于其中位数 1.33 个百分点。 从公司披露的原因看，业绩变动主要来自：2025 年，同比减少 44.68\%。 面对这一局面，公司披露的举措包括：2026 年 9 月，主动缩减部分低毛利业务。 综合看，公司在报告期维持盈利（低于全样本半年 ROE 中位数 0.75\%），利润的可持续性取决于其他种植业景气的延续与成本管控的执行。

\begin{center}\small\setlength{\tabcolsep}{3.4pt}
\captionof{table}[宏辉果蔬（603336）关键财务指标]{宏辉果蔬（603336）关键财务指标（合并报表口径；两个半年报期均未年化；现金短债比 $=$ 货币资金 $\div$ 短期债务）}
\begin{adjustbox}{max width=\textwidth}
\begin{tabular}{@{}lrrrrr@{}}
\toprule
指标 & 2023 年 & 2024 年 & 2025 年 & 2025 年半年报 & 2026 年半年报 \\
\midrule
营业收入（亿元） & 10.86 & 10.80 & 10.52 & 4.70 & 4.18 \\
归母净利润（亿元） & 0.24 & 0.18 & 0.18 & 0.07 & 0.05 \\
毛利率（\%） & 8.0 & 9.1 & 9.1 & 9.1 & 16.2 \\
ROE（\%） & 2.1 & 1.6 & 1.4 & 0.6 & 0.4 \\
资产负债率（\%） & 37.6 & 39.5 & 19.8 & 38.9 & 36.3 \\
经营活动现金流净额（亿元） & -0.03 & 0.35 & 2.05 & 0.38 & 0.94 \\
货币资金（亿元） & 0.39 & 0.59 & 4.09 & 0.78 & 4.04 \\
有息负债合计（亿元） & 6.23 & 6.80 & 2.99 & 6.79 & 7.43 \\
现金短债比（倍） & 0.10 & 0.13 & 2.17 & 0.18 & 2.37 \\
\bottomrule
\end{tabular}
\end{adjustbox}
\end{center}

\begin{center}
\begin{minipage}{\textwidth}\centering
\begin{tikzpicture}
\begin{groupplot}[
  group style={group size=2 by 1, horizontal sep=0.60cm},
  width=6.85cm, height=4.60cm,
  tick label style={font=\tiny},
  title style={font=\scriptsize\heiti},
  yticklabel style={font=\tiny},
  ylabel style={font=\tiny},
  grid=major, grid style={LightGray, line width=0.3pt},
  axis line style={InkGray, line width=0.5pt},
  enlarge x limits={abs=0.8},
  legend style={font=\scriptsize, at={(0.5,-0.20)}, anchor=north,
                legend columns=2, draw=none, fill=none, column sep=6pt},
]
\nextgroupplot[
  ybar, bar width=4pt,
  ymin=-1.13, ymax=12.69,
  xtick={0,1,2,3,4,5.7},
  xticklabels={2021,2022,2023,2024,2025,2026H1},
  ylabel={亿元},
]
\addplot[fill=AgriGreen!75, draw=AgriGreen, bar shift=-2.6pt] table[x=i, y=rev]{data/clean/profiles/fin_603336.dat};
\addplot[fill=SoftRed!60, draw=SoftRed, bar shift=2.6pt] table[x=i, y=ni]{data/clean/profiles/fin_603336.dat};
\addplot[fill=AgriGreen!30, draw=AgriGreen, pattern=north east lines, bar shift=-2.6pt] table[x=x, y=rev]{data/clean/profiles/finh_603336.dat};
\addplot[fill=SoftRed!20, draw=SoftRed, pattern=north east lines, bar shift=2.6pt] table[x=x, y=ni]{data/clean/profiles/finh_603336.dat};
\legend{营业收入（年/半年）, 归母净利润（年/半年）}
\nextgroupplot[
  ymin=-3.2, ymax=43.5,
  xtick={0,1,2,3,4,5.7},
  xticklabels={2021,2022,2023,2024,2025,2026H1},
  ylabel={\%},
]
\addplot[AgriGold, mark=*, mark size=1.3pt] table[x=i, y=roe]{data/clean/profiles/fin_603336.dat};
\addplot[OilBlue, mark=square*, mark size=1.3pt] table[x=i, y=debt]{data/clean/profiles/fin_603336.dat};
\addplot[only marks, AgriGold, mark=*, mark size=2.0pt] table[x=x, y=roe]{data/clean/profiles/finh_603336.dat};
\addplot[only marks, OilBlue, mark=square*, mark size=2.0pt] table[x=x, y=debt]{data/clean/profiles/finh_603336.dat};
\legend{ROE, 资产负债率}
\end{groupplot}
\end{tikzpicture}
\begin{tikzpicture}
\begin{axis}[
  name=finbot,
  width=13.75cm, height=3.10cm, scale only axis,
  ymin=-0.74, ymax=8.33,
  xtick={0,1,2,3,4,5.7},
  xticklabels={2021,2022,2023,2024,2025,2026H1},
  tick label style={font=\tiny},
  yticklabel style={font=\tiny},
  ylabel={亿元}, ylabel style={font=\tiny},
  ybar, bar width=4pt,
  grid=major, grid style={LightGray, line width=0.3pt},
  axis line style={InkGray, line width=0.5pt},
  enlarge x limits={abs=0.8},
  legend style={font=\scriptsize, at={(0.5,-0.26)}, anchor=north,
                legend columns=2, draw=none, fill=none, column sep=10pt},
]
\addplot[fill=AgriGreen!55, draw=AgriGreen, bar shift=-3.0pt] table[x=i, y=cash]{data/clean/profiles/fin_603336.dat};
\addplot[fill=SoftRed!45, draw=SoftRed, bar shift=3.0pt] table[x=i, y=ibd]{data/clean/profiles/fin_603336.dat};
\addplot[fill=AgriGreen!25, draw=AgriGreen, pattern=north east lines, bar shift=-3.0pt] table[x=x, y=cash]{data/clean/profiles/finh_603336.dat};
\addplot[fill=SoftRed!20, draw=SoftRed, pattern=north east lines, bar shift=3.0pt] table[x=x, y=ibd]{data/clean/profiles/finh_603336.dat};
\legend{货币资金（左轴）, 有息负债（左轴）}
\end{axis}
\begin{axis}[
  at={(finbot.south east)}, anchor=south east,
  width=13.75cm, height=3.10cm, scale only axis,
  axis y line*=right, axis x line*=none, xtick=\empty,
  ymin=-1.02, ymax=2.78,
  tick label style={font=\tiny},
  yticklabel style={font=\tiny}, scaled y ticks=false,
  legend style={font=\scriptsize, at={(0.5,-0.44)}, anchor=north,
                legend columns=1, draw=none, fill=none},
]
\addplot[OilBlue, mark=*, mark size=2.2pt, line width=1.3pt] table[x=i, y=ocf]{data/clean/profiles/fin_603336.dat};
\addplot[only marks, OilBlue, mark=*, mark size=3.2pt] table[x=x, y=ocf]{data/clean/profiles/finh_603336.dat};
\legend{经营活动现金流净额（右轴，亿元，蓝点）}
\end{axis}
\end{tikzpicture}

\captionof{figure}[宏辉果蔬（603336）近年主要财务指标]{宏辉果蔬（603336）近年主要财务指标（左上：营业收入与归母净利润；右上：ROE 与资产负债率；下：货币资金、有息负债为左轴柱状，经营活动现金流净额为其右轴的蓝点折线。
2021---2025 年为年报口径，2026H1 为 2026 年半年报/半年末，与其前一年报之间留出间隔、以斜纹或空心点区分；半年报数据未年化）}
\end{minipage}
\end{center}

\pfl{财务分析} 2026 年 6 月末资产负债率 36.3\%，较 2025 年末上升 16.6 个百分点（样本中位数 52.2\%，所处三级行业内按由低到高排第\zhnumber{1}位）。 2026 年 6 月末货币资金 4.04 亿元、短期债务（短期借款与一年内到期非流动负债）1.71 亿元、长期债务 5.73 亿元，有息负债合计 7.43 亿元，现金短债比 2.37 倍——覆盖充分。 净有息负债 3.40 亿元，相当于其 2026 年上半年营业收入的 81\%。 流动比率 2.96、速动比率 2.70，短期流动性无明显缺口。 2026 年上半年经营活动现金流净额 0.94 亿元，经营现金流与利润同向为正，内生资金可以支撑运营。 2026 年 6 月末商誉 4.07 亿元，占归母净资产的 21.8\%，是净资产中最脆弱的一块。 综合看，现金短债比 2.37 倍，短期偿付压力可控；资产负债率低于样本中位数 15.9 个百分点；有息负债中长期债务占比更高，期限结构对短债滚动的压力较小。 公司披露的股东与质押风险为：2025 年第一季度，报告显示黄俊辉持股 2.52 亿股（44.20\%），其中质押 5，840.00 万股。

\pfl{重大战略转型} 公司的战略动作可归为以下几类：并购与重组（2 项）：2026 年 6 月，支付第三期股权转让款 2.10 亿元（对价的 30\%）；2026 年 4 月，第六届董事会第六次会议审议通过收购江西施美药业股份有限公司部分股份，以自有及自筹资金 7.00 亿元受让自然人江鸿持有的施美药业 41.128\% 股份。 公告口径上，近三年（2023 年 9 月---2026 年 9 月）公司披露重大重组与并购类公告 5 条、对外投资与转型类公告 1 条、回购与股权激励类公告 0 条，合计公告 385 条。 第一大行业“果蔬服务业”的收入占比由 2021 年年报的 96.8\% 变为 72.1\%（24.7 个百分点），收入结构出现再平衡。 综合判断，报告期内公司有 6 条与战略相关的公告落地（重大重组与并购 5 条、对外投资与转型 1 条），资本运作与主业调整之间存在直接关联。

\begin{center}\small\setlength{\tabcolsep}{4pt}
\captionof{table}[宏辉果蔬（603336）历史融资与资本运作]{宏辉果蔬（603336）历史融资与资本运作（据公司公告与公开报道整理；金额为披露值，含首次公开发行、再融资、债券、银行借款与重整投资等）}
\begin{adjustbox}{max width=\textwidth}
\begin{tabular}{@{}p{1.45cm}p{2.55cm}p{4.05cm}p{4.95cm}@{}}
\toprule
时间 & 方式 & 规模 & 用途与说明 \\
\midrule
2016-11 & IPO（上交所主板） & 发行 3，\allowbreak{}335 万股，\allowbreak{}发行价 9.31 元/\allowbreak{}股，\allowbreak{}募集资金总额 3.10 亿元 & 天津、\allowbreak{}上海、\allowbreak{}广州果蔬加工配送基地建设项目及宏辉果蔬信息化系统建设项目；占发行后总股本 25.01\% \\
2025-06 & 控股股东协议转让（控制权变更，非上市公司融资） & 申泽瑞泰以 5.68 元/\allowbreak{}股受让黄俊辉持有的 1.51 亿股（26.54\%） & 股权转让款由转让方取得，\allowbreak{}不进入上市公司；转让后黄俊辉持股由 2.52 亿股（44.19\%）降至 1.01 亿股（17.66\%） \\
2026-04 & 对外投资（自有及自筹资金） & 7.00 亿元受让江西施美药业股份有限公司 41.128\% 股份 & 布局医药第二增长曲线，\allowbreak{}取得施美药业实质管理控制权与并表；标的 100\% 股权作价 17.02 亿元 \\
\bottomrule
\end{tabular}
\end{adjustbox}
\end{center}

\pfl{历史融投资情况} 从现金流量表看，2025 年公司取得借款收到的现金 3.71 亿元、偿还债务支付的现金 5.18 亿元、吸收权益性投资 0.01 亿元，筹资活动现金流净额 -1.67 亿元（2023 年为取得借款 4.43 亿元、偿还债务 3.89 亿元）。 2026 年上半年筹资活动现金流净额为 4.36 亿元（取得借款 5.47 亿元、偿还债务 1.06 亿元，未年化），仍处于净融入状态。 2025 年分配股利、利润或偿付利息支付的现金合计 0.21 亿元。 公开披露的融资记录共 3 笔，工具类型以 IPO 为主；金额与用途见下表。 其中规模最大的两笔为：2016 年 11 月，IPO，发行 3，335 万股，发行价 9.31 元/股，募集资金总额 3.10 亿元；2026 年 4 月，对外投资，7.00 亿元受让江西施美药业股份有限公司 41.128\% 股份。 股本方面，近三年披露股本变动 10 次（可转债转股 9 次、送股 1 次、转增 1 次）；总股本由 2021 年末的 4.39 亿股变为 5.70 亿股（+30.02\%）。 分红方面，近三年现金分红 3 次、合计每股派现 0.05 元，最近一次实施在 2025 年 12 月。 近三年“再融资”类公告 3 条，涉及可转债、募集资金使用。 综合看，公司融资结构上以债务滚动为主、股权融资为补充；2025 年筹资活动现金流净额 -1.67 亿元，融资节奏仍以维持存量债务周转为主，与前述经营与投资现金流的方向基本一致。

\pfl{未来潜在融资需求分析} 2026 年 6 月末货币资金 4.04 亿元，而短期债务 1.71 亿元（现金短债比 2.37 倍），2026 年上半年经营现金流净额 0.94 亿元。账面货币资金可以覆盖短期债务，缺口不体现在静态偿付层面。 公司 2026 年 6 月末资产负债率 36.3\%、2026 年上半年 ROE 0.4\%（未年化），资产端与负债端都没有明显压力，融资需求以相机抉择的扩张型用途为主。 公开披露的后续资金安排线索包括：债务与授信安排：2026 年度，公司及子（孙）公司预计互为提供总额不超过 16.00 亿元的担保，用于满足经营与融资需要。 以现有货币资金与经营现金流衡量，公司不存在刚性补血的紧迫性，融资动作更多是主动型的。




\subsubsection{神农种业（300189）}
\label{co:300189}

\pfl{公司基本情况} 神农种业成立于 2000 年，海南起家的水稻、玉米、油菜等农作物种子企业，2011 年创业板上市（原名神农大丰，历经神农基因、神农科技更名）；因信披违法被证监会处罚、主业长期亏损，2025 年靠处置联营企业股权实现账面盈利，2026 年上半年再度亏损。 公司于 2011-03-16 在深交所创业板通过重大资产重组（借壳）上市，注册资本 10.24 亿元，总股本 10.24 亿股。 截至 2026 年 9 月 24 日收盘价 6.17 元，流通市值 54.7 亿元，近一年+41.2\%，流通市值在三级行业“种子”的 8 家公司中按由大到小排第\zhnumber{4}位，近一年涨跌幅高于全样本中位数（-10.1\%）51.3 个百分点。 按 2026 年半年报披露的行业构成，收入占比前两位为农业 95.6\%；其他(补充) 4.4\%。

\pfl{历史股价} 自 2021 年 1 月至 2026 年 9 月，股价由 4.15 元变为 6.17 元，累计 +48.7\%，区间最高 8.66 元（2026 年 1 月）、最低 1.79 元（2024 年 6 月）；当前价相当于区间最高价的 71.2\%，为区间最低价的 3.45 倍。 月度收益的年化波动率 44.7\%，高于全样本中位数 37.4\%，在三级行业“种子”的 8 家公司中波动率由高到低排第\zhnumber{3}位。 把股价阶段与公司事件对齐看，2024 年 6 月 至 2026 年 1 月股价上涨+383.8\%，同期（2025 年 1 月）公司转让控股子公司福建神农大丰种业科技有限公司 88.67\% 股权给李坤泰。 把这些阶段与公司事件对齐后可以看到，几轮大涨大跌多由公司自身的资本运作与风险事件触发，行业景气只解释了其中一部分。

\begin{center}
\begin{minipage}{\textwidth}\centering
\begin{tikzpicture}
\begin{axis}[
  width=12.6cm, height=3.9cm, scale only axis,
  xmin=2020.922, xmax=2026.888, ymin=1.66, ymax=9.27,
  xtick={2021,2022,2023,2024,2025,2026},
  yearticks,
  yticklabel style={font=\scriptsize},
  ylabel={元/股}, ylabel style={font=\scriptsize},
  grid=major, grid style={LightGray, line width=0.3pt},
  axis line style={InkGray, line width=0.5pt},
  every axis plot/.append style={line width=0.9pt},
]
\addplot[AgriGreen] table[x=x,y=y]{data/clean/profiles/px_300189.dat};
\addplot[SoftRed, dashed, line width=0.7pt] table[x=x,y=y]{data/clean/profiles/pxzz_300189.dat};
\addplot[only marks, mark=*, mark size=1.1pt, SoftRed] table[x=x,y=y]{data/clean/profiles/pxzz_300189.dat};
\end{axis}
\end{tikzpicture}

\captionof{figure}[神农种业（300189）股价与周期骨架]{神农种业（300189）月度收盘价（元/股）与 ZigZag 周期骨架（阈值 25\%，圆点为周期转折点）}
\end{minipage}
\end{center}

\pfl{过去周期分析} 以 25\% 的 ZigZag 阈值划分，样本区间内股价形成 3 个主升段与 2 个主跌段。 最大主升段出现在 2024 年 6 月 至 2026 年 1 月，19 个月+383.8\%；最大主跌段为 2022 年 5 月 至 2024 年 6 月，25 个月-68.2\%。 最近一次转折出现在 2026 年 9 月（上涨段自 2026 年 6 月起，3 个月+42.8\%），当前处于该段之后的震荡之中。 公司股价较申万农林牧渔指数提前约 23 个月见底，是板块中较早定价周期反转的公司之一。 种子行业的价格由制种成本、粮价与政策（转基因商业化、种业振兴）共同决定（第 \ref{sec:planting} 节），与生猪价格的联动有限，公司 2026 年 6 月末资产负债率 39.8\%、半年 ROE -1.9\%（未年化），在本轮下行中属于随行业同步调整的一类。 综合区间形态看，该股单段涨跌幅度偏大、以趋势段为主（最大主升 +383.8\%、最大主跌 -68.2\%），年化波动率 44.7\% 与全样本中位数 37.4\% 相比波动明显大于样本中位数。

\pfl{盈利情况} 2026 年上半年营业收入 0.56 亿元（同比 -35.3\%），归母净利润 -0.13 亿元，亏损同比收窄 0.01 亿元。 以年报为趋势对照，2023---2025 年营业收入由 1.67 亿元变为 2.49 亿元（两年年均复合增速 +22.2\%），归母净利润由 -0.36 亿元变为 1.07 亿元。 按半年报口径，2026 年上半年的 ROE 为 -1.9\%（未年化）、销售净利率 -25.9\%、毛利率 38.4\%，ROE 在“种子”8 家公司中按数值由高到低排第\zhnumber{6}位。 同期全样本半年 ROE 中位数 0.75\%，公司与之相差 2.70 个百分点（弱于中位数公司）。 从公司披露的原因看，业绩变动主要来自：2026 年上半年，毛利率-60.37\% 主要系水产种源受赤潮影响造成损失。 综合看，公司仍处亏损区间、亏损同比收窄，半年 ROE 在三级行业内按由高到低排第\zhnumber{6}位，单位盈利能力的修复依赖行业景气与自身成本端的改善，报表弹性主要体现为价格与销量的方向性变化。

\begin{center}\small\setlength{\tabcolsep}{3.4pt}
\captionof{table}[神农种业（300189）关键财务指标]{神农种业（300189）关键财务指标（合并报表口径；两个半年报期均未年化；现金短债比 $=$ 货币资金 $\div$ 短期债务）}
\begin{adjustbox}{max width=\textwidth}
\begin{tabular}{@{}lrrrrr@{}}
\toprule
指标 & 2023 年 & 2024 年 & 2025 年 & 2025 年半年报 & 2026 年半年报 \\
\midrule
营业收入（亿元） & 1.67 & 1.57 & 2.49 & 0.87 & 0.56 \\
归母净利润（亿元） & -0.36 & -0.51 & 1.07 & -0.14 & -0.13 \\
毛利率（\%） & 35.6 & 37.5 & 35.4 & 30.2 & 38.4 \\
ROE（\%） & -4.6 & -6.8 & 14.0 & -2.0 & -1.9 \\
资产负债率（\%） & 14.5 & 31.5 & 38.9 & 33.4 & 39.8 \\
经营活动现金流净额（亿元） & 0.14 & 0.07 & 0.23 & -0.33 & -0.32 \\
货币资金（亿元） & 1.11 & 0.57 & 0.97 & 0.54 & 0.99 \\
有息负债合计（亿元） & 0.08 & 2.23 & 2.83 & 2.42 & 3.05 \\
现金短债比（倍） & 2,779.09 & 3.57 & 1.63 & 3.84 & 1.11 \\
\bottomrule
\end{tabular}
\end{adjustbox}
\end{center}

\begin{center}
\begin{minipage}{\textwidth}\centering
\begin{tikzpicture}
\begin{groupplot}[
  group style={group size=2 by 1, horizontal sep=0.60cm},
  width=6.85cm, height=4.60cm,
  tick label style={font=\tiny},
  title style={font=\scriptsize\heiti},
  yticklabel style={font=\tiny},
  ylabel style={font=\tiny},
  grid=major, grid style={LightGray, line width=0.3pt},
  axis line style={InkGray, line width=0.5pt},
  enlarge x limits={abs=0.8},
  legend style={font=\scriptsize, at={(0.5,-0.20)}, anchor=north,
                legend columns=2, draw=none, fill=none, column sep=6pt},
]
\nextgroupplot[
  ybar, bar width=4pt,
  ymin=-0.96, ymax=2.87,
  xtick={0,1,2,3,4,5.7},
  xticklabels={2021,2022,2023,2024,2025,2026H1},
  ylabel={亿元},
]
\addplot[fill=AgriGreen!75, draw=AgriGreen, bar shift=-2.6pt] table[x=i, y=rev]{data/clean/profiles/fin_300189.dat};
\addplot[fill=SoftRed!60, draw=SoftRed, bar shift=2.6pt] table[x=i, y=ni]{data/clean/profiles/fin_300189.dat};
\addplot[fill=AgriGreen!30, draw=AgriGreen, pattern=north east lines, bar shift=-2.6pt] table[x=x, y=rev]{data/clean/profiles/finh_300189.dat};
\addplot[fill=SoftRed!20, draw=SoftRed, pattern=north east lines, bar shift=2.6pt] table[x=x, y=ni]{data/clean/profiles/finh_300189.dat};
\legend{营业收入（年/半年）, 归母净利润（年/半年）}
\nextgroupplot[
  ymin=-11.5, ymax=44.6,
  xtick={0,1,2,3,4,5.7},
  xticklabels={2021,2022,2023,2024,2025,2026H1},
  ylabel={\%},
]
\addplot[AgriGold, mark=*, mark size=1.3pt] table[x=i, y=roe]{data/clean/profiles/fin_300189.dat};
\addplot[OilBlue, mark=square*, mark size=1.3pt] table[x=i, y=debt]{data/clean/profiles/fin_300189.dat};
\addplot[only marks, AgriGold, mark=*, mark size=2.0pt] table[x=x, y=roe]{data/clean/profiles/finh_300189.dat};
\addplot[only marks, OilBlue, mark=square*, mark size=2.0pt] table[x=x, y=debt]{data/clean/profiles/finh_300189.dat};
\legend{ROE, 资产负债率}
\end{groupplot}
\end{tikzpicture}
\begin{tikzpicture}
\begin{axis}[
  name=finbot,
  width=13.75cm, height=3.10cm, scale only axis,
  ymin=-0.30, ymax=3.42,
  xtick={0,1,2,3,4,5.7},
  xticklabels={2021,2022,2023,2024,2025,2026H1},
  tick label style={font=\tiny},
  yticklabel style={font=\tiny},
  ylabel={亿元}, ylabel style={font=\tiny},
  ybar, bar width=4pt,
  grid=major, grid style={LightGray, line width=0.3pt},
  axis line style={InkGray, line width=0.5pt},
  enlarge x limits={abs=0.8},
  legend style={font=\scriptsize, at={(0.5,-0.26)}, anchor=north,
                legend columns=2, draw=none, fill=none, column sep=10pt},
]
\addplot[fill=AgriGreen!55, draw=AgriGreen, bar shift=-3.0pt] table[x=i, y=cash]{data/clean/profiles/fin_300189.dat};
\addplot[fill=SoftRed!45, draw=SoftRed, bar shift=3.0pt] table[x=i, y=ibd]{data/clean/profiles/fin_300189.dat};
\addplot[fill=AgriGreen!25, draw=AgriGreen, pattern=north east lines, bar shift=-3.0pt] table[x=x, y=cash]{data/clean/profiles/finh_300189.dat};
\addplot[fill=SoftRed!20, draw=SoftRed, pattern=north east lines, bar shift=3.0pt] table[x=x, y=ibd]{data/clean/profiles/finh_300189.dat};
\legend{货币资金（左轴）, 有息负债（左轴）}
\end{axis}
\begin{axis}[
  at={(finbot.south east)}, anchor=south east,
  width=13.75cm, height=3.10cm, scale only axis,
  axis y line*=right, axis x line*=none, xtick=\empty,
  ymin=-0.51, ymax=0.40,
  tick label style={font=\tiny},
  yticklabel style={font=\tiny}, scaled y ticks=false,
  legend style={font=\scriptsize, at={(0.5,-0.44)}, anchor=north,
                legend columns=1, draw=none, fill=none},
]
\addplot[OilBlue, mark=*, mark size=2.2pt, line width=1.3pt] table[x=i, y=ocf]{data/clean/profiles/fin_300189.dat};
\addplot[only marks, OilBlue, mark=*, mark size=3.2pt] table[x=x, y=ocf]{data/clean/profiles/finh_300189.dat};
\legend{经营活动现金流净额（右轴，亿元，蓝点）}
\end{axis}
\end{tikzpicture}

\captionof{figure}[神农种业（300189）近年主要财务指标]{神农种业（300189）近年主要财务指标（左上：营业收入与归母净利润；右上：ROE 与资产负债率；下：货币资金、有息负债为左轴柱状，经营活动现金流净额为其右轴的蓝点折线。
2021---2025 年为年报口径，2026H1 为 2026 年半年报/半年末，与其前一年报之间留出间隔、以斜纹或空心点区分；半年报数据未年化）}
\end{minipage}
\end{center}

\pfl{财务分析} 2026 年 6 月末资产负债率 39.8\%，较 2025 年末上升 0.9 个百分点（样本中位数 52.2\%，所处三级行业内按由低到高排第\zhnumber{6}位）。 2026 年 6 月末货币资金 0.99 亿元、短期债务（短期借款与一年内到期非流动负债）0.88 亿元、长期债务 2.17 亿元，有息负债合计 3.05 亿元，现金短债比 1.11 倍——覆盖充分。 净有息负债 2.06 亿元，相当于其 2026 年上半年营业收入的 366\%。 流动比率 1.79、速动比率 1.44，短期流动性无明显缺口。 2026 年上半年经营活动现金流净额 -0.32 亿元，经营现金流与利润同时为负，现金消耗较快。 2026 年 6 月末在建工程 2.71 亿元，较上年末增加 0.09 亿元，仍有在建投入。 综合看，现金短债比 1.11 倍，短期偿付压力可控；资产负债率低于样本中位数 12.4 个百分点；有息负债中长期债务占比更高，期限结构对短债滚动的压力较小。

\pfl{重大战略转型} 公司的战略动作可归为以下几类：其他动作（2 项）：2025 年，公司投资收益 1.41 亿元；并购与重组（1 项）：2025 年 1 月，公司转让控股子公司福建神农大丰种业科技有限公司 88.67\% 股权给李坤泰。 公告口径上，近三年（2023 年 9 月---2026 年 9 月）公司披露重大重组与并购类公告 1 条、对外投资与转型类公告 2 条、回购与股权激励类公告 0 条，合计公告 265 条。 第一大行业仍是“农业”，收入占比 95.6\%，与 2021 年年报相比没有方向性变化。 综合判断，报告期内公司有 3 条与战略相关的公告落地（重大重组与并购 1 条、对外投资与转型 2 条），资本运作与主业调整之间存在直接关联。

\begin{center}\small\setlength{\tabcolsep}{4pt}
\captionof{table}[神农种业（300189）历史融资与资本运作]{神农种业（300189）历史融资与资本运作（据公司公告与公开报道整理；金额为披露值，含首次公开发行、再融资、债券、银行借款与重整投资等）}
\begin{adjustbox}{max width=\textwidth}
\begin{tabular}{@{}p{1.45cm}p{2.55cm}p{4.05cm}p{4.95cm}@{}}
\toprule
时间 & 方式 & 规模 & 用途与说明 \\
\midrule
2011-03 & IPO（深交所创业板） & 发行 4，\allowbreak{}000 万股，\allowbreak{}发行价 24.00 元/\allowbreak{}股，\allowbreak{}募集资金总额 9.60 亿元 & 与主营业务相关的项目及主营业务发展所需营运资金；发行后股份总数由 12，\allowbreak{}000 万股增至 16 \\
2015 & 发行股份购买资产（拟） & 拟购买波莲基因 61.52\% 股权 & 购买标的公司股权；交易文件认定不构成重大资产重组、\allowbreak{}不构成借壳上市 \\
2019 & 行政处罚（罚款） & 公司被罚 60.00 万元，\allowbreak{}黄培劲被罚 70.00 万元 & 海南监管局〔2019〕4 号行政处罚决定书，\allowbreak{}非融资事项，\allowbreak{}列此以说明资本运作约束 \\
2024 & 债务融资 & 2024 年度筹资活动产生的现金流量净额为 2.12 亿元，\allowbreak{}同比增加 1908.04\% & 补充运营资金；2025 年财务费用 316.53 万元 \\
\bottomrule
\end{tabular}
\end{adjustbox}
\end{center}

\pfl{历史融投资情况} 从现金流量表看，2025 年公司取得借款收到的现金 0.83 亿元、偿还债务支付的现金 0.21 亿元，筹资活动现金流净额 0.65 亿元（2023 年为取得借款 0.05 亿元、偿还债务 0.14 亿元）。 2026 年上半年筹资活动现金流净额为 0.03 亿元（取得借款 0.28 亿元、偿还债务 0.17 亿元，未年化），仍处于净融入状态。 2025 年分配股利、利润或偿付利息支付的现金合计 0.08 亿元。 公开披露的融资记录共 4 笔，工具类型以 IPO、发行股份购买资产、债务融资为主；金额与用途见下表。 其中规模最大的两笔为：2024 年，债务融资，2024 年度筹资活动产生的现金流量净额为 2.12 亿元，同比增加 1908.04\%；2019 年，行政处罚，公司被罚 60.00 万元，黄培劲被罚 70.00 万元。 股本方面，近三年披露股本变动 3 次；总股本自 2021 年末以来未发生变化（10.24 亿股）。 分红方面，近三年未实施现金分红，在全部样本中属于分红能力偏弱的一端。 近三年“再融资”类公告 14 条，涉及定向增发、募集资金使用。 综合看，公司融资结构上股权与债务融资并举；2025 年筹资活动现金流净额 0.65 亿元，年度内仍为净融入，与前述经营与投资现金流的方向基本一致。

\pfl{未来潜在融资需求分析} 2026 年 6 月末货币资金 0.99 亿元，而短期债务 0.88 亿元（现金短债比 1.11 倍），2026 年上半年经营现金流净额 -0.32 亿元。账面货币资金可以覆盖短期债务，缺口不体现在静态偿付层面。 公司 2026 年上半年归母净利润为负（-0.13 亿元），但 6 月末资产负债率 39.8\% 尚未进入第一档（亏损且负债率 $>$65\%）的区间，当前融资需求以补充流动资金、支撑低谷期经营性支出为主。 公开披露的后续资金安排线索包括：债务与授信安排：2025 年，主要因本期借款增加、利息支出上升，对外部债务资金依赖上升。 资金需求还要叠加在建工程：期末在建工程 2.71 亿元、2026 年上半年购建固定资产等支出 0.10 亿元（未年化），这部分支出在周期底部通常会被主动放缓。 静态偿付不存在缺口，后续融资更可能服务于产能或并购安排，而不是填补流动性窟窿。




\subsubsection{敦煌种业（600354）}
\label{co:600354}

\pfl{公司基本情况} 敦煌种业成立于 1998 年，甘肃酒泉的种业与食品加工企业，2004 年上市的老牌种业公司；2018—2019 年连续亏损曾被实施退市风险警示（*ST 敦种），2021 年撤销；2024 年控股股东两度无偿划转、最终变更为酒钢集团，2025—2026 年种子主业量利齐升，但母公司累计未分配利润仍为大额负数，长期不能分红。 公司于 2004-01-15 在上交所首发上市，注册资本 5.28 亿元，总股本 5.28 亿股。 截至 2026 年 9 月 24 日收盘价 9.53 元，流通市值 50.3 亿元，近一年+53.0\%，流通市值在三级行业“种子”的 8 家公司中按由大到小排第\zhnumber{5}位，近一年涨跌幅高于全样本中位数（-10.1\%）63.1 个百分点。 按 2025 年年报披露的行业构成，收入占比前三位为种子 77.1\%；食品与贸易 16.5\%；其他 6.3\%。

\pfl{历史股价} 以 2021 年 1 月为起点、2026 年 9 月为终点，股价由 4.44 元变为 9.53 元，累计 +114.6\%，区间最高 9.53 元（2026 年 9 月）、最低 3.95 元（2021 年 3 月）；当前价相当于区间最高价的 100.0\%，为区间最低价的 2.41 倍。 月度收益的年化波动率 38.5\%，高于全样本中位数 37.4\%，在三级行业“种子”的 8 家公司中波动率由高到低排第\zhnumber{5}位。 把股价阶段与公司事件对齐看，2021 年 3 月 至 2023 年 1 月股价上涨+110.1\%，同期（2021 年 5 月）公司撤销退市风险警示。 从时间对应关系看，公司层面的资本运作与股权、风险事件解释了区间内多数急涨急跌，与行业指数的同步性相对有限。

\begin{center}
\begin{minipage}{\textwidth}\centering
\begin{tikzpicture}
\begin{axis}[
  width=12.6cm, height=3.9cm, scale only axis,
  xmin=2020.922, xmax=2026.888, ymin=3.67, ymax=10.20,
  xtick={2021,2022,2023,2024,2025,2026},
  yearticks,
  yticklabel style={font=\scriptsize},
  ylabel={元/股}, ylabel style={font=\scriptsize},
  grid=major, grid style={LightGray, line width=0.3pt},
  axis line style={InkGray, line width=0.5pt},
  every axis plot/.append style={line width=0.9pt},
]
\addplot[AgriGreen] table[x=x,y=y]{data/clean/profiles/px_600354.dat};
\addplot[SoftRed, dashed, line width=0.7pt] table[x=x,y=y]{data/clean/profiles/pxzz_600354.dat};
\addplot[only marks, mark=*, mark size=1.1pt, SoftRed] table[x=x,y=y]{data/clean/profiles/pxzz_600354.dat};
\end{axis}
\end{tikzpicture}

\captionof{figure}[敦煌种业（600354）股价与周期骨架]{敦煌种业（600354）月度收盘价（元/股）与 ZigZag 周期骨架（阈值 25\%，圆点为周期转折点）}
\end{minipage}
\end{center}

\pfl{过去周期分析} 以 25\% 的 ZigZag 阈值划分，样本区间内股价形成 3 个主升段与 3 个主跌段。 最大主跌段为 2023 年 1 月 至 2024 年 8 月，19 个月-43.0\%；最大主升段出现在 2021 年 3 月 至 2023 年 1 月，22 个月+110.1\%。 最近一次转折出现在 2026 年 9 月（上涨段自 2026 年 6 月起，3 个月+100.2\%），当前处于该段之后的震荡之中。 公司股价较申万农林牧渔指数提前约 62 个月见底，是板块中较早定价周期反转的公司之一。 种子行业的价格由制种成本、粮价与政策（转基因商业化、种业振兴）共同决定（第 \ref{sec:planting} 节），与生猪价格的联动有限，公司 2026 年上半年实现盈利、半年 ROE 18.7\%（未年化），好于全样本中位数（0.75\%），下行周期对其报表的冲击小于同业。 综合区间形态看，该股单段涨跌幅度偏大、以趋势段为主（最大主升 +110.1\%、最大主跌 -43.0\%），年化波动率 38.5\% 与全样本中位数 37.4\% 相比波动与样本中位数接近。

\pfl{盈利情况} 2026 年上半年营业收入 8.68 亿元（同比 +20.9\%），归母净利润 1.45 亿元，同比 +165.7\%。 以年报为趋势对照，2023---2025 年营业收入由 11.58 亿元变为 11.61 亿元（两年年均复合增速 +0.1\%），归母净利润由 0.41 亿元变为 0.47 亿元。 按半年报口径，2026 年上半年的 ROE 为 18.7\%（未年化）、销售净利率 33.2\%、毛利率 48.6\%，ROE 在“种子”8 家公司中按数值由高到低排第\zhnumber{1}位。 同期全样本半年 ROE 中位数 0.75\%，公司与之相差 17.95 个百分点（略好于中位数公司）。 公司给出的应对方向是：2026 年 8 月，依托精准市场定位与多渠道推广，自有核心优势品种销量稳步攀升。 综合看，公司在报告期维持盈利（高于全样本半年 ROE 中位数 0.75\%），利润的可持续性取决于种子景气的延续与成本管控的执行。

\begin{center}\small\setlength{\tabcolsep}{3.4pt}
\captionof{table}[敦煌种业（600354）关键财务指标]{敦煌种业（600354）关键财务指标（合并报表口径；两个半年报期均未年化；现金短债比 $=$ 货币资金 $\div$ 短期债务）}
\begin{adjustbox}{max width=\textwidth}
\begin{tabular}{@{}lrrrrr@{}}
\toprule
指标 & 2023 年 & 2024 年 & 2025 年 & 2025 年半年报 & 2026 年半年报 \\
\midrule
营业收入（亿元） & 11.58 & 11.56 & 11.61 & 7.18 & 8.68 \\
归母净利润（亿元） & 0.41 & 0.50 & 0.47 & 0.54 & 1.45 \\
毛利率（\%） & 27.8 & 32.2 & 36.6 & 40.6 & 48.6 \\
ROE（\%） & 7.1 & 8.0 & 7.0 & 8.0 & 18.7 \\
资产负债率（\%） & 59.7 & 58.4 & 57.3 & 44.9 & 37.5 \\
经营活动现金流净额（亿元） & 2.43 & 2.77 & 4.22 & 0.07 & 0.20 \\
货币资金（亿元） & 7.84 & 8.39 & 10.20 & 7.05 & 8.05 \\
有息负债合计（亿元） & 4.65 & 3.14 & 1.94 & 2.66 & 0.65 \\
现金短债比（倍） & 1.87 & 3.08 & 6.34 & 3.16 & 24.20 \\
\bottomrule
\end{tabular}
\end{adjustbox}
\end{center}

\begin{center}
\begin{minipage}{\textwidth}\centering
\begin{tikzpicture}
\begin{groupplot}[
  group style={group size=2 by 1, horizontal sep=0.60cm},
  width=6.85cm, height=4.60cm,
  tick label style={font=\tiny},
  title style={font=\scriptsize\heiti},
  yticklabel style={font=\tiny},
  ylabel style={font=\tiny},
  grid=major, grid style={LightGray, line width=0.3pt},
  axis line style={InkGray, line width=0.5pt},
  enlarge x limits={abs=0.8},
  legend style={font=\scriptsize, at={(0.5,-0.20)}, anchor=north,
                legend columns=2, draw=none, fill=none, column sep=6pt},
]
\nextgroupplot[
  ybar, bar width=4pt,
  ymin=-1.16, ymax=13.00,
  xtick={0,1,2,3,4,5.7},
  xticklabels={2021,2022,2023,2024,2025,2026H1},
  ylabel={亿元},
]
\addplot[fill=AgriGreen!75, draw=AgriGreen, bar shift=-2.6pt] table[x=i, y=rev]{data/clean/profiles/fin_600354.dat};
\addplot[fill=SoftRed!60, draw=SoftRed, bar shift=2.6pt] table[x=i, y=ni]{data/clean/profiles/fin_600354.dat};
\addplot[fill=AgriGreen!30, draw=AgriGreen, pattern=north east lines, bar shift=-2.6pt] table[x=x, y=rev]{data/clean/profiles/finh_600354.dat};
\addplot[fill=SoftRed!20, draw=SoftRed, pattern=north east lines, bar shift=2.6pt] table[x=x, y=ni]{data/clean/profiles/finh_600354.dat};
\legend{营业收入（年/半年）, 归母净利润（年/半年）}
\nextgroupplot[
  ymin=-5.0, ymax=68.8,
  xtick={0,1,2,3,4,5.7},
  xticklabels={2021,2022,2023,2024,2025,2026H1},
  ylabel={\%},
]
\addplot[AgriGold, mark=*, mark size=1.3pt] table[x=i, y=roe]{data/clean/profiles/fin_600354.dat};
\addplot[OilBlue, mark=square*, mark size=1.3pt] table[x=i, y=debt]{data/clean/profiles/fin_600354.dat};
\addplot[only marks, AgriGold, mark=*, mark size=2.0pt] table[x=x, y=roe]{data/clean/profiles/finh_600354.dat};
\addplot[only marks, OilBlue, mark=square*, mark size=2.0pt] table[x=x, y=debt]{data/clean/profiles/finh_600354.dat};
\legend{ROE, 资产负债率}
\end{groupplot}
\end{tikzpicture}
\begin{tikzpicture}
\begin{axis}[
  name=finbot,
  width=13.75cm, height=3.10cm, scale only axis,
  ymin=-1.02, ymax=11.43,
  xtick={0,1,2,3,4,5.7},
  xticklabels={2021,2022,2023,2024,2025,2026H1},
  tick label style={font=\tiny},
  yticklabel style={font=\tiny},
  ylabel={亿元}, ylabel style={font=\tiny},
  ybar, bar width=4pt,
  grid=major, grid style={LightGray, line width=0.3pt},
  axis line style={InkGray, line width=0.5pt},
  enlarge x limits={abs=0.8},
  legend style={font=\scriptsize, at={(0.5,-0.26)}, anchor=north,
                legend columns=2, draw=none, fill=none, column sep=10pt},
]
\addplot[fill=AgriGreen!55, draw=AgriGreen, bar shift=-3.0pt] table[x=i, y=cash]{data/clean/profiles/fin_600354.dat};
\addplot[fill=SoftRed!45, draw=SoftRed, bar shift=3.0pt] table[x=i, y=ibd]{data/clean/profiles/fin_600354.dat};
\addplot[fill=AgriGreen!25, draw=AgriGreen, pattern=north east lines, bar shift=-3.0pt] table[x=x, y=cash]{data/clean/profiles/finh_600354.dat};
\addplot[fill=SoftRed!20, draw=SoftRed, pattern=north east lines, bar shift=3.0pt] table[x=x, y=ibd]{data/clean/profiles/finh_600354.dat};
\legend{货币资金（左轴）, 有息负债（左轴）}
\end{axis}
\begin{axis}[
  at={(finbot.south east)}, anchor=south east,
  width=13.75cm, height=3.10cm, scale only axis,
  axis y line*=right, axis x line*=none, xtick=\empty,
  ymin=-1.10, ymax=5.49,
  tick label style={font=\tiny},
  yticklabel style={font=\tiny}, scaled y ticks=false,
  legend style={font=\scriptsize, at={(0.5,-0.44)}, anchor=north,
                legend columns=1, draw=none, fill=none},
]
\addplot[OilBlue, mark=*, mark size=2.2pt, line width=1.3pt] table[x=i, y=ocf]{data/clean/profiles/fin_600354.dat};
\addplot[only marks, OilBlue, mark=*, mark size=3.2pt] table[x=x, y=ocf]{data/clean/profiles/finh_600354.dat};
\legend{经营活动现金流净额（右轴，亿元，蓝点）}
\end{axis}
\end{tikzpicture}

\captionof{figure}[敦煌种业（600354）近年主要财务指标]{敦煌种业（600354）近年主要财务指标（左上：营业收入与归母净利润；右上：ROE 与资产负债率；下：货币资金、有息负债为左轴柱状，经营活动现金流净额为其右轴的蓝点折线。
2021---2025 年为年报口径，2026H1 为 2026 年半年报/半年末，与其前一年报之间留出间隔、以斜纹或空心点区分；半年报数据未年化）}
\end{minipage}
\end{center}

\pfl{财务分析} 2026 年 6 月末资产负债率 37.5\%，较 2025 年末下降 19.9 个百分点（样本中位数 52.2\%，所处三级行业内按由低到高排第\zhnumber{5}位）。 2026 年 6 月末货币资金 8.05 亿元、短期债务（短期借款与一年内到期非流动负债）0.33 亿元、长期债务 0.32 亿元，有息负债合计 0.65 亿元，现金短债比 24.20 倍——覆盖充分。 净有息负债 -7.39 亿元，相当于其 2026 年上半年营业收入的 -85\%。 流动比率 2.70、速动比率 2.14，短期指标尚可。 2026 年上半年经营活动现金流净额 0.20 亿元，经营现金流与利润同向为正，内生资金可以支撑运营。 综合看，现金短债比 24.20 倍，短期偿付压力可控；资产负债率低于样本中位数 14.7 个百分点；有息负债中短期债务占比更高，期限结构仍以滚动续作为主。

\pfl{重大战略转型} 近三年公司的战略动作集中在：其他动作（3 项）：2025 年 4 月，与酒钢集团财务有限公司签订金融服务协议，由酒钢财务公司提供存款、贷款及其他综合信贷、结算等服务；2024 年 9 月，预计截至 2024 年末与关联方酒钢集团及其所属公司发生日常关联交易不超过 2，910.00 万元；融资与资本运作（2 项）：2026 年，第一次临时股东会通过：为酒泉敦煌种业百佳食品有限公司贷款担保 6，000.00 万元、为瓜州敦种棉业有限公司贷款担保 5；2026 年 4 月，公司担保总额不高于 1.20 亿元。 公告口径上，近三年（2023 年 9 月---2026 年 9 月）公司披露重大重组与并购类公告 7 条、对外投资与转型类公告 0 条、回购与股权激励类公告 1 条，合计公告 281 条。 第一大行业“种子”的收入占比由 2021 年年报的 65.7\% 变为 77.1\%（11.5 个百分点），收入结构出现再平衡。 综合判断，报告期内公司有 7 条与战略相关的公告落地（重大重组与并购 7 条、对外投资与转型 0 条），资本运作与主业调整之间存在直接关联。

\begin{center}\small\setlength{\tabcolsep}{4pt}
\captionof{table}[敦煌种业（600354）历史融资与资本运作]{敦煌种业（600354）历史融资与资本运作（据公司公告与公开报道整理；金额为披露值，含首次公开发行、再融资、债券、银行借款与重整投资等）}
\begin{adjustbox}{max width=\textwidth}
\begin{tabular}{@{}p{1.45cm}p{2.55cm}p{4.05cm}p{4.95cm}@{}}
\toprule
时间 & 方式 & 规模 & 用途与说明 \\
\midrule
2003-12 & IPO（上交所，二级市场投资者定价市值配售） & 发行 7，\allowbreak{}500 万股，\allowbreak{}发行价 5.68 元/\allowbreak{}股，\allowbreak{}募集资金 42，\allowbreak{}600 万元 & 上市募集资金用于公司主业项目建设；发行数量占发行后总股本 40.33\% \\
2011-02-11 & 非公开发行（向5名特定投资者） & 1,\allowbreak{}758 万股，\allowbreak{}发行价 25.00 元/\allowbreak{}股，\allowbreak{}募集资金净额 4.10 亿元 & 玉米种子烘干生产线建设项目 4，\allowbreak{}900.00 万元、\allowbreak{}购买新疆玛纳斯油脂资产 5 \\
2014-10 & 非公开发行预案 & 拟向特定对象非公开发行 8，\allowbreak{}000 万股，\allowbreak{}发行价格 6.00 元/\allowbreak{}股 & 全部用于补充流动资金；该方案后续实施为 2015 年 10 月的 8，\allowbreak{}000 万股非公开发行 \\
2015-01-13 & 中期票据 & 15 敦煌种业 MTN001，\allowbreak{}发行额度 3.00 亿元 & 偿还银行贷款、\allowbreak{}补充流动资金；报告期内已按募集资金要求使用 \\
2015-10-14 & 非公开发行（向特定投资者） & 8，\allowbreak{}000 万股 & 所募集资金全部用于采购种子等原材料；完成后注册资本变更为 5.28 亿元 \\
2025-04 & 银行综合授信额度（年度议案） & 拟向银行申请不超过人民币 10.60 亿元综合授信额度 & 流动资金贷款等综合授信品种，\allowbreak{}担保总额不高于 1.20 亿元 \\
2026-04-28 & 银行综合授信额度（年度议案） & 拟向银行及各金融机构申请不超过人民币 12.00 亿元的综合授信额度 & 流动资金贷款、\allowbreak{}项目并购贷款、\allowbreak{}中长期贷款、\allowbreak{}银行承兑汇票、\allowbreak{}商业承兑汇票、\allowbreak{}保函、\allowbreak{}信用证、\allowbreak{}抵押贷款、\allowbreak{}担保贷款、\allowbreak{}融资租赁、\allowbreak{}网络供应链融资、\allowbreak{}商业信用等；期限自 2025 年年度股东会审议通过之日起至 2026 年年度股东会召开之日止；授信额度可循环使用 \\
\bottomrule
\end{tabular}
\end{adjustbox}
\end{center}

\pfl{历史融投资情况} 从现金流量表看，2025 年公司取得借款收到的现金 2.64 亿元、偿还债务支付的现金 3.73 亿元，筹资活动现金流净额 -1.88 亿元（2023 年为取得借款 6.28 亿元、偿还债务 6.97 亿元）。 2026 年上半年筹资活动现金流净额为 -2.27 亿元（取得借款 1.31 亿元、偿还债务 2.54 亿元，未年化），已转为净流出，处于净还债状态。 2025 年分配股利、利润或偿付利息支付的现金合计 0.73 亿元。 公开披露的融资记录共 7 笔，工具类型以 IPO、定向增发、中期票据、银行授信为主；金额与用途见下表。 其中规模最大的两笔为：2026 年 4 月，银行授信，拟向银行及各金融机构申请不超过人民币 12.00 亿元的综合授信额度；2025 年 4 月，银行授信，拟向银行申请不超过人民币 10.60 亿元综合授信额度。 股本方面，近三年披露股本变动 3 次；总股本自 2021 年末以来未发生变化（5.28 亿股）。 分红方面，近三年未实施现金分红，在全部样本中属于分红能力偏弱的一端。 近三年“再融资”类公告 1 条，涉及定向增发。 综合看，公司融资结构上股权与债务融资并举；2025 年筹资活动现金流净额 -1.88 亿元，融资节奏仍以维持存量债务周转为主，与前述经营与投资现金流的方向基本一致。

\pfl{未来潜在融资需求分析} 2026 年 6 月末货币资金 8.05 亿元，而短期债务 0.33 亿元（现金短债比 24.20 倍），2026 年上半年经营现金流净额 0.20 亿元。账面货币资金可以覆盖短期债务，缺口不体现在静态偿付层面。 按第 \ref{sec:financing-need} 节的三档划分，公司属于第三档“扩张型”：2026 年上半年 ROE 18.7\%（未年化，折合约 37.4\%）、6 月末资产负债率 37.5\%，资产负债表仍具备加杠杆空间；融资用途以产能扩张、品类并购与海外布局为主，单笔规模通常在 5---20 亿元量级，会被市场定价为成长而非补血。 公开披露的后续资金安排线索包括：债务与授信安排：2026 年度，银行综合授信不超过 12.00 亿元、为子公司担保总额不高于 1.20 亿元继续以银行信贷为主要外部资金来源；资本开支安排：2025 年度，政府补助相关递延收益转入 360.85 万元（与资产相关）、商务局外经贸发展专项资金 309.35 万元、核心种源繁育基地建设项目补助资金 120.00 万元。 以现有货币资金与经营现金流衡量，公司不存在刚性补血的紧迫性，融资动作更多是主动型的。




\subsubsection{ST荃银（300087）}
\label{co:300087}

\pfl{公司基本情况} 合肥的水稻、玉米、小麦种子龙头，2010 年创业板上市；2021 年起由先正达集团系的中化现代农业、中种集团控股，2026 年 1 月中种集团完成部分要约收购后持股 40.51\%；2025 年出现上市以来首次年度亏损，2026 年 6 月因 2024 年年报虚假记载被实施其他风险警示并受行政处罚。 公司于 2010-05-26 在深交所创业板首发上市，注册资本 9.47 亿元，总股本 9.47 亿股。 截至 2026 年 9 月 24 日收盘价 5.61 元，流通市值 49.9 亿元，近一年-40.0\%，流通市值在三级行业“种子”的 8 家公司中按由大到小排第\zhnumber{6}位，近一年涨跌幅低于全样本中位数（-10.1\%）29.9 个百分点。 按 2026 年半年报披露的行业构成，收入占比前两位为农业 97.7\%；其他(补充) 2.3\%。 股权结构方面，控股股东为中种集团控股股东为先正达集团股份有限公司；实际控制人为国务院国资委。

\pfl{历史股价} 本报告样本区间（2021 年 1 月 至 2026 年 9 月）内，股价由 31.89 元变为 5.61 元，累计 -82.4\%，区间最高 35.18 元（2021 年 8 月）、最低 5.12 元（2026 年 7 月）；当前价相当于区间最高价的 15.9\%，为区间最低价的 1.10 倍。 月度收益的年化波动率 57.2\%，高于全样本中位数 37.4\%，在三级行业“种子”的 8 家公司中波动率由高到低排第\zhnumber{1}位。 把股价阶段与公司事件对齐看，2023 年 1 月 至 2024 年 8 月股价下跌-64.8\%，同期 2026 年上半年，公告显示该笔收购最终交易对价为 1.83 亿元。 从时间对应关系看，公司层面的资本运作与股权、风险事件解释了区间内多数急涨急跌，与行业指数的同步性相对有限。

\begin{center}
\begin{minipage}{\textwidth}\centering
\begin{tikzpicture}
\begin{axis}[
  width=12.6cm, height=3.9cm, scale only axis,
  xmin=2020.922, xmax=2026.888, ymin=4.76, ymax=37.64,
  xtick={2021,2022,2023,2024,2025,2026},
  yearticks,
  yticklabel style={font=\scriptsize},
  ylabel={元/股}, ylabel style={font=\scriptsize},
  grid=major, grid style={LightGray, line width=0.3pt},
  axis line style={InkGray, line width=0.5pt},
  every axis plot/.append style={line width=0.9pt},
]
\addplot[AgriGreen] table[x=x,y=y]{data/clean/profiles/px_300087.dat};
\addplot[SoftRed, dashed, line width=0.7pt] table[x=x,y=y]{data/clean/profiles/pxzz_300087.dat};
\addplot[only marks, mark=*, mark size=1.1pt, SoftRed] table[x=x,y=y]{data/clean/profiles/pxzz_300087.dat};
\end{axis}
\end{tikzpicture}

\captionof{figure}[ST荃银（300087）股价与周期骨架]{ST荃银（300087）月度收盘价（元/股）与 ZigZag 周期骨架（阈值 25\%，圆点为周期转折点）}
\end{minipage}
\end{center}

\pfl{过去周期分析} 以 25\% 的 ZigZag 阈值划分，样本区间内股价形成 4 个主升段与 5 个主跌段。 最大主跌段为 2023 年 1 月 至 2024 年 8 月，19 个月-64.8\%；最大主升段出现在 2024 年 8 月 至 2024 年 10 月，2 个月+121.5\%。 最近一次转折出现在 2026 年 7 月（下跌段自 2025 年 12 月起，7 个月-56.2\%），当前处于该段的延续之中。 与申万农林牧渔指数的最近一次底部相比，该公司见底晚约 2 个月，在本轮下行中属于调整更慢的一类。 种子行业的价格由制种成本、粮价与政策（转基因商业化、种业振兴）共同决定（第 \ref{sec:planting} 节），与生猪价格的联动有限，公司 2026 年 6 月末资产负债率 64.8\%、半年 ROE -3.3\%（未年化），在本轮下行中属于随行业同步调整的一类。 综合区间形态看，该股单段涨跌幅度偏大、以趋势段为主（最大主升 +121.5\%、最大主跌 -64.8\%），年化波动率 57.2\% 与全样本中位数 37.4\% 相比波动明显大于样本中位数。

\pfl{盈利情况} 2026 年上半年营业收入 11.79 亿元（同比 -17.9\%），归母净利润 -0.55 亿元，亏损同比扩大 0.14 亿元。 以年报为趋势对照，2023---2025 年营业收入由 41.03 亿元变为 44.95 亿元（两年年均复合增速 +4.7\%），归母净利润由 2.68 亿元变为 -2.12 亿元。 按半年报口径，2026 年上半年的 ROE 为 -3.3\%（未年化）、销售净利率 -7.3\%、毛利率 21.1\%，ROE 在“种子”8 家公司中按数值由高到低排第\zhnumber{7}位。 同期全样本半年 ROE 中位数 0.75\%，公司与之相差 4.09 个百分点（弱于中位数公司）。 从公司披露的原因看，业绩变动主要来自：2026 年 4 月，公司业绩预告显示 2025 年归母净利润亏损 1.80 亿元—2.70 亿元；2026 年 6 月，导致少计提信用减值损失 1，871.51 万元、虚增利润总额 1，871.51 万元。 综合看，公司仍处亏损区间、亏损同比扩大，半年 ROE 在三级行业内按由高到低排第\zhnumber{7}位，单位盈利能力的修复依赖行业景气与自身成本端的改善，报表弹性主要体现为价格与销量的方向性变化。

\begin{center}\small\setlength{\tabcolsep}{3.4pt}
\captionof{table}[ST荃银（300087）关键财务指标]{ST荃银（300087）关键财务指标（合并报表口径；两个半年报期均未年化；现金短债比 $=$ 货币资金 $\div$ 短期债务）}
\begin{adjustbox}{max width=\textwidth}
\begin{tabular}{@{}lrrrrr@{}}
\toprule
指标 & 2023 年 & 2024 年 & 2025 年 & 2025 年半年报 & 2026 年半年报 \\
\midrule
营业收入（亿元） & 41.03 & 47.09 & 44.95 & 14.36 & 11.79 \\
归母净利润（亿元） & 2.68 & 0.97 & -2.12 & -0.41 & -0.55 \\
毛利率（\%） & 26.9 & 25.0 & 21.5 & 13.6 & 21.1 \\
ROE（\%） & 15.0 & 5.0 & -11.6 & -2.1 & -3.3 \\
资产负债率（\%） & 61.6 & 65.9 & 69.0 & 62.3 & 64.8 \\
经营活动现金流净额（亿元） & 2.84 & 0.65 & 1.99 & -7.17 & -7.03 \\
货币资金（亿元） & 15.31 & 16.00 & 16.14 & 8.37 & 7.67 \\
有息负债合计（亿元） & 14.78 & 21.70 & 22.88 & 21.90 & 22.07 \\
现金短债比（倍） & 1.85 & 0.91 & 1.01 & 0.62 & 0.40 \\
\bottomrule
\end{tabular}
\end{adjustbox}
\end{center}

\begin{center}
\begin{minipage}{\textwidth}\centering
\begin{tikzpicture}
\begin{groupplot}[
  group style={group size=2 by 1, horizontal sep=0.60cm},
  width=6.85cm, height=4.60cm,
  tick label style={font=\tiny},
  title style={font=\scriptsize\heiti},
  yticklabel style={font=\tiny},
  ylabel style={font=\tiny},
  grid=major, grid style={LightGray, line width=0.3pt},
  axis line style={InkGray, line width=0.5pt},
  enlarge x limits={abs=0.8},
  legend style={font=\scriptsize, at={(0.5,-0.20)}, anchor=north,
                legend columns=2, draw=none, fill=none, column sep=6pt},
]
\nextgroupplot[
  ybar, bar width=4pt,
  ymin=-7.04, ymax=52.99,
  xtick={0,1,2,3,4,5.7},
  xticklabels={2021,2022,2023,2024,2025,2026H1},
  ylabel={亿元},
]
\addplot[fill=AgriGreen!75, draw=AgriGreen, bar shift=-2.6pt] table[x=i, y=rev]{data/clean/profiles/fin_300087.dat};
\addplot[fill=SoftRed!60, draw=SoftRed, bar shift=2.6pt] table[x=i, y=ni]{data/clean/profiles/fin_300087.dat};
\addplot[fill=AgriGreen!30, draw=AgriGreen, pattern=north east lines, bar shift=-2.6pt] table[x=x, y=rev]{data/clean/profiles/finh_300087.dat};
\addplot[fill=SoftRed!20, draw=SoftRed, pattern=north east lines, bar shift=2.6pt] table[x=x, y=ni]{data/clean/profiles/finh_300087.dat};
\legend{营业收入（年/半年）, 归母净利润（年/半年）}
\nextgroupplot[
  ymin=-18.1, ymax=77.1,
  xtick={0,1,2,3,4,5.7},
  xticklabels={2021,2022,2023,2024,2025,2026H1},
  ylabel={\%},
]
\addplot[AgriGold, mark=*, mark size=1.3pt] table[x=i, y=roe]{data/clean/profiles/fin_300087.dat};
\addplot[OilBlue, mark=square*, mark size=1.3pt] table[x=i, y=debt]{data/clean/profiles/fin_300087.dat};
\addplot[only marks, AgriGold, mark=*, mark size=2.0pt] table[x=x, y=roe]{data/clean/profiles/finh_300087.dat};
\addplot[only marks, OilBlue, mark=square*, mark size=2.0pt] table[x=x, y=debt]{data/clean/profiles/finh_300087.dat};
\legend{ROE, 资产负债率}
\end{groupplot}
\end{tikzpicture}
\begin{tikzpicture}
\begin{axis}[
  name=finbot,
  width=13.75cm, height=3.10cm, scale only axis,
  ymin=-2.29, ymax=25.63,
  xtick={0,1,2,3,4,5.7},
  xticklabels={2021,2022,2023,2024,2025,2026H1},
  tick label style={font=\tiny},
  yticklabel style={font=\tiny},
  ylabel={亿元}, ylabel style={font=\tiny},
  ybar, bar width=4pt,
  grid=major, grid style={LightGray, line width=0.3pt},
  axis line style={InkGray, line width=0.5pt},
  enlarge x limits={abs=0.8},
  legend style={font=\scriptsize, at={(0.5,-0.26)}, anchor=north,
                legend columns=2, draw=none, fill=none, column sep=10pt},
]
\addplot[fill=AgriGreen!55, draw=AgriGreen, bar shift=-3.0pt] table[x=i, y=cash]{data/clean/profiles/fin_300087.dat};
\addplot[fill=SoftRed!45, draw=SoftRed, bar shift=3.0pt] table[x=i, y=ibd]{data/clean/profiles/fin_300087.dat};
\addplot[fill=AgriGreen!25, draw=AgriGreen, pattern=north east lines, bar shift=-3.0pt] table[x=x, y=cash]{data/clean/profiles/finh_300087.dat};
\addplot[fill=SoftRed!20, draw=SoftRed, pattern=north east lines, bar shift=3.0pt] table[x=x, y=ibd]{data/clean/profiles/finh_300087.dat};
\legend{货币资金（左轴）, 有息负债（左轴）}
\end{axis}
\begin{axis}[
  at={(finbot.south east)}, anchor=south east,
  width=13.75cm, height=3.10cm, scale only axis,
  axis y line*=right, axis x line*=none, xtick=\empty,
  ymin=-9.99, ymax=7.77,
  tick label style={font=\tiny},
  yticklabel style={font=\tiny}, scaled y ticks=false,
  legend style={font=\scriptsize, at={(0.5,-0.44)}, anchor=north,
                legend columns=1, draw=none, fill=none},
]
\addplot[OilBlue, mark=*, mark size=2.2pt, line width=1.3pt] table[x=i, y=ocf]{data/clean/profiles/fin_300087.dat};
\addplot[only marks, OilBlue, mark=*, mark size=3.2pt] table[x=x, y=ocf]{data/clean/profiles/finh_300087.dat};
\legend{经营活动现金流净额（右轴，亿元，蓝点）}
\end{axis}
\end{tikzpicture}

\captionof{figure}[ST荃银（300087）近年主要财务指标]{ST荃银（300087）近年主要财务指标（左上：营业收入与归母净利润；右上：ROE 与资产负债率；下：货币资金、有息负债为左轴柱状，经营活动现金流净额为其右轴的蓝点折线。
2021---2025 年为年报口径，2026H1 为 2026 年半年报/半年末，与其前一年报之间留出间隔、以斜纹或空心点区分；半年报数据未年化）}
\end{minipage}
\end{center}

\pfl{财务分析} 2026 年 6 月末资产负债率 64.8\%，较 2025 年末下降 4.2 个百分点（样本中位数 52.2\%，所处三级行业内按由低到高排第\zhnumber{8}位）。 2026 年 6 月末货币资金 7.67 亿元、短期债务（短期借款与一年内到期非流动负债）19.06 亿元、长期债务 3.01 亿元，有息负债合计 22.07 亿元，现金短债比 0.40 倍——覆盖偏紧。 净有息负债 14.40 亿元，相当于其 2026 年上半年营业收入的 122\%。 流动比率 1.29、速动比率 0.55，短期偿债能力处于正常区间。 2026 年上半年经营活动现金流净额 -7.03 亿元，经营现金流与利润同时为负，日常运营对债务与股东投入的依赖度较高。 2026 年 6 月末在建工程 0.53 亿元，较上年末减少 0.95 亿元，在建投入在收缩。 综合看，账面货币资金对有息负债与短期债务的覆盖不足，滚动偿付对新增融资的依赖度较高；资产负债率高于样本中位数 12.6 个百分点；有息负债中短期债务占比更高，期限结构仍以滚动续作为主。 公司披露的监管与合规风险为：2026 年 1 月，收到中国证监会立案告知书，扣非净利润亏损 2.45 亿元—3.50 亿元。

\pfl{重大战略转型} 公司的战略动作可归为以下几类：并购与重组（1 项）：2021 年 1 月，方式包括资产重组、业务调整、委托管理、设立合资公司等。 公告口径上，近三年（2023 年 9 月---2026 年 9 月）公司披露重大重组与并购类公告 18 条、对外投资与转型类公告 0 条、回购与股权激励类公告 8 条，合计公告 318 条。 第一大行业仍是“农业”，收入占比 97.7\%，与 2021 年年报相比没有方向性变化。 综合判断，报告期内公司有 18 条与战略相关的公告落地（重大重组与并购 18 条、对外投资与转型 0 条），资本运作与主业调整之间存在直接关联。

\begin{center}\small\setlength{\tabcolsep}{4pt}
\captionof{table}[ST荃银（300087）历史融资与资本运作]{ST荃银（300087）历史融资与资本运作（据公司公告与公开报道整理；金额为披露值，含首次公开发行、再融资、债券、银行借款与重整投资等）}
\begin{adjustbox}{max width=\textwidth}
\begin{tabular}{@{}p{1.45cm}p{2.55cm}p{4.05cm}p{4.95cm}@{}}
\toprule
时间 & 方式 & 规模 & 用途与说明 \\
\midrule
2010-05 & IPO（深交所创业板） & 发行 1，\allowbreak{}320 万股，\allowbreak{}发行价 35.60 元/\allowbreak{}股，\allowbreak{}实际募集资金 4.70 亿元 & 全部用于与公司主营业务相关的项目及主营业务发展所需的营运资金；占发行后总股本 25\% \\
2015-11 & 发行股份及支付现金购买资产并募集配套资金（未获通过） & 拟购买同路农业 60\% 股权作价 1.44 亿元 & 支付购买资产现金对价及交易费用；交易完成后公司将发行 1，\allowbreak{}049.82 万股 \\
2017-08 & 发行股份及支付现金购买资产并募集配套资金（修订，被否） & 拟购买同路农业 100\% 股权 & 支付购买资产现金对价及本次重组相关费用；2017 年 8 月 25 日临时股东大会 17 项议案全部被否决（中植系、\allowbreak{}大北农系等股东集体反对） \\
2021-07 & 向特定对象发行A股（上市后首次再融资） & 向 6 名特定对象发行 2，\allowbreak{}376.84 万股，\allowbreak{}发行价 23.14 元/\allowbreak{}股 & 研发创新体系建设项目、\allowbreak{}农作物种子海外育繁推一体化建设项目、\allowbreak{}青贮玉米品种产业化及种养结合项目；2021 年 5 月 27 日获证监会同意注册批复，\allowbreak{}2020 年 11 月调减为不超过 5.50 亿元 \\
2023-04-23 & 使用募集资金收购股权 & 第五届董事会第十一次会议通过使用募集资金 2.24 亿元收购河北新纪元种业有限公司股权 & 收购河北新纪元种业股权；2026 年上半年公告显示该笔收购最终交易对价为 1.83 亿元 \\
2025-11-20 & 一致行动关系终止 & 贾桂兰、\allowbreak{}王玉林合计 7，\allowbreak{}435 万股（占总股本 7.85\%）不可撤销承诺预受要约并放弃所持全部股份表决权 & 中种集团表决权比例由 28.36\% 变更为 20.51\%，\allowbreak{}为要约收购前置安排 \\
2026-01 & 控股股东部分要约收购（收购人现金出资，非上市公司融资） & 中种集团以 11.85 元/\allowbreak{}股要约收购 1.89 亿股（占已发行股份 20.00\%） & 提升控股股东持股比例至 40.51\% \\
\bottomrule
\end{tabular}
\end{adjustbox}
\end{center}

\pfl{历史融投资情况} 从现金流量表看，2025 年公司取得借款收到的现金 20.26 亿元、偿还债务支付的现金 19.61 亿元、吸收权益性投资 0.02 亿元，筹资活动现金流净额 -1.16 亿元（2023 年为取得借款 11.99 亿元、偿还债务 8.96 亿元）。 2026 年上半年筹资活动现金流净额为 -0.99 亿元（取得借款 8.82 亿元、偿还债务 9.09 亿元，未年化），已转为净流出，处于净还债状态。 2025 年分配股利、利润或偿付利息支付的现金合计 1.22 亿元。 公开披露的融资记录共 7 笔，工具类型以 IPO、发行股份及支付现金购、定向增发为主；金额与用途见下表。 其中规模最大的两笔为：2023 年 4 月，使用募集资金收购股权，第五届董事会第十一次会议通过使用募集资金 2.24 亿元收购河北新纪元种业有限公司股权；2010 年 5 月，IPO，发行 1，320 万股，发行价 35.60 元/股，实际募集资金 4.70 亿元。 股本方面，近三年披露股本变动 4 次（转增 1 次）；总股本由 2021 年末的 4.54 亿股变为 9.47 亿股（+108.62\%）。 分红方面，近三年现金分红 3 次、合计每股派现 0.25 元，最近一次实施在 2024 年 12 月。 近三年“再融资”类公告 1 条，涉及定向增发。 综合看，公司融资结构上以债务滚动为主、股权融资为补充；2025 年筹资活动现金流净额 -1.16 亿元，融资节奏仍以维持存量债务周转为主，与前述经营与投资现金流的方向基本一致。

\pfl{未来潜在融资需求分析} 2026 年 6 月末货币资金 7.67 亿元，而短期债务 19.06 亿元（现金短债比 0.40 倍），2026 年上半年经营现金流净额 -7.03 亿元。按“短期债务减去账面货币资金”的滚动偿付口径，静态缺口约 11.39 亿元；若再计入上半年亏损的年化额（1.11 亿元），维持一年正常经营所需的外部资金约 11.39---12.50 亿元。 按第 \ref{sec:financing-need} 节的三档划分，公司属于第二档“保壳型”：近三年重大重组与并购类公告 18 条，2026 年上半年归母净利润 -0.55 亿元，融资需求更多体现为“资产置换 + 维持上市地位”，而不是扩产。 资金需求还要叠加在建工程：期末在建工程 0.53 亿元、2026 年上半年购建固定资产等支出 0.73 亿元（未年化），这部分支出在周期底部通常会被主动放缓。 综合看，公司的外部资金需求以补充流动性与项目投入为主，滚动偿付口径下的静态缺口约 11.39 亿元，融资的时点与规模更多取决于信用条件与资本市场窗口。




\subsubsection{万向德农（600371）}
\label{co:600371}

\pfl{公司基本情况} 公司前身可追溯至 1995 年，黑龙江起家、2002 年上市的玉米种子企业（原华冠科技）；主营高度集中于玉米杂交种子，收入自 2025 年起大幅下滑，但资产负债率低、无银行借款、分红与增持持续。 公司于 2002-09-16 在上交所首发上市，注册资本 2.93 亿元，总股本 2.93 亿股。 截至 2026 年 9 月 24 日收盘价 15.85 元，流通市值 46.4 亿元，近一年+76.7\%，流通市值在三级行业“种子”的 8 家公司中按由大到小排第\zhnumber{7}位，近一年涨跌幅高于全样本中位数（-10.1\%）86.8 个百分点。 按 2020 年年报披露的行业构成，收入占比前两位为种子行业 97.4\%；其他(补充) 2.6\%。

\pfl{历史股价} 本报告样本区间（2021 年 1 月 至 2026 年 9 月）内，股价由 12.32 元变为 15.85 元，累计 +28.7\%，区间最高 15.85 元（2026 年 9 月）、最低 6.25 元（2026 年 6 月）；当前价相当于区间最高价的 100.0\%，为区间最低价的 2.54 倍。 月度收益的年化波动率 44.0\%，高于全样本中位数 37.4\%，在三级行业“种子”的 8 家公司中波动率由高到低排第\zhnumber{4}位。

\begin{center}
\begin{minipage}{\textwidth}\centering
\begin{tikzpicture}
\begin{axis}[
  width=12.6cm, height=3.9cm, scale only axis,
  xmin=2020.922, xmax=2026.888, ymin=5.81, ymax=16.96,
  xtick={2021,2022,2023,2024,2025,2026},
  yearticks,
  yticklabel style={font=\scriptsize},
  ylabel={元/股}, ylabel style={font=\scriptsize},
  grid=major, grid style={LightGray, line width=0.3pt},
  axis line style={InkGray, line width=0.5pt},
  every axis plot/.append style={line width=0.9pt},
]
\addplot[AgriGreen] table[x=x,y=y]{data/clean/profiles/px_600371.dat};
\addplot[SoftRed, dashed, line width=0.7pt] table[x=x,y=y]{data/clean/profiles/pxzz_600371.dat};
\addplot[only marks, mark=*, mark size=1.1pt, SoftRed] table[x=x,y=y]{data/clean/profiles/pxzz_600371.dat};
\end{axis}
\end{tikzpicture}

\captionof{figure}[万向德农（600371）股价与周期骨架]{万向德农（600371）月度收盘价（元/股）与 ZigZag 周期骨架（阈值 25\%，圆点为周期转折点）}
\end{minipage}
\end{center}

\pfl{过去周期分析} 以 25\% 的 ZigZag 阈值划分，样本区间内股价形成 4 个主升段与 4 个主跌段。 最大主跌段为 2023 年 1 月 至 2024 年 8 月，19 个月-53.1\%；最大主升段出现在 2026 年 6 月 至 2026 年 9 月，3 个月+153.6\%。 最近一次转折出现在 2026 年 9 月（上涨段自 2026 年 6 月起，3 个月+153.6\%），当前处于该段之后的震荡之中。 公司股价与申万农林牧渔指数几乎同步见底，个股并未走出独立行情。 种子行业的价格由制种成本、粮价与政策（转基因商业化、种业振兴）共同决定（第 \ref{sec:planting} 节），与生猪价格的联动有限，公司 2026 年上半年实现盈利、半年 ROE 2.6\%（未年化），好于全样本中位数（0.75\%），下行周期对其报表的冲击小于同业。 综合区间形态看，该股单段涨跌幅度偏大、以趋势段为主（最大主升 +153.6\%、最大主跌 -53.1\%），年化波动率 44.0\% 与全样本中位数 37.4\% 相比波动明显大于样本中位数。

\pfl{盈利情况} 2026 年上半年营业收入 0.99 亿元（同比 -15.1\%），归母净利润 0.14 亿元，同比 -44.5\%。 以年报为趋势对照，2023---2025 年营业收入由 3.19 亿元变为 1.98 亿元（两年年均复合增速 -21.3\%），归母净利润由 0.65 亿元变为 0.06 亿元。 按半年报口径，2026 年上半年的 ROE 为 2.6\%（未年化）、销售净利率 15.5\%、毛利率 26.5\%，ROE 在“种子”8 家公司中按数值由高到低排第\zhnumber{2}位。 同期全样本半年 ROE 中位数 0.75\%，公司与之相差 1.87 个百分点（略好于中位数公司）。 从公司披露的原因看，业绩变动主要来自：2025 年，营业收入下降主因玉米种子销量、销售价格均下降，营业成本下降主因制种成本下降。 综合看，公司在报告期维持盈利（高于全样本半年 ROE 中位数 0.75\%），利润的可持续性取决于种子景气的延续与成本管控的执行。

\begin{center}\small\setlength{\tabcolsep}{3.4pt}
\captionof{table}[万向德农（600371）关键财务指标]{万向德农（600371）关键财务指标（合并报表口径；两个半年报期均未年化；现金短债比 $=$ 货币资金 $\div$ 短期债务）}
\begin{adjustbox}{max width=\textwidth}
\begin{tabular}{@{}lrrrrr@{}}
\toprule
指标 & 2023 年 & 2024 年 & 2025 年 & 2025 年半年报 & 2026 年半年报 \\
\midrule
营业收入（亿元） & 3.19 & 3.43 & 1.98 & 1.17 & 0.99 \\
归母净利润（亿元） & 0.65 & 0.53 & 0.06 & 0.25 & 0.14 \\
毛利率（\%） & 34.9 & 27.1 & 27.3 & 33.7 & 26.5 \\
ROE（\%） & 11.3 & 9.1 & 1.0 & 4.4 & 2.6 \\
资产负债率（\%） & 32.4 & 22.5 & 21.0 & 15.0 & 14.5 \\
经营活动现金流净额（亿元） & 0.55 & -0.73 & -0.38 & -0.30 & -0.16 \\
货币资金（亿元） & 4.74 & 3.47 & 2.69 & 2.86 & 2.26 \\
有息负债合计（亿元） & --- & --- & --- & --- & 0.00 \\
现金短债比（倍） & --- & --- & --- & --- & --- \\
\bottomrule
\end{tabular}
\end{adjustbox}
\end{center}

\begin{center}
\begin{minipage}{\textwidth}\centering
\begin{tikzpicture}
\begin{groupplot}[
  group style={group size=2 by 1, horizontal sep=0.60cm},
  width=6.85cm, height=4.60cm,
  tick label style={font=\tiny},
  title style={font=\scriptsize\heiti},
  yticklabel style={font=\tiny},
  ylabel style={font=\tiny},
  grid=major, grid style={LightGray, line width=0.3pt},
  axis line style={InkGray, line width=0.5pt},
  enlarge x limits={abs=0.8},
  legend style={font=\scriptsize, at={(0.5,-0.20)}, anchor=north,
                legend columns=2, draw=none, fill=none, column sep=6pt},
]
\nextgroupplot[
  ybar, bar width=4pt,
  ymin=-0.34, ymax=3.84,
  xtick={0,1,2,3,4,5.7},
  xticklabels={2021,2022,2023,2024,2025,2026H1},
  ylabel={亿元},
]
\addplot[fill=AgriGreen!75, draw=AgriGreen, bar shift=-2.6pt] table[x=i, y=rev]{data/clean/profiles/fin_600371.dat};
\addplot[fill=SoftRed!60, draw=SoftRed, bar shift=2.6pt] table[x=i, y=ni]{data/clean/profiles/fin_600371.dat};
\addplot[fill=AgriGreen!30, draw=AgriGreen, pattern=north east lines, bar shift=-2.6pt] table[x=x, y=rev]{data/clean/profiles/finh_600371.dat};
\addplot[fill=SoftRed!20, draw=SoftRed, pattern=north east lines, bar shift=2.6pt] table[x=x, y=ni]{data/clean/profiles/finh_600371.dat};
\legend{营业收入（年/半年）, 归母净利润（年/半年）}
\nextgroupplot[
  ymin=-2.6, ymax=35.6,
  xtick={0,1,2,3,4,5.7},
  xticklabels={2021,2022,2023,2024,2025,2026H1},
  ylabel={\%},
]
\addplot[AgriGold, mark=*, mark size=1.3pt] table[x=i, y=roe]{data/clean/profiles/fin_600371.dat};
\addplot[OilBlue, mark=square*, mark size=1.3pt] table[x=i, y=debt]{data/clean/profiles/fin_600371.dat};
\addplot[only marks, AgriGold, mark=*, mark size=2.0pt] table[x=x, y=roe]{data/clean/profiles/finh_600371.dat};
\addplot[only marks, OilBlue, mark=square*, mark size=2.0pt] table[x=x, y=debt]{data/clean/profiles/finh_600371.dat};
\legend{ROE, 资产负债率}
\end{groupplot}
\end{tikzpicture}
\begin{tikzpicture}
\begin{axis}[
  name=finbot,
  width=13.75cm, height=3.10cm, scale only axis,
  ymin=-0.47, ymax=5.31,
  xtick={0,1,2,3,4,5.7},
  xticklabels={2021,2022,2023,2024,2025,2026H1},
  tick label style={font=\tiny},
  yticklabel style={font=\tiny},
  ylabel={亿元}, ylabel style={font=\tiny},
  ybar, bar width=4pt,
  grid=major, grid style={LightGray, line width=0.3pt},
  axis line style={InkGray, line width=0.5pt},
  enlarge x limits={abs=0.8},
  legend style={font=\scriptsize, at={(0.5,-0.26)}, anchor=north,
                legend columns=2, draw=none, fill=none, column sep=10pt},
]
\addplot[fill=AgriGreen!55, draw=AgriGreen, bar shift=-3.0pt] table[x=i, y=cash]{data/clean/profiles/fin_600371.dat};
\addplot[fill=SoftRed!45, draw=SoftRed, bar shift=3.0pt] table[x=i, y=ibd]{data/clean/profiles/fin_600371.dat};
\addplot[fill=AgriGreen!25, draw=AgriGreen, pattern=north east lines, bar shift=-3.0pt] table[x=x, y=cash]{data/clean/profiles/finh_600371.dat};
\addplot[fill=SoftRed!20, draw=SoftRed, pattern=north east lines, bar shift=3.0pt] table[x=x, y=ibd]{data/clean/profiles/finh_600371.dat};
\legend{货币资金（左轴）, 有息负债（左轴）}
\end{axis}
\begin{axis}[
  at={(finbot.south east)}, anchor=south east,
  width=13.75cm, height=3.10cm, scale only axis,
  axis y line*=right, axis x line*=none, xtick=\empty,
  ymin=-1.16, ymax=1.41,
  tick label style={font=\tiny},
  yticklabel style={font=\tiny}, scaled y ticks=false,
  legend style={font=\scriptsize, at={(0.5,-0.44)}, anchor=north,
                legend columns=1, draw=none, fill=none},
]
\addplot[OilBlue, mark=*, mark size=2.2pt, line width=1.3pt] table[x=i, y=ocf]{data/clean/profiles/fin_600371.dat};
\addplot[only marks, OilBlue, mark=*, mark size=3.2pt] table[x=x, y=ocf]{data/clean/profiles/finh_600371.dat};
\legend{经营活动现金流净额（右轴，亿元，蓝点）}
\end{axis}
\end{tikzpicture}

\captionof{figure}[万向德农（600371）近年主要财务指标]{万向德农（600371）近年主要财务指标（左上：营业收入与归母净利润；右上：ROE 与资产负债率；下：货币资金、有息负债为左轴柱状，经营活动现金流净额为其右轴的蓝点折线。
2021---2025 年为年报口径，2026H1 为 2026 年半年报/半年末，与其前一年报之间留出间隔、以斜纹或空心点区分；半年报数据未年化）}
\end{minipage}
\end{center}

\pfl{财务分析} 2026 年 6 月末资产负债率 14.5\%，较 2025 年末下降 6.5 个百分点（样本中位数 52.2\%，所处三级行业内按由低到高排第\zhnumber{1}位）。 流动比率 4.96、速动比率 2.78，短期流动性无明显缺口。 2026 年上半年经营活动现金流净额 -0.16 亿元，净利润为正但经营现金流为负，利润的现金含量偏低。 综合看，资产负债率低于样本中位数 37.7 个百分点；有息负债中长期债务占比更高，期限结构对短债滚动的压力较小。

\pfl{重大战略转型} 公司的战略动作可归为以下几类：并购与重组（2 项）：2026 年 8 月，公司通过出售闲置房产确认资产处置收益 213.89 万元计入非经常性损益，对当期利润形成一定支撑；2013 年 12 月，万向三农集团有限公司通过上交所大宗交易平台累计出售公司 2，046 万股股份；其他动作（1 项）：2024 年 4 月，第九届董事会第二十一次会议审议通过德农种业放弃万向财务公司增资优先认缴出资权暨关联交易议案。 公告口径上，近三年（2023 年 9 月---2026 年 9 月）公司披露重大重组与并购类公告 3 条、对外投资与转型类公告 3 条、回购与股权激励类公告 0 条，合计公告 235 条。 综合判断，报告期内公司有 6 条与战略相关的公告落地（重大重组与并购 3 条、对外投资与转型 3 条），资本运作与主业调整之间存在直接关联。

\begin{center}\small\setlength{\tabcolsep}{4pt}
\captionof{table}[万向德农（600371）历史融资与资本运作]{万向德农（600371）历史融资与资本运作（据公司公告与公开报道整理；金额为披露值，含首次公开发行、再融资、债券、银行借款与重整投资等）}
\begin{adjustbox}{max width=\textwidth}
\begin{tabular}{@{}p{1.45cm}p{2.55cm}p{4.05cm}p{4.95cm}@{}}
\toprule
时间 & 方式 & 规模 & 用途与说明 \\
\midrule
2002-08-29 & IPO（上交所主板） & 发行人民币普通股 4，\allowbreak{}000 万股，\allowbreak{}发行价 3.47 元/\allowbreak{}股 & 公司首发募集资金用于主业项目（招股文件口径）；发行后总股本 10，\allowbreak{}000 万股；公司设立时总股本 6 \\
2024-04-12 & 现金分红（2023年度） & 以权益分派股权登记日总股本为基数每 10 股派现金股利 2.00 元 & 股东回报；最近三个会计年度累计现金分红金额（含税）1.61 亿元 \\
2025-2026 & 控股股东增持（二级市场集中竞价） & 万向三农集团拟增持总金额不低于人民币 2，\allowbreak{}500.00 万元、\allowbreak{}不超过人民币 5 & 控股股东增持；公司未发行可转换公司债券 \\
\bottomrule
\end{tabular}
\end{adjustbox}
\end{center}

\pfl{历史融投资情况} 从现金流量表看，2025 年公司，筹资活动现金流净额 -0.49 亿元。 2026 年上半年筹资活动现金流净额为 -0.29 亿元（取得借款 — 亿元、偿还债务 — 亿元，未年化），已转为净流出，处于净还债状态。 2025 年分配股利、利润或偿付利息支付的现金合计 0.49 亿元。 公开披露的融资记录共 3 笔，工具类型以 IPO 为主；金额与用途见下表。 其中规模最大的两笔为：2025---2026 年，控股股东增持，万向三农集团拟增持总金额不低于人民币 2，500.00 万元、不超过人民币 5；2002 年 8 月，IPO，发行人民币普通股 4，000 万股，发行价 3.47 元/股。 股本方面，近三年披露股本变动 3 次；总股本自 2021 年末以来未发生变化（2.93 亿股）。 分红方面，近三年现金分红 4 次、合计每股派现 0.65 元，最近一次实施在 2025 年 12 月。 综合看，公司融资结构上股权与债务融资并举；2025 年筹资活动现金流净额 -0.49 亿元，融资节奏仍以维持存量债务周转为主，与前述经营与投资现金流的方向基本一致。

\pfl{未来潜在融资需求分析} 公司 2026 年 6 月末资产负债率 14.5\%、2026 年上半年 ROE 2.6\%（未年化），资产端与负债端都没有明显压力，融资需求以相机抉择的扩张型用途为主。




\subsubsection{众兴菌业（002772）}
\label{co:002772}

\pfl{公司基本情况} 甘肃天水的工厂化食用菌龙头，2015 年深交所上市，以金针菇、双孢菇为核心品种（双孢菇满产后日产能达 420 吨/天）；2025 年下半年起行业企稳回暖，价格修复带来利润大幅弹性，同时以技改方式布局冬虫夏草工厂化生态繁育，推进多品种协同与西南总部基地建设。 公司于 2015-06-26 在深交所主板首发上市，注册资本 3.75 亿元，总股本 3.93 亿股。 截至 2026 年 9 月 24 日收盘价 10.64 元，流通市值 39.5 亿元，近一年+13.2\%，流通市值在三级行业“食用菌”的 4 家公司中按由大到小排第\zhnumber{1}位，近一年涨跌幅高于全样本中位数（-10.1\%）23.3 个百分点。 按 2026 年半年报披露的行业构成，收入占比前两位为农业种植业 99.6\%；其他(补充) 0.4\%。

\pfl{历史股价} 自 2021 年 1 月至 2026 年 9 月，股价由 8.46 元变为 10.64 元，累计 +25.8\%，区间最高 18.37 元（2026 年 2 月）、最低 5.88 元（2024 年 7 月）；当前价相当于区间最高价的 57.9\%，为区间最低价的 1.81 倍。 月度收益的年化波动率 50.1\%，高于全样本中位数 37.4\%，在三级行业“食用菌”的 4 家公司中波动率由高到低排第\zhnumber{2}位。

\begin{center}
\begin{minipage}{\textwidth}\centering
\begin{tikzpicture}
\begin{axis}[
  width=12.6cm, height=3.9cm, scale only axis,
  xmin=2020.922, xmax=2026.888, ymin=5.47, ymax=19.66,
  xtick={2021,2022,2023,2024,2025,2026},
  yearticks,
  yticklabel style={font=\scriptsize},
  ylabel={元/股}, ylabel style={font=\scriptsize},
  grid=major, grid style={LightGray, line width=0.3pt},
  axis line style={InkGray, line width=0.5pt},
  every axis plot/.append style={line width=0.9pt},
]
\addplot[AgriGreen] table[x=x,y=y]{data/clean/profiles/px_002772.dat};
\addplot[SoftRed, dashed, line width=0.7pt] table[x=x,y=y]{data/clean/profiles/pxzz_002772.dat};
\addplot[only marks, mark=*, mark size=1.1pt, SoftRed] table[x=x,y=y]{data/clean/profiles/pxzz_002772.dat};
\end{axis}
\end{tikzpicture}

\captionof{figure}[众兴菌业（002772）股价与周期骨架]{众兴菌业（002772）月度收盘价（元/股）与 ZigZag 周期骨架（阈值 25\%，圆点为周期转折点）}
\end{minipage}
\end{center}

\pfl{过去周期分析} 以 25\% 的 ZigZag 阈值划分，样本区间内股价形成 3 个主升段与 3 个主跌段。 最大主跌段为 2021 年 6 月 至 2022 年 4 月，10 个月-59.2\%；最大主升段出现在 2024 年 7 月 至 2026 年 2 月，19 个月+212.4\%。 最近一次转折出现在 2026 年 7 月（下跌段自 2026 年 2 月起，5 个月-44.3\%），当前处于该段的延续之中。 公司股价较申万农林牧渔指数提前约 22 个月见底，是板块中较早定价周期反转的公司之一。 食用菌是工厂化生产的制造型农业（第 \ref{sec:soft} 章），产能投放节奏与金针菇等单品价格是盈利的主要变量，与养殖周期基本无关，公司 2026 年上半年实现盈利、半年 ROE 5.3\%（未年化），好于全样本中位数（0.75\%），下行周期对其报表的冲击小于同业。 综合区间形态看，该股单段涨跌幅度偏大、以趋势段为主（最大主升 +212.4\%、最大主跌 -59.2\%），年化波动率 50.1\% 与全样本中位数 37.4\% 相比波动明显大于样本中位数。

\pfl{盈利情况} 2026 年上半年营业收入 10.85 亿元（同比 +19.6\%），归母净利润 1.91 亿元，同比 +176.5\%。 以年报为趋势对照，2023---2025 年营业收入由 19.31 亿元变为 20.98 亿元（两年年均复合增速 +4.2\%），归母净利润由 1.60 亿元变为 3.35 亿元。 按半年报口径，2026 年上半年的 ROE 为 5.3\%（未年化）、销售净利率 17.6\%、毛利率 30.9\%，ROE 在“食用菌”4 家公司中按数值由高到低排第\zhnumber{1}位。 同期全样本半年 ROE 中位数 0.75\%，公司与之相差 4.59 个百分点（略好于中位数公司）。 综合看，公司在报告期维持盈利（高于全样本半年 ROE 中位数 0.75\%），利润的可持续性取决于食用菌景气的延续与成本管控的执行。

\begin{center}\small\setlength{\tabcolsep}{3.4pt}
\captionof{table}[众兴菌业（002772）关键财务指标]{众兴菌业（002772）关键财务指标（合并报表口径；两个半年报期均未年化；现金短债比 $=$ 货币资金 $\div$ 短期债务）}
\begin{adjustbox}{max width=\textwidth}
\begin{tabular}{@{}lrrrrr@{}}
\toprule
指标 & 2023 年 & 2024 年 & 2025 年 & 2025 年半年报 & 2026 年半年报 \\
\midrule
营业收入（亿元） & 19.31 & 19.35 & 20.98 & 9.08 & 10.85 \\
归母净利润（亿元） & 1.60 & 1.28 & 3.35 & 0.69 & 1.91 \\
毛利率（\%） & 23.0 & 20.8 & 28.6 & 19.4 & 30.9 \\
ROE（\%） & 4.8 & 3.8 & 9.9 & 2.1 & 5.3 \\
资产负债率（\%） & 43.4 & 46.4 & 50.1 & 49.7 & 54.0 \\
经营活动现金流净额（亿元） & 5.02 & 4.95 & 4.96 & 1.50 & 1.72 \\
货币资金（亿元） & 10.39 & 11.91 & 14.50 & 12.86 & 16.69 \\
有息负债合计（亿元） & 20.84 & 23.70 & 28.90 & 26.90 & 35.65 \\
现金短债比（倍） & 1.03 & 0.91 & 1.57 & 1.17 & 1.35 \\
\bottomrule
\end{tabular}
\end{adjustbox}
\end{center}

\begin{center}
\begin{minipage}{\textwidth}\centering
\begin{tikzpicture}
\begin{groupplot}[
  group style={group size=2 by 1, horizontal sep=0.60cm},
  width=6.85cm, height=4.60cm,
  tick label style={font=\tiny},
  title style={font=\scriptsize\heiti},
  yticklabel style={font=\tiny},
  ylabel style={font=\tiny},
  grid=major, grid style={LightGray, line width=0.3pt},
  axis line style={InkGray, line width=0.5pt},
  enlarge x limits={abs=0.8},
  legend style={font=\scriptsize, at={(0.5,-0.20)}, anchor=north,
                legend columns=2, draw=none, fill=none, column sep=6pt},
]
\nextgroupplot[
  ybar, bar width=4pt,
  ymin=-2.10, ymax=23.50,
  xtick={0,1,2,3,4,5.7},
  xticklabels={2021,2022,2023,2024,2025,2026H1},
  ylabel={亿元},
]
\addplot[fill=AgriGreen!75, draw=AgriGreen, bar shift=-2.6pt] table[x=i, y=rev]{data/clean/profiles/fin_002772.dat};
\addplot[fill=SoftRed!60, draw=SoftRed, bar shift=2.6pt] table[x=i, y=ni]{data/clean/profiles/fin_002772.dat};
\addplot[fill=AgriGreen!30, draw=AgriGreen, pattern=north east lines, bar shift=-2.6pt] table[x=x, y=rev]{data/clean/profiles/finh_002772.dat};
\addplot[fill=SoftRed!20, draw=SoftRed, pattern=north east lines, bar shift=2.6pt] table[x=x, y=ni]{data/clean/profiles/finh_002772.dat};
\legend{营业收入（年/半年）, 归母净利润（年/半年）}
\nextgroupplot[
  ymin=-4.3, ymax=59.4,
  xtick={0,1,2,3,4,5.7},
  xticklabels={2021,2022,2023,2024,2025,2026H1},
  ylabel={\%},
]
\addplot[AgriGold, mark=*, mark size=1.3pt] table[x=i, y=roe]{data/clean/profiles/fin_002772.dat};
\addplot[OilBlue, mark=square*, mark size=1.3pt] table[x=i, y=debt]{data/clean/profiles/fin_002772.dat};
\addplot[only marks, AgriGold, mark=*, mark size=2.0pt] table[x=x, y=roe]{data/clean/profiles/finh_002772.dat};
\addplot[only marks, OilBlue, mark=square*, mark size=2.0pt] table[x=x, y=debt]{data/clean/profiles/finh_002772.dat};
\legend{ROE, 资产负债率}
\end{groupplot}
\end{tikzpicture}
\begin{tikzpicture}
\begin{axis}[
  name=finbot,
  width=13.75cm, height=3.10cm, scale only axis,
  ymin=-3.57, ymax=39.93,
  xtick={0,1,2,3,4,5.7},
  xticklabels={2021,2022,2023,2024,2025,2026H1},
  tick label style={font=\tiny},
  yticklabel style={font=\tiny},
  ylabel={亿元}, ylabel style={font=\tiny},
  ybar, bar width=4pt,
  grid=major, grid style={LightGray, line width=0.3pt},
  axis line style={InkGray, line width=0.5pt},
  enlarge x limits={abs=0.8},
  legend style={font=\scriptsize, at={(0.5,-0.26)}, anchor=north,
                legend columns=2, draw=none, fill=none, column sep=10pt},
]
\addplot[fill=AgriGreen!55, draw=AgriGreen, bar shift=-3.0pt] table[x=i, y=cash]{data/clean/profiles/fin_002772.dat};
\addplot[fill=SoftRed!45, draw=SoftRed, bar shift=3.0pt] table[x=i, y=ibd]{data/clean/profiles/fin_002772.dat};
\addplot[fill=AgriGreen!25, draw=AgriGreen, pattern=north east lines, bar shift=-3.0pt] table[x=x, y=cash]{data/clean/profiles/finh_002772.dat};
\addplot[fill=SoftRed!20, draw=SoftRed, pattern=north east lines, bar shift=3.0pt] table[x=x, y=ibd]{data/clean/profiles/finh_002772.dat};
\legend{货币资金（左轴）, 有息负债（左轴）}
\end{axis}
\begin{axis}[
  at={(finbot.south east)}, anchor=south east,
  width=13.75cm, height=3.10cm, scale only axis,
  axis y line*=right, axis x line*=none, xtick=\empty,
  ymin=-6.52, ymax=7.77,
  tick label style={font=\tiny},
  yticklabel style={font=\tiny}, scaled y ticks=false,
  legend style={font=\scriptsize, at={(0.5,-0.44)}, anchor=north,
                legend columns=1, draw=none, fill=none},
]
\addplot[OilBlue, mark=*, mark size=2.2pt, line width=1.3pt] table[x=i, y=ocf]{data/clean/profiles/fin_002772.dat};
\addplot[only marks, OilBlue, mark=*, mark size=3.2pt] table[x=x, y=ocf]{data/clean/profiles/finh_002772.dat};
\legend{经营活动现金流净额（右轴，亿元，蓝点）}
\end{axis}
\end{tikzpicture}

\captionof{figure}[众兴菌业（002772）近年主要财务指标]{众兴菌业（002772）近年主要财务指标（左上：营业收入与归母净利润；右上：ROE 与资产负债率；下：货币资金、有息负债为左轴柱状，经营活动现金流净额为其右轴的蓝点折线。
2021---2025 年为年报口径，2026H1 为 2026 年半年报/半年末，与其前一年报之间留出间隔、以斜纹或空心点区分；半年报数据未年化）}
\end{minipage}
\end{center}

\pfl{财务分析} 2026 年 6 月末资产负债率 54.0\%，较 2025 年末上升 3.8 个百分点（样本中位数 52.2\%，所处三级行业内按由低到高排第\zhnumber{3}位）。 2026 年 6 月末货币资金 16.69 亿元、短期债务（短期借款与一年内到期非流动负债）12.38 亿元、长期债务 23.27 亿元，有息负债合计 35.65 亿元，现金短债比 1.35 倍——覆盖充分。 净有息负债 18.96 亿元，相当于其 2026 年上半年营业收入的 175\%。 流动比率 2.24、速动比率 1.24，短期指标尚可。 2026 年上半年经营活动现金流净额 1.72 亿元，经营现金流与利润同向为正，内生资金可以支撑运营。 2026 年 6 月末在建工程 5.95 亿元，较上年末增加 1.91 亿元，仍有在建投入。 综合看，现金短债比 1.35 倍，短期偿付压力可控；资产负债率高于样本中位数 1.8 个百分点；有息负债中长期债务占比更高，期限结构对短债滚动的压力较小。 公司披露的股东与质押风险为：2026 年 8 月，控股股东陶军持股比例 29.39\%（1.10 亿股），其中质押 2，889.00 万股。

\pfl{重大战略转型} 从公开披露看，公司的战略推进集中在：产能与项目（4 项）：2025 年 6 月，第五届董事会第十二次会议通过拟签订项目投资协议暨对外投资；2025 年 2 月，公司投资实施冬虫夏草工厂化仿生培育更新改造项目，预计新增投资 1.23 亿元（其中固定资产投资 1.04 亿元）；融资与资本运作（1 项）：2020 年 4 月，非公开发行募投项目年产 7，500 吨蟹味菇生产线、年产 7。 公告口径上，近三年（2023 年 9 月---2026 年 9 月）公司披露重大重组与并购类公告 0 条、对外投资与转型类公告 14 条、回购与股权激励类公告 23 条，合计公告 440 条。 第一大行业仍是“农业种植业”，收入占比 99.6\%，与 2021 年年报相比没有方向性变化。 综合判断，报告期内公司有 14 条与战略相关的公告落地（重大重组与并购 0 条、对外投资与转型 14 条），资本运作与主业调整之间存在直接关联。

\begin{center}\small\setlength{\tabcolsep}{4pt}
\captionof{table}[众兴菌业（002772）历史融资与资本运作]{众兴菌业（002772）历史融资与资本运作（据公司公告与公开报道整理；金额为披露值，含首次公开发行、再融资、债券、银行借款与重整投资等）}
\begin{adjustbox}{max width=\textwidth}
\begin{tabular}{@{}p{1.45cm}p{2.55cm}p{4.05cm}p{4.95cm}@{}}
\toprule
时间 & 方式 & 规模 & 用途与说明 \\
\midrule
2015-06 & IPO（深交所） & 发行 3，\allowbreak{}725 万股，\allowbreak{}发行价 13.00 元/\allowbreak{}股，\allowbreak{}募集资金总额 4.84 亿元 & 年产 12，\allowbreak{}600 吨金针菇生产线建设项目、\allowbreak{}年产 7 \\
2016-03-18 & 限制性股票（股权激励） & 筹资净额 9,\allowbreak{}361.94 万元 & 股权激励认购；公司历次筹资记录中的一项 \\
2016-08-22 & 非公开发行（向特定投资者） & 发行 5，\allowbreak{}368.33 万股，\allowbreak{}每股 21.00 元，\allowbreak{}募集资金总额 11.27 亿元 & 非公开发行募投项目（后部分变更）；截至 2020 年 12 月 31 日，\allowbreak{}非公开发行募集资金本息合计 11.76 亿元 \\
2017-12-13 & 可转换公司债券（众兴转债，128026） & 发行 920 万张、\allowbreak{}每张面值 100.00 元，\allowbreak{}募集资金 9.20 亿元 & 年产 2 万吨双孢蘑菇及 11 万吨堆肥工厂化生产项目、\allowbreak{}年产 32，\allowbreak{}400 吨金针菇生产线建设项目、\allowbreak{}年产 32 \\
2025-02-07 & 股份回购（回购专用证券账户集中竞价） & 累计回购 1，\allowbreak{}849.88 万股，\allowbreak{}占公司总股本 4.7042\% & 股份回购（公司曾披露取得金融机构回购专项贷款承诺函）；回购专户股份不参与 2025 年度利润分配基数 \\
\bottomrule
\end{tabular}
\end{adjustbox}
\end{center}

\pfl{历史融投资情况} 从现金流量表看，2025 年公司取得借款收到的现金 18.26 亿元、偿还债务支付的现金 13.54 亿元、吸收权益性投资 0.07 亿元，筹资活动现金流净额 2.67 亿元（2023 年为取得借款 8.15 亿元、偿还债务 13.59 亿元）。 2026 年上半年筹资活动现金流净额为 3.34 亿元（取得借款 10.96 亿元、偿还债务 5.37 亿元，未年化），仍处于净融入状态。 2025 年分配股利、利润或偿付利息支付的现金合计 1.87 亿元。 公开披露的融资记录共 5 笔，工具类型以 IPO、定向增发、可转债为主；金额与用途见下表。 其中规模最大的两笔为：2017 年 12 月，可转债，发行 920 万张、每张面值 100.00 元，募集资金 9.20 亿元；2016 年 8 月，定向增发，发行 5，368.33 万股，每股 21.00 元，募集资金总额 11.27 亿元。 股本方面，近三年披露股本变动 7 次（可转债转股 5 次、股份回购 1 次）；总股本由 2021 年末的 4.07 亿股变为 3.93 亿股（-3.48\%）。 分红方面，近三年现金分红 6 次、合计每股派现 1.50 元，最近一次实施在 2026 年 6 月。 近三年“再融资”类公告 1 条，涉及可转债。 综合看，公司融资结构上以债务滚动为主、股权融资为补充；2025 年筹资活动现金流净额 2.67 亿元，年度内仍为净融入，与前述经营与投资现金流的方向基本一致。

\pfl{未来潜在融资需求分析} 2026 年 6 月末货币资金 16.69 亿元，而短期债务 12.38 亿元（现金短债比 1.35 倍），2026 年上半年经营现金流净额 1.72 亿元。账面货币资金可以覆盖短期债务，缺口不体现在静态偿付层面。 按第 \ref{sec:financing-need} 节的三档划分，公司属于第三档“扩张型”：2026 年上半年 ROE 5.3\%（未年化，折合约 10.7\%）、6 月末资产负债率 54.0\%，资产负债表仍具备加杠杆空间；融资用途以产能扩张、品类并购与海外布局为主，单笔规模通常在 5---20 亿元量级，会被市场定价为成长而非补血。 公开披露的后续资金安排线索包括：债务与授信安排：2025 年 2 月，股份回购方案资金安排曾取得金融机构回购专项贷款承诺函；2026 年度，公司为合并报表范围内全资子公司及全资孙公司提供担保的预计额度为 11.00 亿元，用于子公司日常经营及项目建设。 资金需求还要叠加在建工程：期末在建工程 5.95 亿元、2026 年上半年购建固定资产等支出 3.02 亿元（未年化），这部分支出在周期底部通常会被主动放缓。 综合看，公司在静态偿付层面不存在缺口，外部融资更多是配合扩张节奏与并购整合的主动性安排，而非偿债压力驱动。




\subsubsection{国投丰乐（000713）}
\label{co:000713}

\pfl{公司基本情况} 中国种业首家上市公司（1997 年深交所上市，原合肥丰乐种业），主营种子、农化、香料三大板块；2024 年国投种业受让 20\% 股份成为控股股东、实际控制人变更为国务院国资委，2025 年更名为国投丰乐，2026 年 1 月完成向控股股东定增 10.89 亿元，但 2026 年上半年仍亏损。 公司于 1997-04-22 在深交所主板首发上市，注册资本 7.98 亿元，总股本 6.14 亿股。 截至 2026 年 9 月 24 日收盘价 6.35 元，流通市值 39.0 亿元，近一年-8.0\%，流通市值在三级行业“种子”的 8 家公司中按由大到小排第\zhnumber{8}位，近一年涨跌幅高于全样本中位数（-10.1\%）2.1 个百分点。 按 2026 年半年报披露的行业构成，收入占比前三位为农化类 69.6\%；种子类 16.5\%；香料类 13.8\%。

\pfl{历史股价} 自 2021 年 1 月至 2026 年 9 月，股价由 15.87 元变为 6.35 元，累计 -60.0\%，区间最高 17.08 元（2021 年 2 月）、最低 5.22 元（2026 年 6 月）；当前价相当于区间最高价的 37.2\%，为区间最低价的 1.22 倍。 月度收益的年化波动率 33.6\%，低于全样本中位数 37.4\%，在三级行业“种子”的 8 家公司中波动率由高到低排第\zhnumber{6}位。 把股价阶段与公司事件对齐看，2022 年 5 月 至 2024 年 6 月股价下跌-50.1\%，同期（2024 年 4 月）完成股份过户登记取得丰乐种业 20\% 股权。 从时间对应关系看，公司层面的资本运作与股权、风险事件解释了区间内多数急涨急跌，与行业指数的同步性相对有限。

\begin{center}
\begin{minipage}{\textwidth}\centering
\begin{tikzpicture}
\begin{axis}[
  width=12.6cm, height=3.9cm, scale only axis,
  xmin=2020.922, xmax=2026.888, ymin=4.85, ymax=18.28,
  xtick={2021,2022,2023,2024,2025,2026},
  yearticks,
  yticklabel style={font=\scriptsize},
  ylabel={元/股}, ylabel style={font=\scriptsize},
  grid=major, grid style={LightGray, line width=0.3pt},
  axis line style={InkGray, line width=0.5pt},
  every axis plot/.append style={line width=0.9pt},
]
\addplot[AgriGreen] table[x=x,y=y]{data/clean/profiles/px_000713.dat};
\addplot[SoftRed, dashed, line width=0.7pt] table[x=x,y=y]{data/clean/profiles/pxzz_000713.dat};
\addplot[only marks, mark=*, mark size=1.1pt, SoftRed] table[x=x,y=y]{data/clean/profiles/pxzz_000713.dat};
\end{axis}
\end{tikzpicture}

\captionof{figure}[国投丰乐（000713）股价与周期骨架]{国投丰乐（000713）月度收盘价（元/股）与 ZigZag 周期骨架（阈值 25\%，圆点为周期转折点）}
\end{minipage}
\end{center}

\pfl{过去周期分析} 以 25\% 的 ZigZag 阈值划分，样本区间内股价形成 2 个主升段与 3 个主跌段。 最大主跌段为 2022 年 5 月 至 2024 年 6 月，25 个月-50.1\%；最大主升段出现在 2024 年 6 月 至 2024 年 11 月，5 个月+37.6\%。 最近一次转折出现在 2026 年 6 月（下跌段自 2024 年 11 月起，19 个月-30.1\%），当前处于该段的延续之中。 公司股价与申万农林牧渔指数几乎同步见底，个股并未走出独立行情。 种子行业的价格由制种成本、粮价与政策（转基因商业化、种业振兴）共同决定（第 \ref{sec:planting} 节），与生猪价格的联动有限，公司 2026 年 6 月末资产负债率 32.1\%、半年 ROE -1.0\%（未年化），在本轮下行中属于随行业同步调整的一类。 综合区间形态看，该股单段涨跌幅度偏大、以趋势段为主（最大主升 +37.6\%、最大主跌 -50.1\%），年化波动率 33.6\% 与全样本中位数 37.4\% 相比波动小于样本中位数。

\pfl{盈利情况} 2026 年上半年营业收入 11.95 亿元（同比 +3.9\%），归母净利润 -0.31 亿元，亏损同比扩大 0.03 亿元。 以年报为趋势对照，2023---2025 年营业收入由 31.14 亿元变为 29.03 亿元（两年年均复合增速 -3.4\%），归母净利润由 0.40 亿元变为 0.66 亿元。 按半年报口径，2026 年上半年的 ROE 为 -1.0\%（未年化）、销售净利率 -3.0\%、毛利率 12.0\%，ROE 在“种子”8 家公司中按数值由高到低排第\zhnumber{5}位。 同期全样本半年 ROE 中位数 0.75\%，公司与之相差 1.75 个百分点（弱于中位数公司）。 从公司披露的原因看，业绩变动主要来自：2026 年上半年，同比减少 561.36 万元、下降 605.29\%。 综合看，公司仍处亏损区间、亏损同比扩大，半年 ROE 在三级行业内按由高到低排第\zhnumber{5}位，单位盈利能力的修复依赖行业景气与自身成本端的改善，报表弹性主要体现为价格与销量的方向性变化。

\begin{center}\small\setlength{\tabcolsep}{3.4pt}
\captionof{table}[国投丰乐（000713）关键财务指标]{国投丰乐（000713）关键财务指标（合并报表口径；两个半年报期均未年化；现金短债比 $=$ 货币资金 $\div$ 短期债务）}
\begin{adjustbox}{max width=\textwidth}
\begin{tabular}{@{}lrrrrr@{}}
\toprule
指标 & 2023 年 & 2024 年 & 2025 年 & 2025 年半年报 & 2026 年半年报 \\
\midrule
营业收入（亿元） & 31.14 & 29.26 & 29.03 & 11.50 & 11.95 \\
归母净利润（亿元） & 0.40 & 0.70 & 0.66 & -0.28 & -0.31 \\
毛利率（\%） & 12.7 & 15.7 & 16.5 & 10.7 & 12.0 \\
ROE（\%） & 2.1 & 3.5 & 3.2 & -1.4 & -1.0 \\
资产负债率（\%） & 30.6 & 35.8 & 29.7 & 37.1 & 32.1 \\
经营活动现金流净额（亿元） & 1.66 & -0.28 & 2.92 & -1.05 & -0.97 \\
货币资金（亿元） & 1.45 & 2.99 & 16.71 & 2.22 & 15.50 \\
有息负债合计（亿元） & 1.57 & 2.37 & 3.03 & 3.68 & 3.72 \\
现金短债比（倍） & 26.60 & 4.64 & 9.73 & 1.07 & 6.52 \\
\bottomrule
\end{tabular}
\end{adjustbox}
\end{center}

\begin{center}
\begin{minipage}{\textwidth}\centering
\begin{tikzpicture}
\begin{groupplot}[
  group style={group size=2 by 1, horizontal sep=0.60cm},
  width=6.85cm, height=4.60cm,
  tick label style={font=\tiny},
  title style={font=\scriptsize\heiti},
  yticklabel style={font=\tiny},
  ylabel style={font=\tiny},
  grid=major, grid style={LightGray, line width=0.3pt},
  axis line style={InkGray, line width=0.5pt},
  enlarge x limits={abs=0.8},
  legend style={font=\scriptsize, at={(0.5,-0.20)}, anchor=north,
                legend columns=2, draw=none, fill=none, column sep=6pt},
]
\nextgroupplot[
  ybar, bar width=4pt,
  ymin=-3.46, ymax=34.91,
  xtick={0,1,2,3,4,5.7},
  xticklabels={2021,2022,2023,2024,2025,2026H1},
  ylabel={亿元},
]
\addplot[fill=AgriGreen!75, draw=AgriGreen, bar shift=-2.6pt] table[x=i, y=rev]{data/clean/profiles/fin_000713.dat};
\addplot[fill=SoftRed!60, draw=SoftRed, bar shift=2.6pt] table[x=i, y=ni]{data/clean/profiles/fin_000713.dat};
\addplot[fill=AgriGreen!30, draw=AgriGreen, pattern=north east lines, bar shift=-2.6pt] table[x=x, y=rev]{data/clean/profiles/finh_000713.dat};
\addplot[fill=SoftRed!20, draw=SoftRed, pattern=north east lines, bar shift=2.6pt] table[x=x, y=ni]{data/clean/profiles/finh_000713.dat};
\legend{营业收入（年/半年）, 归母净利润（年/半年）}
\nextgroupplot[
  ymin=-3.9, ymax=39.5,
  xtick={0,1,2,3,4,5.7},
  xticklabels={2021,2022,2023,2024,2025,2026H1},
  ylabel={\%},
]
\addplot[AgriGold, mark=*, mark size=1.3pt] table[x=i, y=roe]{data/clean/profiles/fin_000713.dat};
\addplot[OilBlue, mark=square*, mark size=1.3pt] table[x=i, y=debt]{data/clean/profiles/fin_000713.dat};
\addplot[only marks, AgriGold, mark=*, mark size=2.0pt] table[x=x, y=roe]{data/clean/profiles/finh_000713.dat};
\addplot[only marks, OilBlue, mark=square*, mark size=2.0pt] table[x=x, y=debt]{data/clean/profiles/finh_000713.dat};
\legend{ROE, 资产负债率}
\end{groupplot}
\end{tikzpicture}
\begin{tikzpicture}
\begin{axis}[
  name=finbot,
  width=13.75cm, height=3.10cm, scale only axis,
  ymin=-1.67, ymax=18.72,
  xtick={0,1,2,3,4,5.7},
  xticklabels={2021,2022,2023,2024,2025,2026H1},
  tick label style={font=\tiny},
  yticklabel style={font=\tiny},
  ylabel={亿元}, ylabel style={font=\tiny},
  ybar, bar width=4pt,
  grid=major, grid style={LightGray, line width=0.3pt},
  axis line style={InkGray, line width=0.5pt},
  enlarge x limits={abs=0.8},
  legend style={font=\scriptsize, at={(0.5,-0.26)}, anchor=north,
                legend columns=2, draw=none, fill=none, column sep=10pt},
]
\addplot[fill=AgriGreen!55, draw=AgriGreen, bar shift=-3.0pt] table[x=i, y=cash]{data/clean/profiles/fin_000713.dat};
\addplot[fill=SoftRed!45, draw=SoftRed, bar shift=3.0pt] table[x=i, y=ibd]{data/clean/profiles/fin_000713.dat};
\addplot[fill=AgriGreen!25, draw=AgriGreen, pattern=north east lines, bar shift=-3.0pt] table[x=x, y=cash]{data/clean/profiles/finh_000713.dat};
\addplot[fill=SoftRed!20, draw=SoftRed, pattern=north east lines, bar shift=3.0pt] table[x=x, y=ibd]{data/clean/profiles/finh_000713.dat};
\legend{货币资金（左轴）, 有息负债（左轴）}
\end{axis}
\begin{axis}[
  at={(finbot.south east)}, anchor=south east,
  width=13.75cm, height=3.10cm, scale only axis,
  axis y line*=right, axis x line*=none, xtick=\empty,
  ymin=-2.25, ymax=5.44,
  tick label style={font=\tiny},
  yticklabel style={font=\tiny}, scaled y ticks=false,
  legend style={font=\scriptsize, at={(0.5,-0.44)}, anchor=north,
                legend columns=1, draw=none, fill=none},
]
\addplot[OilBlue, mark=*, mark size=2.2pt, line width=1.3pt] table[x=i, y=ocf]{data/clean/profiles/fin_000713.dat};
\addplot[only marks, OilBlue, mark=*, mark size=3.2pt] table[x=x, y=ocf]{data/clean/profiles/finh_000713.dat};
\legend{经营活动现金流净额（右轴，亿元，蓝点）}
\end{axis}
\end{tikzpicture}

\captionof{figure}[国投丰乐（000713）近年主要财务指标]{国投丰乐（000713）近年主要财务指标（左上：营业收入与归母净利润；右上：ROE 与资产负债率；下：货币资金、有息负债为左轴柱状，经营活动现金流净额为其右轴的蓝点折线。
2021---2025 年为年报口径，2026H1 为 2026 年半年报/半年末，与其前一年报之间留出间隔、以斜纹或空心点区分；半年报数据未年化）}
\end{minipage}
\end{center}

\pfl{财务分析} 2026 年 6 月末资产负债率 32.1\%，较 2025 年末上升 2.4 个百分点（样本中位数 52.2\%，所处三级行业内按由低到高排第\zhnumber{3}位）。 2026 年 6 月末货币资金 15.50 亿元、短期债务（短期借款与一年内到期非流动负债）2.38 亿元、长期债务 1.34 亿元，有息负债合计 3.72 亿元，现金短债比 6.52 倍——覆盖充分。 净有息负债 -11.78 亿元，相当于其 2026 年上半年营业收入的 -99\%。 流动比率 2.45、速动比率 1.77，短期偿债能力处于正常区间。 2026 年上半年经营活动现金流净额 -0.97 亿元，经营现金流与利润同时为负，现金消耗较快。 2026 年 6 月末在建工程 1.95 亿元，较上年末增加 0.51 亿元，仍有在建投入。 2026 年 6 月末商誉 2.21 亿元，占归母净资产的 7.0\%，是净资产中最脆弱的一块。 综合看，现金短债比 6.52 倍，短期偿付压力可控；资产负债率低于样本中位数 20.1 个百分点；有息负债中短期债务占比更高，期限结构仍以滚动续作为主。

\pfl{重大战略转型} 公司的战略动作可归为以下几类：并购与重组（2 项）：2026 年 6 月，与国投种业签署国投（张掖）金种科技有限公司股权转让协议，以 3，376.48 万元现金收购国投金种 60\% 股权；2024 年 4 月，完成股份过户登记取得丰乐种业 20\% 股权；其他动作（2 项）：2026 年上半年，玉米产业销量稳步增长、产品结构不断优化。 公告口径上，近三年（2023 年 9 月---2026 年 9 月）公司披露重大重组与并购类公告 7 条、对外投资与转型类公告 0 条、回购与股权激励类公告 0 条，合计公告 404 条。 第一大行业仍是“农化类”，收入占比 69.6\%，与 2021 年年报相比没有方向性变化。 综合判断，报告期内公司有 7 条与战略相关的公告落地（重大重组与并购 7 条、对外投资与转型 0 条），资本运作与主业调整之间存在直接关联。

\begin{center}\small\setlength{\tabcolsep}{4pt}
\captionof{table}[国投丰乐（000713）历史融资与资本运作]{国投丰乐（000713）历史融资与资本运作（据公司公告与公开报道整理；金额为披露值，含首次公开发行、再融资、债券、银行借款与重整投资等）}
\begin{adjustbox}{max width=\textwidth}
\begin{tabular}{@{}p{1.45cm}p{2.55cm}p{4.05cm}p{4.95cm}@{}}
\toprule
时间 & 方式 & 规模 & 用途与说明 \\
\midrule
1997-04-22 & IPO（深交所） & 发行价 6.75 元/\allowbreak{}股 & 公司为中国种业首家上市公司 \\
2024-04-29 & 控股股东协议受让（控制权变更，非上市公司融资） & 国投种业以现金受让合肥建投持有的 1.23 亿股无限售流通股（占股份总数 20.00\%） & 取得上市公司控股权；2024 年 1 月国投种业收到市场监管总局经营者集中反垄断审查不实施进一步审查决定书（反执二审查决定〔2024〕31 号）；2024 年 3 月 25 日国务院国资委出具国资产权〔2024〕106 号批复原则同意 \\
2025-10-17 & 要约收购义务豁免安排 & 国投种业认购本次发行股票将触发要约收购义务 & 顺利完成向控股股东定增；发行完成后国投种业持有股份比例超过 30\% \\
2025-12-25 & 向特定对象发行A股股票（发行对象为控股股东国投种业，以现金认购） & 发行 1.84 亿股，\allowbreak{}发行价 5.91 元/\allowbreak{}股，\allowbreak{}募集资金总额 10.89 亿元 & 全部用于补充流动资金及偿还银行借款；募集资金于 2025 年 12 月 26 日到账；2026 年 1 月 21 日新增股份上市；发行完成后国投种业持股比例由 20.00\% 升至 38.46\% \\
2026-06-25 & 现金收购关联方股权（关联交易） & 以现金支付方式收购控股股东国投种业持有的国投（张掖）金种科技有限公司 60\% 股权 & 成为国投金种控股股东并纳入合并报表；国投金种另由张掖现代种业集团有限公司持有 40\%；国投种业直接持有公司 38.46\% 股权 \\
\bottomrule
\end{tabular}
\end{adjustbox}
\end{center}

\pfl{历史融投资情况} 从现金流量表看，2025 年公司取得借款收到的现金 2.53 亿元、偿还债务支付的现金 1.70 亿元、吸收权益性投资 10.81 亿元，筹资活动现金流净额 11.42 亿元（2023 年为取得借款 1.65 亿元、偿还债务 1.61 亿元）。 2026 年上半年筹资活动现金流净额为 0.42 亿元（取得借款 1.48 亿元、偿还债务 0.76 亿元，未年化），仍处于净融入状态。 2025 年分配股利、利润或偿付利息支付的现金合计 0.24 亿元。 公开披露的融资记录共 5 笔，工具类型以 IPO、定向增发为主；金额与用途见下表。 其中规模最大的两笔为：2026 年 6 月，现金收购关联方股权，以现金支付方式收购控股股东国投种业持有的国投（张掖）金种科技有限公司 60\% 股权。 股本方面，近三年披露股本变动 3 次；总股本自 2021 年末以来未发生变化（6.14 亿股）。 分红方面，近三年现金分红 4 次、合计每股派现 0.09 元，最近一次实施在 2025 年 12 月。 近三年“再融资”类公告 45 条，涉及定向增发、募集资金使用。 综合看，公司融资结构上以债务滚动为主、股权融资为补充；2025 年筹资活动现金流净额 11.42 亿元，年度内仍为净融入，与前述经营与投资现金流的方向基本一致。

\pfl{未来潜在融资需求分析} 2026 年 6 月末货币资金 15.50 亿元，而短期债务 2.38 亿元（现金短债比 6.52 倍），2026 年上半年经营现金流净额 -0.97 亿元。账面货币资金可以覆盖短期债务，缺口不体现在静态偿付层面。 公司 2026 年上半年归母净利润为负（-0.31 亿元），但 6 月末资产负债率 32.1\% 尚未进入第一档（亏损且负债率 $>$65\%）的区间，当前融资需求以补充流动资金、支撑低谷期经营性支出为主。 公开披露的后续资金安排线索包括：股权融资线索：2026 年 4 月，发行数量不超过 1.84 亿股，募集资金总额不超过 10.89 亿元（含本数）；2025 年度，另支付发行费用合计 307.55 万元（不含税）；债务与授信安排：2026 年 8 月，公司募集资金用途明确为补充流动资金及偿还银行借款。 资金需求还要叠加在建工程：期末在建工程 1.95 亿元、2026 年上半年购建固定资产等支出 0.45 亿元（未年化），这部分支出在周期底部通常会被主动放缓。 静态偿付不存在缺口，后续融资更可能服务于产能或并购安排，而不是填补流动性窟窿。




\subsubsection{新赛股份（600540）}
\label{co:600540}

\pfl{公司基本情况} 2003 年设立至今，新赛股份是新疆兵团系棉花全产业链企业（皮棉、棉籽精深加工、棉油棉蛋白），2004 年上交所上市、历史上曾被实施退市风险警示（*ST 新赛）；2023 年 9 月控股股东变更为中新建物流、实际控制人为兵团国资委；2024—2025 年连续亏损，2026 年上半年扭亏且净利同比增逾 18 倍，同时处于证监会立案调查期间。 公司于 2004-01-07 在上交所首发上市，注册资本 5.81 亿元，总股本 5.81 亿股。 截至 2026 年 9 月 24 日收盘价 5.43 元，流通市值 31.6 亿元，近一年+16.8\%，流通市值在三级行业“其他种植业”的 6 家公司中按由大到小排第\zhnumber{5}位，近一年涨跌幅高于全样本中位数（-10.1\%）26.9 个百分点。 股权结构方面，控股股东为新疆艾比湖农工商联合企业总公司。

\pfl{历史股价} 本报告样本区间（2021 年 1 月 至 2026 年 9 月）内，股价由 3.95 元变为 5.43 元，累计 +37.5\%，区间最高 7.61 元（2021 年 3 月）、最低 3.31 元（2026 年 6 月）；当前价相当于区间最高价的 71.4\%，为区间最低价的 1.64 倍。 月度收益的年化波动率 48.3\%，高于全样本中位数 37.4\%，在三级行业“其他种植业”的 6 家公司中波动率由高到低排第\zhnumber{1}位。

\begin{center}
\begin{minipage}{\textwidth}\centering
\begin{tikzpicture}
\begin{axis}[
  width=12.6cm, height=3.9cm, scale only axis,
  xmin=2020.922, xmax=2026.888, ymin=3.08, ymax=8.14,
  xtick={2021,2022,2023,2024,2025,2026},
  yearticks,
  yticklabel style={font=\scriptsize},
  ylabel={元/股}, ylabel style={font=\scriptsize},
  grid=major, grid style={LightGray, line width=0.3pt},
  axis line style={InkGray, line width=0.5pt},
  every axis plot/.append style={line width=0.9pt},
]
\addplot[AgriGreen] table[x=x,y=y]{data/clean/profiles/px_600540.dat};
\addplot[SoftRed, dashed, line width=0.7pt] table[x=x,y=y]{data/clean/profiles/pxzz_600540.dat};
\addplot[only marks, mark=*, mark size=1.1pt, SoftRed] table[x=x,y=y]{data/clean/profiles/pxzz_600540.dat};
\end{axis}
\end{tikzpicture}

\captionof{figure}[新赛股份（600540）股价与周期骨架]{新赛股份（600540）月度收盘价（元/股）与 ZigZag 周期骨架（阈值 25\%，圆点为周期转折点）}
\end{minipage}
\end{center}

\pfl{过去周期分析} 以 25\% 的 ZigZag 阈值划分，样本区间内股价形成 5 个主升段与 4 个主跌段。 最大主升段出现在 2021 年 1 月 至 2021 年 3 月，2 个月+92.7\%；最大主跌段为 2023 年 9 月 至 2024 年 6 月，9 个月-47.1\%。 最近一次转折出现在 2026 年 9 月（上涨段自 2026 年 6 月起，3 个月+64.0\%），当前处于该段之后的震荡之中。 公司股价与申万农林牧渔指数几乎同步见底，个股并未走出独立行情。 该子行业的周期来源与猪周期不同（第 \ref{sec:grain} 章、第 \ref{sec:soft} 章），价格更多由自身供给、进口政策与全球定价决定，公司 2026 年 6 月末资产负债率 82.7\%，高于全样本中位数 52.2\%，其股价因此更多反映资金面与信用风险，而不是价格弹性本身。 综合区间形态看，该股单段涨跌幅度偏大、以趋势段为主（最大主升 +92.7\%、最大主跌 -47.1\%），年化波动率 48.3\% 与全样本中位数 37.4\% 相比波动明显大于样本中位数。

\pfl{盈利情况} 2026 年上半年营业收入 25.71 亿元（同比 +77.2\%），归母净利润 1.25 亿元，同比 +1819.7\%。 以年报为趋势对照，2023---2025 年营业收入由 9.72 亿元变为 30.12 亿元（两年年均复合增速 +76.1\%），归母净利润由 0.15 亿元变为 -1.42 亿元。 按半年报口径，2026 年上半年的 ROE 为 30.8\%（未年化）、销售净利率 4.9\%、毛利率 8.7\%，ROE 在“其他种植业”6 家公司中按数值由高到低排第\zhnumber{1}位。 同期全样本半年 ROE 中位数 0.75\%，公司与之相差 30.10 个百分点（略好于中位数公司）。 综合看，公司在报告期维持盈利（高于全样本半年 ROE 中位数 0.75\%），利润的可持续性取决于其他种植业景气的延续与成本管控的执行。

\begin{center}\small\setlength{\tabcolsep}{3.4pt}
\captionof{table}[新赛股份（600540）关键财务指标]{新赛股份（600540）关键财务指标（合并报表口径；两个半年报期均未年化；现金短债比 $=$ 货币资金 $\div$ 短期债务）}
\begin{adjustbox}{max width=\textwidth}
\begin{tabular}{@{}lrrrrr@{}}
\toprule
指标 & 2023 年 & 2024 年 & 2025 年 & 2025 年半年报 & 2026 年半年报 \\
\midrule
营业收入（亿元） & 9.72 & 15.55 & 30.12 & 14.51 & 25.71 \\
归母净利润（亿元） & 0.15 & -2.44 & -1.42 & 0.07 & 1.25 \\
毛利率（\%） & 3.4 & -4.4 & 0.5 & 2.3 & 8.7 \\
ROE（\%） & 2.0 & -40.2 & -34.4 & 1.3 & 30.8 \\
资产负债率（\%） & 75.1 & 89.3 & 95.1 & 86.2 & 82.7 \\
经营活动现金流净额（亿元） & -4.29 & -8.19 & -13.17 & 18.27 & 24.20 \\
货币资金（亿元） & 8.72 & 10.69 & 17.05 & 16.53 & 8.39 \\
有息负债合计（亿元） & 18.55 & 35.62 & 57.97 & 22.40 & 15.08 \\
现金短债比（倍） & 0.50 & 0.31 & 0.30 & 0.77 & 0.58 \\
\bottomrule
\end{tabular}
\end{adjustbox}
\end{center}

\begin{center}
\begin{minipage}{\textwidth}\centering
\begin{tikzpicture}
\begin{groupplot}[
  group style={group size=2 by 1, horizontal sep=0.60cm},
  width=6.85cm, height=4.60cm,
  tick label style={font=\tiny},
  title style={font=\scriptsize\heiti},
  yticklabel style={font=\tiny},
  ylabel style={font=\tiny},
  grid=major, grid style={LightGray, line width=0.3pt},
  axis line style={InkGray, line width=0.5pt},
  enlarge x limits={abs=0.8},
  legend style={font=\scriptsize, at={(0.5,-0.20)}, anchor=north,
                legend columns=2, draw=none, fill=none, column sep=6pt},
]
\nextgroupplot[
  ybar, bar width=4pt,
  ymin=-6.07, ymax=34.07,
  xtick={0,1,2,3,4,5.7},
  xticklabels={2021,2022,2023,2024,2025,2026H1},
  ylabel={亿元},
]
\addplot[fill=AgriGreen!75, draw=AgriGreen, bar shift=-2.6pt] table[x=i, y=rev]{data/clean/profiles/fin_600540.dat};
\addplot[fill=SoftRed!60, draw=SoftRed, bar shift=2.6pt] table[x=i, y=ni]{data/clean/profiles/fin_600540.dat};
\addplot[fill=AgriGreen!30, draw=AgriGreen, pattern=north east lines, bar shift=-2.6pt] table[x=x, y=rev]{data/clean/profiles/finh_600540.dat};
\addplot[fill=SoftRed!20, draw=SoftRed, pattern=north east lines, bar shift=2.6pt] table[x=x, y=ni]{data/clean/profiles/finh_600540.dat};
\legend{营业收入（年/半年）, 归母净利润（年/半年）}
\nextgroupplot[
  ymin=-51.0, ymax=108.6,
  xtick={0,1,2,3,4,5.7},
  xticklabels={2021,2022,2023,2024,2025,2026H1},
  ylabel={\%},
]
\addplot[AgriGold, mark=*, mark size=1.3pt] table[x=i, y=roe]{data/clean/profiles/fin_600540.dat};
\addplot[OilBlue, mark=square*, mark size=1.3pt] table[x=i, y=debt]{data/clean/profiles/fin_600540.dat};
\addplot[only marks, AgriGold, mark=*, mark size=2.0pt] table[x=x, y=roe]{data/clean/profiles/finh_600540.dat};
\addplot[only marks, OilBlue, mark=square*, mark size=2.0pt] table[x=x, y=debt]{data/clean/profiles/finh_600540.dat};
\legend{ROE, 资产负债率}
\end{groupplot}
\end{tikzpicture}
\begin{tikzpicture}
\begin{axis}[
  name=finbot,
  width=13.75cm, height=3.10cm, scale only axis,
  ymin=-5.80, ymax=64.92,
  xtick={0,1,2,3,4,5.7},
  xticklabels={2021,2022,2023,2024,2025,2026H1},
  tick label style={font=\tiny},
  yticklabel style={font=\tiny},
  ylabel={亿元}, ylabel style={font=\tiny},
  ybar, bar width=4pt,
  grid=major, grid style={LightGray, line width=0.3pt},
  axis line style={InkGray, line width=0.5pt},
  enlarge x limits={abs=0.8},
  legend style={font=\scriptsize, at={(0.5,-0.26)}, anchor=north,
                legend columns=2, draw=none, fill=none, column sep=10pt},
]
\addplot[fill=AgriGreen!55, draw=AgriGreen, bar shift=-3.0pt] table[x=i, y=cash]{data/clean/profiles/fin_600540.dat};
\addplot[fill=SoftRed!45, draw=SoftRed, bar shift=3.0pt] table[x=i, y=ibd]{data/clean/profiles/fin_600540.dat};
\addplot[fill=AgriGreen!25, draw=AgriGreen, pattern=north east lines, bar shift=-3.0pt] table[x=x, y=cash]{data/clean/profiles/finh_600540.dat};
\addplot[fill=SoftRed!20, draw=SoftRed, pattern=north east lines, bar shift=3.0pt] table[x=x, y=ibd]{data/clean/profiles/finh_600540.dat};
\legend{货币资金（左轴）, 有息负债（左轴）}
\end{axis}
\begin{axis}[
  at={(finbot.south east)}, anchor=south east,
  width=13.75cm, height=3.10cm, scale only axis,
  axis y line*=right, axis x line*=none, xtick=\empty,
  ymin=-22.89, ymax=35.42,
  tick label style={font=\tiny},
  yticklabel style={font=\tiny}, scaled y ticks=false,
  legend style={font=\scriptsize, at={(0.5,-0.44)}, anchor=north,
                legend columns=1, draw=none, fill=none},
]
\addplot[OilBlue, mark=*, mark size=2.2pt, line width=1.3pt] table[x=i, y=ocf]{data/clean/profiles/fin_600540.dat};
\addplot[only marks, OilBlue, mark=*, mark size=3.2pt] table[x=x, y=ocf]{data/clean/profiles/finh_600540.dat};
\legend{经营活动现金流净额（右轴，亿元，蓝点）}
\end{axis}
\end{tikzpicture}

\captionof{figure}[新赛股份（600540）近年主要财务指标]{新赛股份（600540）近年主要财务指标（左上：营业收入与归母净利润；右上：ROE 与资产负债率；下：货币资金、有息负债为左轴柱状，经营活动现金流净额为其右轴的蓝点折线。
2021---2025 年为年报口径，2026H1 为 2026 年半年报/半年末，与其前一年报之间留出间隔、以斜纹或空心点区分；半年报数据未年化）}
\end{minipage}
\end{center}

\pfl{财务分析} 2026 年 6 月末资产负债率 82.7\%，较 2025 年末下降 12.4 个百分点（样本中位数 52.2\%，所处三级行业内按由低到高排第\zhnumber{6}位）。 2026 年 6 月末货币资金 8.39 亿元、短期债务（短期借款与一年内到期非流动负债）14.35 亿元、长期债务 0.73 亿元，有息负债合计 15.08 亿元，现金短债比 0.58 倍——覆盖偏紧。 净有息负债 6.69 亿元，相当于其 2026 年上半年营业收入的 26\%。 流动比率 0.79、速动比率 0.70，短期偿付压力较大。 2026 年上半年经营活动现金流净额 24.20 亿元，经营现金流与利润同向为正，内生资金可以支撑运营。 综合看，现金短债比 0.58 倍，短期资金安排偏紧；资产负债率高于样本中位数 30.5 个百分点；有息负债中短期债务占比更高，期限结构仍以滚动续作为主。 公司披露的监管与合规风险为：2026 年 6 月，收到中国证监会立案告知书；2026 年 8 月，上海证券交易所对公司 2025 年年度报告信息披露监管问询函回复。

\pfl{重大战略转型} 公司的战略动作可归为以下几类：产能与项目（2 项）：2023 年 7 月，公司披露控股股东和实际控制人拟发生变更的提示性公告；2023 年 9 月，中新建物流增资完成后注册资本至 68.24 亿元；其他动作（2 项）：2026 年上半年，截至 6 月末子企业亏损户数较 2025 年末明显减少、亏损面大幅收窄。 公告口径上，近三年（2023 年 9 月---2026 年 9 月）公司披露重大重组与并购类公告 1 条、对外投资与转型类公告 0 条、回购与股权激励类公告 0 条，合计公告 422 条。 综合判断，报告期内公司有 1 条与战略相关的公告落地（重大重组与并购 1 条、对外投资与转型 0 条），资本运作与主业调整之间存在直接关联。

\begin{center}\small\setlength{\tabcolsep}{4pt}
\captionof{table}[新赛股份（600540）历史融资与资本运作]{新赛股份（600540）历史融资与资本运作（据公司公告与公开报道整理；金额为披露值，含首次公开发行、再融资、债券、银行借款与重整投资等）}
\begin{adjustbox}{max width=\textwidth}
\begin{tabular}{@{}p{1.45cm}p{2.55cm}p{4.05cm}p{4.95cm}@{}}
\toprule
时间 & 方式 & 规模 & 用途与说明 \\
\midrule
2003-12 & IPO（上交所，二级市场投资者定价配售） & 发行 5，\allowbreak{}000 万股，\allowbreak{}发行价 6.67 元/\allowbreak{}股，\allowbreak{}募集资金总额 3.33 亿元 & 首次公开发行募集资金（其中记入股本 5，\allowbreak{}000.00 万元、\allowbreak{}资本公积 2.64 亿元）；网上配号总数 71 \\
2014-11-20 & 非公开发行（向特定投资者，主承销商广州证券） & 发行人民币普通股（A 股）5，\allowbreak{}954 万股，\allowbreak{}每股发行价格 8.70 元 & 非公开发行募投项目（后续部分用途变更）；经中国证监会核准 \\
2018 & 银行借款与募集资金暂时补流 & 全年通过银行贷款融资 7.92 亿元（其中协助棉业公司 9 & 满足分子公司在用款高峰期的资金需求、\allowbreak{}缓解企业资金压力；同期全年实际收回各类应收款项 2.44 亿元 \\
2023-08 & 变更部分募集资金用途 & 将原募投项目霍城县可利煤炭物流配送有限公司专用线扩建项目部分募集资金用途变更为偿还总部银行贷款项目，\allowbreak{}600.00 万元 & 偿还总部银行贷款；其余部分募集资金用途未发生改变 \\
2026-04-28 & 借款及担保额度议案（2026年度） & 预计 2026 年度公司及合并报表范围内子公司所需借款额度 79.00 亿元（可上浮 15\%）；拟提供合计不超过 35.00 亿元的担保（含现有担保、\allowbreak{}展期或续保及新增担保 & 生产经营及子公司资金需求；2026 年 5 月 21 日 2025 年年度股东会审议通过；截至 2026 年 4 月 28 日公司及控股子公司对外担保总额 35.00 亿元 \\
\bottomrule
\end{tabular}
\end{adjustbox}
\end{center}

\pfl{历史融投资情况} 从现金流量表看，2025 年公司取得借款收到的现金 79.41 亿元、偿还债务支付的现金 54.78 亿元，筹资活动现金流净额 23.71 亿元（2023 年为取得借款 25.78 亿元、偿还债务 19.46 亿元）。 2026 年上半年筹资活动现金流净额为 -43.34 亿元（取得借款 8.81 亿元、偿还债务 51.67 亿元，未年化），已转为净流出，处于净还债状态。 2025 年分配股利、利润或偿付利息支付的现金合计 0.91 亿元。 公开披露的融资记录共 5 笔，工具类型以 IPO、定向增发、银行借款为主；金额与用途见下表。 其中规模最大的两笔为：2026 年 4 月，银行借款，预计 2026 年度公司及合并报表范围内子公司所需借款额度 79.00 亿元（可上浮 15\%），拟提供合计不超过 35.00 亿元的担保（含现有担保、展期或续保及新增担保；2023 年 8 月，变更部分募集资金用途，将原募投项目霍城县可利煤炭物流配送有限公司专用线扩建项目部分募集资金用途变更为偿还总部银行贷款项目，600.00 万元。 股本方面，近三年披露股本变动 3 次；总股本由 2021 年末的 4.71 亿股变为 5.81 亿股（+23.45\%）。 分红方面，近三年未实施现金分红，在全部样本中属于分红能力偏弱的一端。 近三年“再融资”类公告 1 条，涉及定向增发。 综合看，公司融资结构上股权与债务融资并举；2025 年筹资活动现金流净额 23.71 亿元，年度内仍为净融入，与前述经营与投资现金流的方向基本一致。

\pfl{未来潜在融资需求分析} 2026 年 6 月末货币资金 8.39 亿元，而短期债务 14.35 亿元（现金短债比 0.58 倍），2026 年上半年经营现金流净额 24.20 亿元。按“短期债务减去账面货币资金”的滚动偿付口径，静态缺口约 5.96 亿元；上半年盈利或接近盈亏平衡，该缺口不随经营亏损进一步扩大。 公司 2026 年 6 月末资产负债率 82.7\%、2026 年上半年 ROE 30.8\%（未年化），资产负债率偏高（82.7\%），融资需求首先用于置换与滚动存量债务、补充营运资金，属于“补血型”而非扩张型。 公开披露的后续资金安排线索包括：债务与授信安排：2026 年度，公司及控股子公司借款额度 79.00 亿元、担保额度 35.00 亿元（均可上浮 15\%）公司对外部债务融资依赖度高；2026 年 4 月，公告公司及子公司申请银行综合授信额度部分募集资金投资项目延期对子公司流动资金借款担保计划。 综合看，公司的外部资金需求以补充流动性与项目投入为主，滚动偿付口径下的静态缺口约 5.96 亿元，能否落地取决于授信续作与发行窗口的配合。




\subsubsection{新农开发（600359）}
\label{co:600359}

\pfl{公司基本情况} 1999 年设立至今，新农开发是新疆生产建设兵团第一师国资委控制的农业上市平台，主营作物种子、大米谷物加工、甘草饮品与肉牛养殖；2023 年剥离乳业，2025 年控股股东变更为农垦集团。 公司于 1999-04-29 在上交所首发上市，注册资本 3.82 亿元，总股本 3.82 亿股。 截至 2026 年 9 月 24 日收盘价 7.62 元，流通市值 29.1 亿元，近一年+11.9\%，流通市值在三级行业“其他种植业”的 6 家公司中按由大到小排第\zhnumber{6}位，近一年涨跌幅高于全样本中位数（-10.1\%）22.0 个百分点。 按 2026 年半年报披露的行业构成，收入占比前三位为服务业及其他 72.7\%；农业 20.6\%；工业 6.7\%。 与 2021 年年报相比，第一大行业口径由“农业”（55.1\%）变为“服务业及其他”（72.7\%）；两期披露的分类口径有调整，占比差异中同时包含分类因素。 股权结构方面，控股股东为阿拉尔塔河投资有限责任公司；实际控制人为仍为新疆生产建设兵团第一师国资委。

\pfl{历史股价} 以 2021 年 1 月为起点、2026 年 9 月为终点，股价由 7.68 元变为 7.62 元，累计 -0.8\%，区间最高 10.11 元（2022 年 5 月）、最低 5.40 元（2026 年 6 月）；当前价相当于区间最高价的 75.4\%，为区间最低价的 1.41 倍。 月度收益的年化波动率 29.2\%，低于全样本中位数 37.4\%，在三级行业“其他种植业”的 6 家公司中波动率由高到低排第\zhnumber{4}位。 把股价阶段与公司事件对齐看，2024 年 6 月 至 2026 年 1 月股价上涨+41.9\%，同期 2025 年半年报，公司披露截至报告期末尚有 3，176.09 万元新农乳业股权转让款未收回。 从时间对应关系看，公司层面的资本运作与股权、风险事件解释了区间内多数急涨急跌，与行业指数的同步性相对有限。

\begin{center}
\begin{minipage}{\textwidth}\centering
\begin{tikzpicture}
\begin{axis}[
  width=12.6cm, height=3.9cm, scale only axis,
  xmin=2020.922, xmax=2026.888, ymin=5.02, ymax=10.82,
  xtick={2021,2022,2023,2024,2025,2026},
  yearticks,
  yticklabel style={font=\scriptsize},
  ylabel={元/股}, ylabel style={font=\scriptsize},
  grid=major, grid style={LightGray, line width=0.3pt},
  axis line style={InkGray, line width=0.5pt},
  every axis plot/.append style={line width=0.9pt},
]
\addplot[AgriGreen] table[x=x,y=y]{data/clean/profiles/px_600359.dat};
\addplot[SoftRed, dashed, line width=0.7pt] table[x=x,y=y]{data/clean/profiles/pxzz_600359.dat};
\addplot[only marks, mark=*, mark size=1.1pt, SoftRed] table[x=x,y=y]{data/clean/profiles/pxzz_600359.dat};
\end{axis}
\end{tikzpicture}

\captionof{figure}[新农开发（600359）股价与周期骨架]{新农开发（600359）月度收盘价（元/股）与 ZigZag 周期骨架（阈值 25\%，圆点为周期转折点）}
\end{minipage}
\end{center}

\pfl{过去周期分析} 以 25\% 的 ZigZag 阈值划分，样本区间内股价形成 3 个主升段与 2 个主跌段。 最大主跌段为 2022 年 5 月 至 2024 年 6 月，25 个月-44.7\%；最大主升段出现在 2024 年 6 月 至 2026 年 1 月，19 个月+41.9\%。 最近一次转折出现在 2026 年 9 月（上涨段自 2026 年 6 月起，3 个月+41.1\%），当前处于该段之后的震荡之中。 公司股价与申万农林牧渔指数几乎同步见底，个股并未走出独立行情。 该子行业的周期来源与猪周期不同（第 \ref{sec:grain} 章、第 \ref{sec:soft} 章），价格更多由自身供给、进口政策与全球定价决定，公司 2026 年上半年实现盈利、半年 ROE 2.8\%（未年化），好于全样本中位数（0.75\%），下行周期对其报表的冲击小于同业。 综合区间形态看，该股单段涨跌幅度相对温和、以震荡为主（最大主升 +41.9\%、最大主跌 -44.7\%），年化波动率 29.2\% 与全样本中位数 37.4\% 相比波动小于样本中位数。

\pfl{盈利情况} 2026 年上半年营业收入 1.96 亿元（同比 -31.1\%），归母净利润 0.19 亿元，同比 -1.9\%。 以年报为趋势对照，2023---2025 年营业收入由 5.48 亿元变为 4.84 亿元（两年年均复合增速 -6.1\%），归母净利润由 0.69 亿元变为 -0.08 亿元。 按半年报口径，2026 年上半年的 ROE 为 2.8\%（未年化）、销售净利率 10.3\%、毛利率 30.5\%，ROE 在“其他种植业”6 家公司中按数值由高到低排第\zhnumber{3}位。 按同一口径，三级行业“其他种植业”6 家公司的半年 ROE 中位数为 1.70\%，该公司高于其中位数 1.11 个百分点。 综合看，公司在报告期维持盈利（高于全样本半年 ROE 中位数 0.75\%），利润的可持续性取决于其他种植业景气的延续与成本管控的执行。

\begin{center}\small\setlength{\tabcolsep}{3.4pt}
\captionof{table}[新农开发（600359）关键财务指标]{新农开发（600359）关键财务指标（合并报表口径；两个半年报期均未年化；现金短债比 $=$ 货币资金 $\div$ 短期债务）}
\begin{adjustbox}{max width=\textwidth}
\begin{tabular}{@{}lrrrrr@{}}
\toprule
指标 & 2023 年 & 2024 年 & 2025 年 & 2025 年半年报 & 2026 年半年报 \\
\midrule
营业收入（亿元） & 5.48 & 6.22 & 4.84 & 2.84 & 1.96 \\
归母净利润（亿元） & 0.69 & 0.22 & -0.08 & 0.19 & 0.19 \\
毛利率（\%） & 26.7 & 23.0 & 15.6 & 22.8 & 30.5 \\
ROE（\%） & 11.3 & 3.3 & -1.2 & 2.8 & 2.8 \\
资产负债率（\%） & 47.0 & 46.3 & 54.1 & 51.8 & 55.6 \\
经营活动现金流净额（亿元） & -0.47 & 0.95 & 0.35 & 0.76 & -0.26 \\
货币资金（亿元） & 2.53 & 3.42 & 2.74 & 3.84 & 3.42 \\
有息负债合计（亿元） & 2.11 & 4.15 & 5.16 & 5.11 & 6.52 \\
现金短债比（倍） & 1.22 & 0.97 & 0.77 & 0.94 & 0.70 \\
\bottomrule
\end{tabular}
\end{adjustbox}
\end{center}

\begin{center}
\begin{minipage}{\textwidth}\centering
\begin{tikzpicture}
\begin{groupplot}[
  group style={group size=2 by 1, horizontal sep=0.60cm},
  width=6.85cm, height=4.60cm,
  tick label style={font=\tiny},
  title style={font=\scriptsize\heiti},
  yticklabel style={font=\tiny},
  ylabel style={font=\tiny},
  grid=major, grid style={LightGray, line width=0.3pt},
  axis line style={InkGray, line width=0.5pt},
  enlarge x limits={abs=0.8},
  legend style={font=\scriptsize, at={(0.5,-0.20)}, anchor=north,
                legend columns=2, draw=none, fill=none, column sep=6pt},
]
\nextgroupplot[
  ybar, bar width=4pt,
  ymin=-0.77, ymax=7.69,
  xtick={0,1,2,3,4,5.7},
  xticklabels={2021,2022,2023,2024,2025,2026H1},
  ylabel={亿元},
]
\addplot[fill=AgriGreen!75, draw=AgriGreen, bar shift=-2.6pt] table[x=i, y=rev]{data/clean/profiles/fin_600359.dat};
\addplot[fill=SoftRed!60, draw=SoftRed, bar shift=2.6pt] table[x=i, y=ni]{data/clean/profiles/fin_600359.dat};
\addplot[fill=AgriGreen!30, draw=AgriGreen, pattern=north east lines, bar shift=-2.6pt] table[x=x, y=rev]{data/clean/profiles/finh_600359.dat};
\addplot[fill=SoftRed!20, draw=SoftRed, pattern=north east lines, bar shift=2.6pt] table[x=x, y=ni]{data/clean/profiles/finh_600359.dat};
\legend{营业收入（年/半年）, 归母净利润（年/半年）}
\nextgroupplot[
  ymin=-7.0, ymax=78.6,
  xtick={0,1,2,3,4,5.7},
  xticklabels={2021,2022,2023,2024,2025,2026H1},
  ylabel={\%},
]
\addplot[AgriGold, mark=*, mark size=1.3pt] table[x=i, y=roe]{data/clean/profiles/fin_600359.dat};
\addplot[OilBlue, mark=square*, mark size=1.3pt] table[x=i, y=debt]{data/clean/profiles/fin_600359.dat};
\addplot[only marks, AgriGold, mark=*, mark size=2.0pt] table[x=x, y=roe]{data/clean/profiles/finh_600359.dat};
\addplot[only marks, OilBlue, mark=square*, mark size=2.0pt] table[x=x, y=debt]{data/clean/profiles/finh_600359.dat};
\legend{ROE, 资产负债率}
\end{groupplot}
\end{tikzpicture}
\begin{tikzpicture}
\begin{axis}[
  name=finbot,
  width=13.75cm, height=3.10cm, scale only axis,
  ymin=-0.76, ymax=8.53,
  xtick={0,1,2,3,4,5.7},
  xticklabels={2021,2022,2023,2024,2025,2026H1},
  tick label style={font=\tiny},
  yticklabel style={font=\tiny},
  ylabel={亿元}, ylabel style={font=\tiny},
  ybar, bar width=4pt,
  grid=major, grid style={LightGray, line width=0.3pt},
  axis line style={InkGray, line width=0.5pt},
  enlarge x limits={abs=0.8},
  legend style={font=\scriptsize, at={(0.5,-0.26)}, anchor=north,
                legend columns=2, draw=none, fill=none, column sep=10pt},
]
\addplot[fill=AgriGreen!55, draw=AgriGreen, bar shift=-3.0pt] table[x=i, y=cash]{data/clean/profiles/fin_600359.dat};
\addplot[fill=SoftRed!45, draw=SoftRed, bar shift=3.0pt] table[x=i, y=ibd]{data/clean/profiles/fin_600359.dat};
\addplot[fill=AgriGreen!25, draw=AgriGreen, pattern=north east lines, bar shift=-3.0pt] table[x=x, y=cash]{data/clean/profiles/finh_600359.dat};
\addplot[fill=SoftRed!20, draw=SoftRed, pattern=north east lines, bar shift=3.0pt] table[x=x, y=ibd]{data/clean/profiles/finh_600359.dat};
\legend{货币资金（左轴）, 有息负债（左轴）}
\end{axis}
\begin{axis}[
  at={(finbot.south east)}, anchor=south east,
  width=13.75cm, height=3.10cm, scale only axis,
  axis y line*=right, axis x line*=none, xtick=\empty,
  ymin=-0.88, ymax=1.62,
  tick label style={font=\tiny},
  yticklabel style={font=\tiny}, scaled y ticks=false,
  legend style={font=\scriptsize, at={(0.5,-0.44)}, anchor=north,
                legend columns=1, draw=none, fill=none},
]
\addplot[OilBlue, mark=*, mark size=2.2pt, line width=1.3pt] table[x=i, y=ocf]{data/clean/profiles/fin_600359.dat};
\addplot[only marks, OilBlue, mark=*, mark size=3.2pt] table[x=x, y=ocf]{data/clean/profiles/finh_600359.dat};
\legend{经营活动现金流净额（右轴，亿元，蓝点）}
\end{axis}
\end{tikzpicture}

\captionof{figure}[新农开发（600359）近年主要财务指标]{新农开发（600359）近年主要财务指标（左上：营业收入与归母净利润；右上：ROE 与资产负债率；下：货币资金、有息负债为左轴柱状，经营活动现金流净额为其右轴的蓝点折线。
2021---2025 年为年报口径，2026H1 为 2026 年半年报/半年末，与其前一年报之间留出间隔、以斜纹或空心点区分；半年报数据未年化）}
\end{minipage}
\end{center}

\pfl{财务分析} 2026 年 6 月末资产负债率 55.6\%，较 2025 年末上升 1.5 个百分点（样本中位数 52.2\%，所处三级行业内按由低到高排第\zhnumber{2}位）。 2026 年 6 月末货币资金 3.42 亿元、短期债务（短期借款与一年内到期非流动负债）4.87 亿元、长期债务 1.66 亿元，有息负债合计 6.52 亿元，现金短债比 0.70 倍——覆盖偏紧。 净有息负债 3.10 亿元，相当于其 2026 年上半年营业收入的 159\%。 流动比率 1.15、速动比率 0.62，短期指标尚可。 2026 年上半年经营活动现金流净额 -0.26 亿元，净利润为正但经营现金流为负，利润的现金含量偏低。 2026 年 6 月末在建工程 2.27 亿元，较上年末增加 0.45 亿元，仍有在建投入。 综合看，现金短债比 0.70 倍，短期资金安排偏紧；资产负债率高于样本中位数 3.4 个百分点；有息负债中短期债务占比更高，期限结构仍以滚动续作为主。

\pfl{重大战略转型} 从公开披露看，公司的战略推进集中在：产能与项目（1 项）：2026 年度，公司拟投资建设项目 6 个、总投资 4.14 亿元；并购与重组（1 项）：2026 年 6 月，培育 X 类新兴增长产业、剥离非优势产业。 公告口径上，近三年（2023 年 9 月---2026 年 9 月）公司披露重大重组与并购类公告 7 条、对外投资与转型类公告 0 条、回购与股权激励类公告 0 条，合计公告 312 条。 主营结构的口径变化已经落到报表上：2021 年年报的第一大行业为“农业”（55.1\%），2026 年半年报为“服务业及其他”（72.7\%）；公司在这两期的分部划分方式有过调整。 综合判断，报告期内公司有 7 条与战略相关的公告落地（重大重组与并购 7 条、对外投资与转型 0 条），资本运作与主业调整之间存在直接关联。

\begin{center}\small\setlength{\tabcolsep}{4pt}
\captionof{table}[新农开发（600359）历史融资与资本运作]{新农开发（600359）历史融资与资本运作（据公司公告与公开报道整理；金额为披露值，含首次公开发行、再融资、债券、银行借款与重整投资等）}
\begin{adjustbox}{max width=\textwidth}
\begin{tabular}{@{}p{1.45cm}p{2.55cm}p{4.05cm}p{4.95cm}@{}}
\toprule
时间 & 方式 & 规模 & 用途与说明 \\
\midrule
1999-04 & IPO（上海证券交易所，网上定价发行） & 发行 9，\allowbreak{}000 万股，\allowbreak{}发行价 3.86 元/\allowbreak{}股，\allowbreak{}募集资金 34，\allowbreak{}740 万元 & 司 \\
2001-09 & 配股 & 配售 2,\allowbreak{}700 万股 & 经证监公司字[2001]89 号文核准 \\
2008-10 & 非公开发行 A 股股票（预案） & 拟发行 4，\allowbreak{}000 万—8，\allowbreak{}000 万股，\allowbreak{}发行价格不低于 6.82 元/\allowbreak{}股 & 8 万吨棉浆粕生产线项目 1.59 亿元（与山东海龙按 55\%/\allowbreak{}45\% 共同投入）、\allowbreak{}日处理 900 吨鲜奶奶粉加工生产线建设项目 2.34 亿元；预案经第三届董事会第 24 次会议通过 \\
2014-10 & 非公开发行股票 & 6,\allowbreak{}051.28 万股，\allowbreak{}发行价格 9.75 元/\allowbreak{}股（按此计算约 5.90 亿元） & 2014-12 完成验资并出具 XYZH/\allowbreak{}2014URA1028 号验资报告 \\
2023-04 & 重大资产出售（现金交易） & 3.18 亿元 & 向天润乳业出售所持新农乳业 97.4359\% 股权 \\
2025-01 & 现金收购股权（关联交易） & 3,\allowbreak{}944.76 万元 & 收购控股股东农垦集团所持天山雪 100\% 股权；以自有资金分三期支付（80\%/\allowbreak{}10\%/\allowbreak{}10\%）；2025 年半年报披露已支付 80\% 即 3 \\
\bottomrule
\end{tabular}
\end{adjustbox}
\end{center}

\pfl{历史融投资情况} 从现金流量表看，2025 年公司取得借款收到的现金 5.60 亿元、偿还债务支付的现金 4.59 亿元，筹资活动现金流净额 0.57 亿元（2023 年为取得借款 2.58 亿元、偿还债务 4.56 亿元）。 2026 年上半年筹资活动现金流净额为 1.23 亿元（取得借款 3.34 亿元、偿还债务 2.03 亿元，未年化），仍处于净融入状态。 2025 年分配股利、利润或偿付利息支付的现金合计 0.12 亿元。 公开披露的融资记录共 6 笔，工具类型以 IPO、配股、定向增发为主；金额与用途见下表。 其中规模最大的两笔为：2023 年 4 月，重大资产出售，3.18 亿元；1999 年 4 月，IPO，发行 9，000 万股，发行价 3.86 元/股，募集资金 34，740 万元。 股本方面，近三年披露股本变动 3 次；总股本自 2021 年末以来未发生变化（3.82 亿股）。 分红方面，近三年未实施现金分红，在全部样本中属于分红能力偏弱的一端。 综合看，公司融资结构上股权与债务融资并举；2025 年筹资活动现金流净额 0.57 亿元，年度内仍为净融入，与前述经营与投资现金流的方向基本一致。

\pfl{未来潜在融资需求分析} 2026 年 6 月末货币资金 3.42 亿元，而短期债务 4.87 亿元（现金短债比 0.70 倍），2026 年上半年经营现金流净额 -0.26 亿元。按“短期债务减去账面货币资金”的滚动偿付口径，静态缺口约 1.44 亿元；上半年盈利或接近盈亏平衡，该缺口不随经营亏损进一步扩大。 公司 2026 年 6 月末资产负债率 55.6\%、2026 年上半年 ROE 2.8\%（未年化），现金短债比 0.70 倍、短债覆盖不足，融资需求首先用于置换与滚动存量债务、补充营运资金，属于“补血型”而非扩张型。 公开披露的后续资金安排线索包括：债务与授信安排：2026 年度，公司投资计划中部分项目建设资金来源明确为“企业自筹、银行贷款”；2026 年，第一次临时股东会审议通过 2026 年度贷款计划 2026 年度预计为子公司借款提供担保 2026 年度为控股子公司提供财务资助预计 2026 年度日常关联交易。 资金需求还要叠加在建工程：期末在建工程 2.27 亿元、2026 年上半年购建固定资产等支出 0.29 亿元（未年化），这部分支出在周期底部通常会被主动放缓。 综合看，公司的外部资金需求以营运资金周转与在建项目投入为主，滚动偿付口径下的静态缺口约 1.44 亿元，融资节奏将受制于银行续贷意愿与股权融资窗口。




\subsubsection{雪榕生物（300511）}
\label{co:300511}

\pfl{公司基本情况} 公司前身可追溯至 2016 年，国内产能最大的食用菌工厂化企业，金针菇日产能全国第一；2021—2024 年连亏、2025 年扭亏，2025 年 11 月控股股东与实控人变更为万紫千鸿、蒋智。 公司于 2016-05-04 在深交所创业板首发上市，注册资本 6.41 亿元，总股本 4.99 亿股。 截至 2026 年 9 月 24 日收盘价 4.84 元，流通市值 27.6 亿元，近一年-18.7\%，流通市值在三级行业“食用菌”的 4 家公司中按由大到小排第\zhnumber{2}位，近一年涨跌幅低于全样本中位数（-10.1\%）8.6 个百分点。 按 2025 年年报披露的行业构成，收入占比前两位为食用菌行业 98.9\%；其他(补充) 1.1\%。

\pfl{历史股价} 本报告样本区间（2021 年 1 月 至 2026 年 9 月）内，股价由 10.83 元变为 4.84 元，累计 -55.3\%，区间最高 11.29 元（2021 年 2 月）、最低 2.59 元（2024 年 6 月）；当前价相当于区间最高价的 42.9\%，为区间最低价的 1.87 倍。 月度收益的年化波动率 41.9\%，高于全样本中位数 37.4\%，在三级行业“食用菌”的 4 家公司中波动率由高到低排第\zhnumber{3}位。 把股价阶段与公司事件对齐看，2024 年 6 月 至 2026 年 4 月股价上涨+210.0\%，同期（2025 年 11 月）万紫千鸿取得控制权的目的被披露为基于对上市公司内在价值和未来发展前景的信心。 把这些阶段与公司事件对齐后可以看到，几轮大涨大跌多由公司自身的资本运作与风险事件触发，行业景气只解释了其中一部分。

\begin{center}
\begin{minipage}{\textwidth}\centering
\begin{tikzpicture}
\begin{axis}[
  width=12.6cm, height=3.9cm, scale only axis,
  xmin=2020.922, xmax=2026.888, ymin=2.41, ymax=12.08,
  xtick={2021,2022,2023,2024,2025,2026},
  yearticks,
  yticklabel style={font=\scriptsize},
  ylabel={元/股}, ylabel style={font=\scriptsize},
  grid=major, grid style={LightGray, line width=0.3pt},
  axis line style={InkGray, line width=0.5pt},
  every axis plot/.append style={line width=0.9pt},
]
\addplot[AgriGreen] table[x=x,y=y]{data/clean/profiles/px_300511.dat};
\addplot[SoftRed, dashed, line width=0.7pt] table[x=x,y=y]{data/clean/profiles/pxzz_300511.dat};
\addplot[only marks, mark=*, mark size=1.1pt, SoftRed] table[x=x,y=y]{data/clean/profiles/pxzz_300511.dat};
\end{axis}
\end{tikzpicture}

\captionof{figure}[雪榕生物（300511）股价与周期骨架]{雪榕生物（300511）月度收盘价（元/股）与 ZigZag 周期骨架（阈值 25\%，圆点为周期转折点）}
\end{minipage}
\end{center}

\pfl{过去周期分析} 以 25\% 的 ZigZag 阈值划分，样本区间内股价形成 2 个主升段与 3 个主跌段。 最大主跌段为 2023 年 7 月 至 2024 年 6 月，11 个月-66.8\%；最大主升段出现在 2024 年 6 月 至 2026 年 4 月，22 个月+210.0\%。 最近一次转折出现在 2026 年 9 月（下跌段自 2026 年 4 月起，5 个月-39.7\%），当前处于该段的延续之中。 公司股价较申万农林牧渔指数提前约 23 个月见底，是板块中较早定价周期反转的公司之一。 食用菌是工厂化生产的制造型农业（第 \ref{sec:soft} 章），产能投放节奏与金针菇等单品价格是盈利的主要变量，与养殖周期基本无关，公司 2026 年上半年实现盈利、半年 ROE 1.1\%（未年化），好于全样本中位数（0.75\%），下行周期对其报表的冲击小于同业。 综合区间形态看，该股单段涨跌幅度偏大、以趋势段为主（最大主升 +210.0\%、最大主跌 -66.8\%），年化波动率 41.9\% 与全样本中位数 37.4\% 相比波动与样本中位数接近。

\pfl{盈利情况} 2026 年上半年营业收入 11.56 亿元（同比 +46.2\%），归母净利润 0.18 亿元，实现扭亏。 以年报为趋势对照，2023---2025 年营业收入由 25.65 亿元变为 19.88 亿元（两年年均复合增速 -12.0\%），归母净利润由 -1.88 亿元变为 0.30 亿元。 按半年报口径，2026 年上半年的 ROE 为 1.1\%（未年化）、销售净利率 1.2\%、毛利率 17.2\%，ROE 在“食用菌”4 家公司中按数值由高到低排第\zhnumber{4}位。 按同一口径，三级行业“食用菌”4 家公司的半年 ROE 中位数为 2.15\%，该公司低于其中位数 1.05 个百分点。 从公司披露的原因看，业绩变动主要来自：2024 年度，公司主动收缩整体业务规模导致产能利用率下降、单位成本升高，并对相关资产及资产组计提减值（跌价）准备。 综合看，公司在报告期维持盈利（高于全样本半年 ROE 中位数 0.75\%），利润的可持续性取决于食用菌景气的延续与成本管控的执行。

\begin{center}\small\setlength{\tabcolsep}{3.4pt}
\captionof{table}[雪榕生物（300511）关键财务指标]{雪榕生物（300511）关键财务指标（合并报表口径；两个半年报期均未年化；现金短债比 $=$ 货币资金 $\div$ 短期债务）}
\begin{adjustbox}{max width=\textwidth}
\begin{tabular}{@{}lrrrrr@{}}
\toprule
指标 & 2023 年 & 2024 年 & 2025 年 & 2025 年半年报 & 2026 年半年报 \\
\midrule
营业收入（亿元） & 25.65 & 21.62 & 19.88 & 7.91 & 11.56 \\
归母净利润（亿元） & -1.88 & -6.17 & 0.30 & -1.03 & 0.18 \\
毛利率（\%） & 8.4 & 3.6 & 13.6 & 1.2 & 17.2 \\
ROE（\%） & -11.4 & -49.5 & 2.3 & -9.8 & 1.1 \\
资产负债率（\%） & 68.9 & 78.6 & 58.0 & 67.5 & 56.5 \\
经营活动现金流净额（亿元） & 3.58 & 4.37 & 4.16 & 1.38 & 2.21 \\
货币资金（亿元） & 2.53 & 1.68 & 1.61 & 0.76 & 1.76 \\
有息负债合计（亿元） & 22.37 & 21.52 & 12.20 & 15.21 & 12.52 \\
现金短债比（倍） & 0.20 & 0.14 & 0.16 & 0.11 & 0.17 \\
\bottomrule
\end{tabular}
\end{adjustbox}
\end{center}

\begin{center}
\begin{minipage}{\textwidth}\centering
\begin{tikzpicture}
\begin{groupplot}[
  group style={group size=2 by 1, horizontal sep=0.60cm},
  width=6.85cm, height=4.60cm,
  tick label style={font=\tiny},
  title style={font=\scriptsize\heiti},
  yticklabel style={font=\tiny},
  ylabel style={font=\tiny},
  grid=major, grid style={LightGray, line width=0.3pt},
  axis line style={InkGray, line width=0.5pt},
  enlarge x limits={abs=0.8},
  legend style={font=\scriptsize, at={(0.5,-0.20)}, anchor=north,
                legend columns=2, draw=none, fill=none, column sep=6pt},
]
\nextgroupplot[
  ybar, bar width=4pt,
  ymin=-9.35, ymax=29.47,
  xtick={0,1,2,3,4,5.7},
  xticklabels={2021,2022,2023,2024,2025,2026H1},
  ylabel={亿元},
]
\addplot[fill=AgriGreen!75, draw=AgriGreen, bar shift=-2.6pt] table[x=i, y=rev]{data/clean/profiles/fin_300511.dat};
\addplot[fill=SoftRed!60, draw=SoftRed, bar shift=2.6pt] table[x=i, y=ni]{data/clean/profiles/fin_300511.dat};
\addplot[fill=AgriGreen!30, draw=AgriGreen, pattern=north east lines, bar shift=-2.6pt] table[x=x, y=rev]{data/clean/profiles/finh_300511.dat};
\addplot[fill=SoftRed!20, draw=SoftRed, pattern=north east lines, bar shift=2.6pt] table[x=x, y=ni]{data/clean/profiles/finh_300511.dat};
\legend{营业收入（年/半年）, 归母净利润（年/半年）}
\nextgroupplot[
  ymin=-59.8, ymax=91.5,
  xtick={0,1,2,3,4,5.7},
  xticklabels={2021,2022,2023,2024,2025,2026H1},
  ylabel={\%},
]
\addplot[AgriGold, mark=*, mark size=1.3pt] table[x=i, y=roe]{data/clean/profiles/fin_300511.dat};
\addplot[OilBlue, mark=square*, mark size=1.3pt] table[x=i, y=debt]{data/clean/profiles/fin_300511.dat};
\addplot[only marks, AgriGold, mark=*, mark size=2.0pt] table[x=x, y=roe]{data/clean/profiles/finh_300511.dat};
\addplot[only marks, OilBlue, mark=square*, mark size=2.0pt] table[x=x, y=debt]{data/clean/profiles/finh_300511.dat};
\legend{ROE, 资产负债率}
\end{groupplot}
\end{tikzpicture}
\begin{tikzpicture}
\begin{axis}[
  name=finbot,
  width=13.75cm, height=3.10cm, scale only axis,
  ymin=-2.24, ymax=25.06,
  xtick={0,1,2,3,4,5.7},
  xticklabels={2021,2022,2023,2024,2025,2026H1},
  tick label style={font=\tiny},
  yticklabel style={font=\tiny},
  ylabel={亿元}, ylabel style={font=\tiny},
  ybar, bar width=4pt,
  grid=major, grid style={LightGray, line width=0.3pt},
  axis line style={InkGray, line width=0.5pt},
  enlarge x limits={abs=0.8},
  legend style={font=\scriptsize, at={(0.5,-0.26)}, anchor=north,
                legend columns=2, draw=none, fill=none, column sep=10pt},
]
\addplot[fill=AgriGreen!55, draw=AgriGreen, bar shift=-3.0pt] table[x=i, y=cash]{data/clean/profiles/fin_300511.dat};
\addplot[fill=SoftRed!45, draw=SoftRed, bar shift=3.0pt] table[x=i, y=ibd]{data/clean/profiles/fin_300511.dat};
\addplot[fill=AgriGreen!25, draw=AgriGreen, pattern=north east lines, bar shift=-3.0pt] table[x=x, y=cash]{data/clean/profiles/finh_300511.dat};
\addplot[fill=SoftRed!20, draw=SoftRed, pattern=north east lines, bar shift=3.0pt] table[x=x, y=ibd]{data/clean/profiles/finh_300511.dat};
\legend{货币资金（左轴）, 有息负债（左轴）}
\end{axis}
\begin{axis}[
  at={(finbot.south east)}, anchor=south east,
  width=13.75cm, height=3.10cm, scale only axis,
  axis y line*=right, axis x line*=none, xtick=\empty,
  ymin=-1.14, ymax=5.68,
  tick label style={font=\tiny},
  yticklabel style={font=\tiny}, scaled y ticks=false,
  legend style={font=\scriptsize, at={(0.5,-0.44)}, anchor=north,
                legend columns=1, draw=none, fill=none},
]
\addplot[OilBlue, mark=*, mark size=2.2pt, line width=1.3pt] table[x=i, y=ocf]{data/clean/profiles/fin_300511.dat};
\addplot[only marks, OilBlue, mark=*, mark size=3.2pt] table[x=x, y=ocf]{data/clean/profiles/finh_300511.dat};
\legend{经营活动现金流净额（右轴，亿元，蓝点）}
\end{axis}
\end{tikzpicture}

\captionof{figure}[雪榕生物（300511）近年主要财务指标]{雪榕生物（300511）近年主要财务指标（左上：营业收入与归母净利润；右上：ROE 与资产负债率；下：货币资金、有息负债为左轴柱状，经营活动现金流净额为其右轴的蓝点折线。
2021---2025 年为年报口径，2026H1 为 2026 年半年报/半年末，与其前一年报之间留出间隔、以斜纹或空心点区分；半年报数据未年化）}
\end{minipage}
\end{center}

\pfl{财务分析} 2026 年 6 月末资产负债率 56.5\%，较 2025 年末下降 1.4 个百分点（样本中位数 52.2\%，所处三级行业内按由低到高排第\zhnumber{4}位）。 2026 年 6 月末货币资金 1.76 亿元、短期债务（短期借款与一年内到期非流动负债）10.20 亿元、长期债务 2.32 亿元，有息负债合计 12.52 亿元，现金短债比 0.17 倍——缺口明显。 净有息负债 10.76 亿元，相当于其 2026 年上半年营业收入的 93\%。 流动比率 0.37、速动比率 0.15，短期偿付压力较大。 2026 年上半年经营活动现金流净额 2.21 亿元，经营现金流与利润同向为正，内生资金可以支撑运营。 2026 年 6 月末在建工程 0.70 亿元，较上年末增加 0.50 亿元，仍有在建投入。 综合看，账面货币资金对有息负债与短期债务的覆盖不足，滚动偿付对新增融资的依赖度较高；资产负债率高于样本中位数 4.4 个百分点；有息负债中短期债务占比更高，期限结构仍以滚动续作为主。 公司披露的股东与质押风险为：2026 年 9 月，控股股东上海万紫千鸿智能科技有限公司质押公司 1，900 万股。

\pfl{重大战略转型} 近三年公司的战略动作集中在：产能与项目（2 项）：2025 年 12 月，公司实现扭亏为盈；2024 年，主动收缩整体业务规模以应对行业产能过剩与价格低迷；融资与资本运作（1 项）：2025 年 11 月，万紫千鸿取得控制权的目的被披露为基于对上市公司内在价值和未来发展前景的信心。 公告口径上，近三年（2023 年 9 月---2026 年 9 月）公司披露重大重组与并购类公告 0 条、对外投资与转型类公告 5 条、回购与股权激励类公告 17 条，合计公告 516 条。 第一大行业仍是“食用菌行业”，收入占比 98.9\%，与 2021 年年报相比没有方向性变化。 综合判断，报告期内公司有 5 条与战略相关的公告落地（重大重组与并购 0 条、对外投资与转型 5 条），资本运作与主业调整之间存在直接关联。

\begin{center}\small\setlength{\tabcolsep}{4pt}
\captionof{table}[雪榕生物（300511）历史融资与资本运作]{雪榕生物（300511）历史融资与资本运作（据公司公告与公开报道整理；金额为披露值，含首次公开发行、再融资、债券、银行借款与重整投资等）}
\begin{adjustbox}{max width=\textwidth}
\begin{tabular}{@{}p{1.45cm}p{2.55cm}p{4.05cm}p{4.95cm}@{}}
\toprule
时间 & 方式 & 规模 & 用途与说明 \\
\midrule
2016-05 & IPO（深圳证券交易所创业板） & 发行 3，\allowbreak{}750 万股，\allowbreak{}发行价 16.82 元/\allowbreak{}股，\allowbreak{}募集资金总额 6.31 亿元 & 承销保荐费用 4,\allowbreak{}084.50 万元 \\
2020-06-24 & 公开发行可转换公司债券“雪榕转债”（债券代码 123056） & 85.00 亿元（585 万张，\allowbreak{}每张面值 100.00 元），\allowbreak{}发行信用等级 AA- & 含泰国食用菌工厂化生产车间建设项目等及补充流动资金；2020-10-13 公告变更部分募集资金用途，\allowbreak{}初始转股价 11.89 元/\allowbreak{}股 \\
2025-03 & 向特定对象发行股票（定增，预案） & 发行价格 3.00 元/\allowbreak{}股；若“雪榕转债”未转股 & 扣除发行费用后全部用于补充流动资金；发行对象为控股股东万紫千鸿（现金全额认购） \\
2025-09-03 & 可转债提前赎回 & 赎回价格 100.93 元/\allowbreak{}张（含税）；截至赎回登记日（2025-10-14）收市后 5 & 终止“雪榕转债”存续 \\
2026-03 & 2026 年限制性股票激励计划（第二类限制性股票） & 拟授予 2，\allowbreak{}704 万股，\allowbreak{}占当时总股本比例 4.22\%，\allowbreak{}每股转让价 6.26 元 & 向激励对象定向发行 A 股普通股股票；2026-03-30 披露向激励对象授予的公告；股权激励计划草案 2026-02-11 披露 \\
\bottomrule
\end{tabular}
\end{adjustbox}
\end{center}

\pfl{历史融投资情况} 从现金流量表看，2025 年公司取得借款收到的现金 4.54 亿元、偿还债务支付的现金 7.79 亿元、吸收权益性投资 0.01 亿元，筹资活动现金流净额 -3.58 亿元（2023 年为取得借款 12.01 亿元、偿还债务 8.96 亿元）。 2026 年上半年筹资活动现金流净额为 -0.90 亿元（取得借款 2.43 亿元、偿还债务 2.81 亿元，未年化），已转为净流出，处于净还债状态。 2025 年分配股利、利润或偿付利息支付的现金合计 0.43 亿元。 公开披露的融资记录共 5 笔，工具类型以 IPO、可转债、定向增发为主；金额与用途见下表。 其中规模最大的两笔为：2025 年 3 月，定向增发，发行价格 3.00 元/股，若“雪榕转债”未转股；2020 年 6 月，可转债，85.00 亿元（585 万张，每张面值 100.00 元），发行信用等级 AA-。 股本方面，近三年披露股本变动 13 次（可转债转股 9 次、其他 1 次、股份回购 1 次）；总股本由 2021 年末的 4.42 亿股变为 4.99 亿股（+12.87\%）。 分红方面，近三年未实施现金分红，在全部样本中属于分红能力偏弱的一端。 近三年“再融资”类公告 29 条，涉及定向增发、可转债、募集资金使用。 综合看，公司融资结构上以债务滚动为主、股权融资为补充；2025 年筹资活动现金流净额 -3.58 亿元，融资节奏仍以维持存量债务周转为主，与前述经营与投资现金流的方向基本一致。

\pfl{未来潜在融资需求分析} 2026 年 6 月末货币资金 1.76 亿元，而短期债务 10.20 亿元（现金短债比 0.17 倍），2026 年上半年经营现金流净额 2.21 亿元。按“短期债务减去账面货币资金”的滚动偿付口径，静态缺口约 8.43 亿元；上半年盈利或接近盈亏平衡，该缺口不随经营亏损进一步扩大。 公司 2026 年 6 月末资产负债率 56.5\%、2026 年上半年 ROE 1.1\%（未年化），现金短债比 0.17 倍、短债覆盖不足，融资需求首先用于置换与滚动存量债务、补充营运资金，属于“补血型”而非扩张型。 公开披露的后续资金安排线索包括：股权融资线索：2025 年度，公司向特定对象发行股票预案：拟向控股股东万紫千鸿发行不超过 1.50 亿股（雪榕转债未转股情形）、发行价 3.00 元/股、募集资金总额不超过 4.49 亿元。 资金需求还要叠加在建工程：期末在建工程 0.70 亿元、2026 年上半年购建固定资产等支出 1.40 亿元（未年化），这部分支出在周期底部通常会被主动放缓。 综合看，公司的外部资金需求以补充流动性与项目投入为主，滚动偿付口径下的静态缺口约 8.43 亿元，能否落地取决于授信续作与发行窗口的配合。




\subsubsection{华绿生物（300970）}
\label{co:300970}

\pfl{公司基本情况} 公司前身可追溯至 2018 年，江苏泗阳的食用菌工厂化企业，以金针菇、真姬菇为主并拓展鹿茸菇等珍稀品种；2021 年创业板上市，2024 年由盈转亏、2025 年随行业回暖大幅回升。 公司于 2021-04-12 在深交所创业板首发上市，注册资本 1.60 亿元，总股本 1.18 亿股。 截至 2026 年 9 月 24 日收盘价 17.62 元，流通市值 21.8 亿元，近一年+25.8\%，流通市值在三级行业“食用菌”的 4 家公司中按由大到小排第\zhnumber{3}位，近一年涨跌幅高于全样本中位数（-10.1\%）35.9 个百分点。 按 2025 年年报披露的行业构成，收入几乎全部来自食用菌 100.0\%。

\pfl{历史股价} 自 2021 年 4 月至 2026 年 9 月，股价由 57.00 元变为 17.62 元，累计 -69.1\%，区间最高 63.31 元（2021 年 5 月）、最低 11.23 元（2024 年 8 月）；当前价相当于区间最高价的 27.8\%，为区间最低价的 1.57 倍。 月度收益的年化波动率 65.5\%，高于全样本中位数 37.4\%，在三级行业“食用菌”的 4 家公司中波动率由高到低排第\zhnumber{1}位。

\begin{center}
\begin{minipage}{\textwidth}\centering
\begin{tikzpicture}
\begin{axis}[
  width=12.6cm, height=3.9cm, scale only axis,
  xmin=2021.172, xmax=2026.888, ymin=10.44, ymax=67.74,
  xtick={2021,2022,2023,2024,2025,2026},
  yearticks,
  yticklabel style={font=\scriptsize},
  ylabel={元/股}, ylabel style={font=\scriptsize},
  grid=major, grid style={LightGray, line width=0.3pt},
  axis line style={InkGray, line width=0.5pt},
  every axis plot/.append style={line width=0.9pt},
]
\addplot[AgriGreen] table[x=x,y=y]{data/clean/profiles/px_300970.dat};
\addplot[SoftRed, dashed, line width=0.7pt] table[x=x,y=y]{data/clean/profiles/pxzz_300970.dat};
\addplot[only marks, mark=*, mark size=1.1pt, SoftRed] table[x=x,y=y]{data/clean/profiles/pxzz_300970.dat};
\end{axis}
\end{tikzpicture}

\captionof{figure}[华绿生物（300970）股价与周期骨架]{华绿生物（300970）月度收盘价（元/股）与 ZigZag 周期骨架（阈值 25\%，圆点为周期转折点）}
\end{minipage}
\end{center}

\pfl{过去周期分析} 以 25\% 的 ZigZag 阈值划分，样本区间内股价形成 3 个主升段与 4 个主跌段。 最大主升段出现在 2024 年 8 月 至 2026 年 4 月，20 个月+214.3\%；最大主跌段为 2021 年 4 月 至 2021 年 7 月，3 个月-63.9\%。 最近一次转折出现在 2026 年 7 月（下跌段自 2026 年 4 月起，3 个月-54.3\%），当前处于该段的延续之中。 公司股价较申万农林牧渔指数提前约 21 个月见底，是板块中较早定价周期反转的公司之一。 食用菌是工厂化生产的制造型农业（第 \ref{sec:soft} 章），产能投放节奏与金针菇等单品价格是盈利的主要变量，与养殖周期基本无关，公司 2026 年上半年实现盈利、半年 ROE 2.8\%（未年化），好于全样本中位数（0.75\%），下行周期对其报表的冲击小于同业。 综合区间形态看，该股单段涨跌幅度偏大、以趋势段为主（最大主升 +214.3\%、最大主跌 -63.9\%），年化波动率 65.5\% 与全样本中位数 37.4\% 相比波动明显大于样本中位数。

\pfl{盈利情况} 2026 年上半年营业收入 6.94 亿元（同比 +39.2\%），归母净利润 0.46 亿元，实现扭亏。 以年报为趋势对照，2023---2025 年营业收入由 9.96 亿元变为 12.92 亿元（两年年均复合增速 +13.9\%），归母净利润由 0.30 亿元变为 1.10 亿元。 按半年报口径，2026 年上半年的 ROE 为 2.8\%（未年化）、销售净利率 8.0\%、毛利率 15.4\%，ROE 在“食用菌”4 家公司中按数值由高到低排第\zhnumber{2}位。 同期全样本半年 ROE 中位数 0.75\%，公司与之相差 2.04 个百分点（略好于中位数公司）。 从公司披露的原因看，业绩变动主要来自：2024 年，由盈转亏：主要产品金针菇整体销售价格较上年同期下降较多、部分项目建成转固导致折旧增加、2022/2023 两期股权激励计划导致报告期累计计提期间费用约 2，000.00 万元；2026 年上半年，营业成本 5.87 亿元。 综合看，公司在报告期维持盈利（高于全样本半年 ROE 中位数 0.75\%），利润的可持续性取决于食用菌景气的延续与成本管控的执行。

\begin{center}\small\setlength{\tabcolsep}{3.4pt}
\captionof{table}[华绿生物（300970）关键财务指标]{华绿生物（300970）关键财务指标（合并报表口径；两个半年报期均未年化；现金短债比 $=$ 货币资金 $\div$ 短期债务）}
\begin{adjustbox}{max width=\textwidth}
\begin{tabular}{@{}lrrrrr@{}}
\toprule
指标 & 2023 年 & 2024 年 & 2025 年 & 2025 年半年报 & 2026 年半年报 \\
\midrule
营业收入（亿元） & 9.96 & 10.32 & 12.92 & 4.99 & 6.94 \\
归母净利润（亿元） & 0.30 & -0.48 & 1.10 & -0.54 & 0.46 \\
毛利率（\%） & 13.5 & 6.0 & 16.1 & -3.4 & 15.4 \\
ROE（\%） & 2.0 & -3.2 & 7.1 & -3.7 & 2.8 \\
资产负债率（\%） & 26.0 & 34.1 & 32.3 & 34.7 & 35.4 \\
经营活动现金流净额（亿元） & 1.67 & 1.38 & 3.45 & 0.47 & 1.65 \\
货币资金（亿元） & 4.10 & 3.61 & 5.90 & 3.52 & 5.89 \\
有息负债合计（亿元） & 3.33 & 5.66 & 5.72 & 5.72 & 7.00 \\
现金短债比（倍） & 6.42 & 7.39 & 8.30 & 13.68 & 19.04 \\
\bottomrule
\end{tabular}
\end{adjustbox}
\end{center}

\begin{center}
\begin{minipage}{\textwidth}\centering
\begin{tikzpicture}
\begin{groupplot}[
  group style={group size=2 by 1, horizontal sep=0.60cm},
  width=6.85cm, height=4.60cm,
  tick label style={font=\tiny},
  title style={font=\scriptsize\heiti},
  yticklabel style={font=\tiny},
  ylabel style={font=\tiny},
  grid=major, grid style={LightGray, line width=0.3pt},
  axis line style={InkGray, line width=0.5pt},
  enlarge x limits={abs=0.8},
  legend style={font=\scriptsize, at={(0.5,-0.20)}, anchor=north,
                legend columns=2, draw=none, fill=none, column sep=6pt},
]
\nextgroupplot[
  ybar, bar width=4pt,
  ymin=-1.82, ymax=14.53,
  xtick={0,1,2,3,4,5.7},
  xticklabels={2021,2022,2023,2024,2025,2026H1},
  ylabel={亿元},
]
\addplot[fill=AgriGreen!75, draw=AgriGreen, bar shift=-2.6pt] table[x=i, y=rev]{data/clean/profiles/fin_300970.dat};
\addplot[fill=SoftRed!60, draw=SoftRed, bar shift=2.6pt] table[x=i, y=ni]{data/clean/profiles/fin_300970.dat};
\addplot[fill=AgriGreen!30, draw=AgriGreen, pattern=north east lines, bar shift=-2.6pt] table[x=x, y=rev]{data/clean/profiles/finh_300970.dat};
\addplot[fill=SoftRed!20, draw=SoftRed, pattern=north east lines, bar shift=2.6pt] table[x=x, y=ni]{data/clean/profiles/finh_300970.dat};
\legend{营业收入（年/半年）, 归母净利润（年/半年）}
\nextgroupplot[
  ymin=-6.3, ymax=39.3,
  xtick={0,1,2,3,4,5.7},
  xticklabels={2021,2022,2023,2024,2025,2026H1},
  ylabel={\%},
]
\addplot[AgriGold, mark=*, mark size=1.3pt] table[x=i, y=roe]{data/clean/profiles/fin_300970.dat};
\addplot[OilBlue, mark=square*, mark size=1.3pt] table[x=i, y=debt]{data/clean/profiles/fin_300970.dat};
\addplot[only marks, AgriGold, mark=*, mark size=2.0pt] table[x=x, y=roe]{data/clean/profiles/finh_300970.dat};
\addplot[only marks, OilBlue, mark=square*, mark size=2.0pt] table[x=x, y=debt]{data/clean/profiles/finh_300970.dat};
\legend{ROE, 资产负债率}
\end{groupplot}
\end{tikzpicture}
\begin{tikzpicture}
\begin{axis}[
  name=finbot,
  width=13.75cm, height=3.10cm, scale only axis,
  ymin=-0.70, ymax=7.84,
  xtick={0,1,2,3,4,5.7},
  xticklabels={2021,2022,2023,2024,2025,2026H1},
  tick label style={font=\tiny},
  yticklabel style={font=\tiny},
  ylabel={亿元}, ylabel style={font=\tiny},
  ybar, bar width=4pt,
  grid=major, grid style={LightGray, line width=0.3pt},
  axis line style={InkGray, line width=0.5pt},
  enlarge x limits={abs=0.8},
  legend style={font=\scriptsize, at={(0.5,-0.26)}, anchor=north,
                legend columns=2, draw=none, fill=none, column sep=10pt},
]
\addplot[fill=AgriGreen!55, draw=AgriGreen, bar shift=-3.0pt] table[x=i, y=cash]{data/clean/profiles/fin_300970.dat};
\addplot[fill=SoftRed!45, draw=SoftRed, bar shift=3.0pt] table[x=i, y=ibd]{data/clean/profiles/fin_300970.dat};
\addplot[fill=AgriGreen!25, draw=AgriGreen, pattern=north east lines, bar shift=-3.0pt] table[x=x, y=cash]{data/clean/profiles/finh_300970.dat};
\addplot[fill=SoftRed!20, draw=SoftRed, pattern=north east lines, bar shift=3.0pt] table[x=x, y=ibd]{data/clean/profiles/finh_300970.dat};
\legend{货币资金（左轴）, 有息负债（左轴）}
\end{axis}
\begin{axis}[
  at={(finbot.south east)}, anchor=south east,
  width=13.75cm, height=3.10cm, scale only axis,
  axis y line*=right, axis x line*=none, xtick=\empty,
  ymin=-0.90, ymax=4.49,
  tick label style={font=\tiny},
  yticklabel style={font=\tiny}, scaled y ticks=false,
  legend style={font=\scriptsize, at={(0.5,-0.44)}, anchor=north,
                legend columns=1, draw=none, fill=none},
]
\addplot[OilBlue, mark=*, mark size=2.2pt, line width=1.3pt] table[x=i, y=ocf]{data/clean/profiles/fin_300970.dat};
\addplot[only marks, OilBlue, mark=*, mark size=3.2pt] table[x=x, y=ocf]{data/clean/profiles/finh_300970.dat};
\legend{经营活动现金流净额（右轴，亿元，蓝点）}
\end{axis}
\end{tikzpicture}

\captionof{figure}[华绿生物（300970）近年主要财务指标]{华绿生物（300970）近年主要财务指标（左上：营业收入与归母净利润；右上：ROE 与资产负债率；下：货币资金、有息负债为左轴柱状，经营活动现金流净额为其右轴的蓝点折线。
2021---2025 年为年报口径，2026H1 为 2026 年半年报/半年末，与其前一年报之间留出间隔、以斜纹或空心点区分；半年报数据未年化）}
\end{minipage}
\end{center}

\pfl{财务分析} 2026 年 6 月末资产负债率 35.4\%，较 2025 年末上升 3.1 个百分点（样本中位数 52.2\%，所处三级行业内按由低到高排第\zhnumber{2}位）。 2026 年 6 月末货币资金 5.89 亿元、短期债务（短期借款与一年内到期非流动负债）0.31 亿元、长期债务 6.69 亿元，有息负债合计 7.00 亿元，现金短债比 19.04 倍——覆盖充分。 净有息负债 1.11 亿元，相当于其 2026 年上半年营业收入的 16\%。 流动比率 4.21、速动比率 2.99，短期指标尚可。 2026 年上半年经营活动现金流净额 1.65 亿元，经营现金流与利润同向为正，内生资金可以支撑运营。 综合看，现金短债比 19.04 倍，短期偿付压力可控；资产负债率低于样本中位数 16.8 个百分点；有息负债中长期债务占比更高，期限结构对短债滚动的压力较小。

\pfl{重大战略转型} 从公开披露看，公司的战略推进集中在：产能与项目（3 项）：2026 年 4 月，媒体披露公司年内合计投资 5.00 亿元加码高端产能；2026 年 3 月，公司拟与文山市人民政府签订投资协议，在云南省文山壮族苗族自治州文山市新建日产 50 吨鲜品黄牛肝菌（见手青）工厂化生产项目。 公告口径上，近三年（2023 年 9 月---2026 年 9 月）公司披露重大重组与并购类公告 1 条、对外投资与转型类公告 5 条、回购与股权激励类公告 69 条，合计公告 311 条。 第一大行业仍是“食用菌”，收入占比 100.0\%，与 2021 年年报相比没有方向性变化。 综合判断，报告期内公司有 6 条与战略相关的公告落地（重大重组与并购 1 条、对外投资与转型 5 条），资本运作与主业调整之间存在直接关联。

\begin{center}\small\setlength{\tabcolsep}{4pt}
\captionof{table}[华绿生物（300970）历史融资与资本运作]{华绿生物（300970）历史融资与资本运作（据公司公告与公开报道整理；金额为披露值，含首次公开发行、再融资、债券、银行借款与重整投资等）}
\begin{adjustbox}{max width=\textwidth}
\begin{tabular}{@{}p{1.45cm}p{2.55cm}p{4.05cm}p{4.95cm}@{}}
\toprule
时间 & 方式 & 规模 & 用途与说明 \\
\midrule
2021-04 & IPO（深圳证券交易所创业板） & 发行 1，\allowbreak{}459.00 万股，\allowbreak{}发行价 44.77 元/\allowbreak{}股，\allowbreak{}募集资金总额 6.53 亿元 & 江苏华绿生物科技股份有限公司一厂技术改造项目、\allowbreak{}泗阳华茂农业发展有限公司技术改造项目、\allowbreak{}年产 3 万吨真姬菇工厂化栽培建设项目；；计入资本公积（股本溢价）5.82 亿元 \\
2022-09-30 & 终止募投项目并将剩余募集资金永久补充流动资金 & 剩余募集资金 2.87 亿元及利息 & 终止一厂技术改造项目、\allowbreak{}泗阳华茂农业发展有限公司技术改造项目 \\
2022/\allowbreak{}2023 & 两期股权激励计划（第二类限制性股票，股票来源为向激励对象定向增发） & 2022 年计划以 10.68 元/\allowbreak{}股向 58 名激励对象授予 400.00 万股（授予日 2022-09-30）；2023 年计划以 11.46 元/\allowbreak{}股向 41 名激励对象授予 329.80 万股、\allowbreak{}预留 70.20 万股（首次授予日 2023-08-21） & 员工激励；2023-10-31 第一期 118.86 万股归属上市；2024-10-22 第二期 142.52 万股归属上市 \\
2023-11-30 & 募投项目结项并将节余募集资金永久补充流动资金 & 节余募集资金 856.19 万元（截至 2023-11-30） & 河北华绿之珍年产 5.4 万吨鲜品金针菇工厂化生产项目结项 \\
2024-03-04 & 股份回购 & 计划使用自有资金不低于 3，\allowbreak{}000.00 万元、\allowbreak{}不超过 6 & 在适宜时机用于实施股权激励计划或员工持股计划；2024 年年度报告列示回购股份金额 3，\allowbreak{}001.27 万元；2025 年年度报告以现有总股本 1.23 亿股剔除回购专用账户 234.66 万股为分红基数 \\
2025-04-22 & 2025 年度利润分配预案（含转增） & 以 1.20 亿股为基数，\allowbreak{}每 10 股派发现金红利 5.00 元（含税）、\allowbreak{}送红股 0 股 & 股东回报；2026-04-22 董事会审议通过（公司现有总股本 1.23 亿股 \\
\bottomrule
\end{tabular}
\end{adjustbox}
\end{center}

\pfl{历史融投资情况} 从现金流量表看，2025 年公司取得借款收到的现金 1.49 亿元、偿还债务支付的现金 1.51 亿元、吸收权益性投资 0.26 亿元，筹资活动现金流净额 0.07 亿元（2023 年为取得借款 0.75 亿元、偿还债务 1.30 亿元）。 2026 年上半年筹资活动现金流净额为 -1.16 亿元（取得借款 0.09 亿元、偿还债务 0.44 亿元，未年化），已转为净流出，处于净还债状态。 2025 年分配股利、利润或偿付利息支付的现金合计 0.13 亿元。 公开披露的融资记录共 6 笔，工具类型以 IPO 为主；金额与用途见下表。 其中规模最大的两笔为：2023 年 11 月，募投项目结项并将节余，节余募集资金 856.19 万元（截至 2023-11-30）；2021 年 4 月，IPO，发行 1，459.00 万股，发行价 44.77 元/股，募集资金总额 6.53 亿元。 股本方面，近三年披露股本变动 5 次（增发新股上市 1 次、股权激励 1 次、限售股份上市 1 次）；总股本由 2021 年末的 1.17 亿股变为 1.18 亿股（+1.02\%）。 分红方面，近三年现金分红 3 次、合计每股派现 0.90 元，最近一次实施在 2025 年 12 月。 综合看，公司融资结构上股权与债务融资并举；2025 年筹资活动现金流净额 0.07 亿元，年度内仍为净融入，与前述经营与投资现金流的方向基本一致。

\pfl{未来潜在融资需求分析} 2026 年 6 月末货币资金 5.89 亿元，而短期债务 0.31 亿元（现金短债比 19.04 倍），2026 年上半年经营现金流净额 1.65 亿元。账面货币资金可以覆盖短期债务，缺口不体现在静态偿付层面。 公司 2026 年 6 月末资产负债率 35.4\%、2026 年上半年 ROE 2.8\%（未年化），资产端与负债端都没有明显压力，融资需求以相机抉择的扩张型用途为主。 公开披露的后续资金安排线索包括：债务与授信安排：2026 年 4 月，公司对子公司担保：为南川华绿提供担保额度 1.40 亿元；2026 年 6 月，在建项目资金来源披露为金融机构贷款：南川鹿茸菇项目预算 5，004.54 万元、工程进度 95.00\%。 综合看，公司在静态偿付层面不存在缺口，外部融资更多是配合扩张节奏与并购整合的主动性安排，而非偿债压力驱动。




\subsubsection{博闻科技（600883）}
\label{co:600883}

\pfl{公司基本情况} 博闻科技成立于 1990 年，云南保山起家的老牌上市公司，主业历经水泥—科技—农副食品的多次更迭，现主营云南高原特色农副食品加工（食用菌、火腿、咖啡），并持有新疆众和参股权益与部分证券资产。 公司于 1995-12-08 在上交所首发上市，注册资本 2.36 亿元，总股本 2.36 亿股。 截至 2026 年 9 月 24 日收盘价 8.11 元，流通市值 19.1 亿元，近一年+2.7\%，流通市值在三级行业“食用菌”的 4 家公司中按由大到小排第\zhnumber{4}位，近一年涨跌幅高于全样本中位数（-10.1\%）12.7 个百分点。 按 2025 年年报披露的行业构成，收入占比前三位为食用菌 73.5\%；咖啡 18.0\%；火腿及制品 8.1\%。

\pfl{历史股价} 本报告样本区间（2021 年 1 月 至 2026 年 9 月）内，股价由 7.74 元变为 8.11 元，累计 +4.8\%，区间最高 9.33 元（2022 年 1 月）、最低 5.54 元（2024 年 8 月）；当前价相当于区间最高价的 86.9\%，为区间最低价的 1.46 倍。 月度收益的年化波动率 25.8\%，低于全样本中位数 37.4\%，在三级行业“食用菌”的 4 家公司中波动率由高到低排第\zhnumber{4}位。 把股价阶段与公司事件对齐看，2024 年 8 月 至 2026 年 4 月股价上涨+62.6\%，同期（2025 年 12 月）公司账面同时持有云南白药股份 67.34 万股，并减持众和转债回笼资金。 从时间对应关系看，公司层面的资本运作与股权、风险事件解释了区间内多数急涨急跌，与行业指数的同步性相对有限。

\begin{center}
\begin{minipage}{\textwidth}\centering
\begin{tikzpicture}
\begin{axis}[
  width=12.6cm, height=3.9cm, scale only axis,
  xmin=2020.922, xmax=2026.888, ymin=5.15, ymax=9.98,
  xtick={2021,2022,2023,2024,2025,2026},
  yearticks,
  yticklabel style={font=\scriptsize},
  ylabel={元/股}, ylabel style={font=\scriptsize},
  grid=major, grid style={LightGray, line width=0.3pt},
  axis line style={InkGray, line width=0.5pt},
  every axis plot/.append style={line width=0.9pt},
]
\addplot[AgriGreen] table[x=x,y=y]{data/clean/profiles/px_600883.dat};
\addplot[SoftRed, dashed, line width=0.7pt] table[x=x,y=y]{data/clean/profiles/pxzz_600883.dat};
\addplot[only marks, mark=*, mark size=1.1pt, SoftRed] table[x=x,y=y]{data/clean/profiles/pxzz_600883.dat};
\end{axis}
\end{tikzpicture}

\captionof{figure}[博闻科技（600883）股价与周期骨架]{博闻科技（600883）月度收盘价（元/股）与 ZigZag 周期骨架（阈值 25\%，圆点为周期转折点）}
\end{minipage}
\end{center}

\pfl{过去周期分析} 以 25\% 的 ZigZag 阈值划分，样本区间内股价形成 2 个主升段与 1 个主跌段。 最大主跌段为 2022 年 1 月 至 2024 年 8 月，31 个月-40.6\%；最大主升段出现在 2024 年 8 月 至 2026 年 4 月，20 个月+62.6\%。 最近一次转折出现在 2026 年 4 月（上涨段自 2024 年 8 月起，20 个月+62.6\%），当前处于该段之后的震荡之中。 公司股价较申万农林牧渔指数提前约 21 个月见底，是板块中较早定价周期反转的公司之一。 食用菌是工厂化生产的制造型农业（第 \ref{sec:soft} 章），产能投放节奏与金针菇等单品价格是盈利的主要变量，与养殖周期基本无关，公司 2026 年上半年实现盈利、半年 ROE 1.5\%（未年化），好于全样本中位数（0.75\%），下行周期对其报表的冲击小于同业。 综合区间形态看，该股单段涨跌幅度偏大、以趋势段为主（最大主升 +62.6\%、最大主跌 -40.6\%），年化波动率 25.8\% 与全样本中位数 37.4\% 相比波动小于样本中位数。

\pfl{盈利情况} 2026 年上半年营业收入 0.26 亿元（同比 -8.8\%），归母净利润 0.15 亿元，同比 +5.3\%。 以年报为趋势对照，2023---2025 年营业收入由 0.21 亿元变为 0.56 亿元（两年年均复合增速 +62.2\%），归母净利润由 0.96 亿元变为 0.24 亿元。 按半年报口径，2026 年上半年的 ROE 为 1.5\%（未年化）、销售净利率 56.9\%、毛利率 5.4\%，ROE 在“食用菌”4 家公司中按数值由高到低排第\zhnumber{3}位。 同期全样本半年 ROE 中位数 0.75\%，公司与之相差 0.77 个百分点（略好于中位数公司）。 综合看，公司在报告期维持盈利（高于全样本半年 ROE 中位数 0.75\%），利润的可持续性取决于食用菌景气的延续与成本管控的执行。

\begin{center}\small\setlength{\tabcolsep}{3.4pt}
\captionof{table}[博闻科技（600883）关键财务指标]{博闻科技（600883）关键财务指标（合并报表口径；两个半年报期均未年化；现金短债比 $=$ 货币资金 $\div$ 短期债务）}
\begin{adjustbox}{max width=\textwidth}
\begin{tabular}{@{}lrrrrr@{}}
\toprule
指标 & 2023 年 & 2024 年 & 2025 年 & 2025 年半年报 & 2026 年半年报 \\
\midrule
营业收入（亿元） & 0.21 & 0.43 & 0.56 & 0.29 & 0.26 \\
归母净利润（亿元） & 0.96 & 0.72 & 0.24 & 0.14 & 0.15 \\
毛利率（\%） & -1.0 & 0.9 & 8.0 & 5.6 & 5.4 \\
ROE（\%） & 11.2 & 7.7 & 2.4 & 1.5 & 1.5 \\
资产负债率（\%） & 5.6 & 5.7 & 5.9 & 5.8 & 6.6 \\
经营活动现金流净额（亿元） & -0.33 & -0.40 & -0.11 & -0.05 & -0.05 \\
货币资金（亿元） & 1.07 & 0.35 & 1.63 & 1.07 & 1.43 \\
有息负债合计（亿元） & 0.04 & 0.04 & 0.05 & 0.04 & 0.12 \\
现金短债比（倍） & 30.98 & 8.33 & 35.25 & 25.76 & 30.40 \\
\bottomrule
\end{tabular}
\end{adjustbox}
\end{center}

\begin{center}
\begin{minipage}{\textwidth}\centering
\begin{tikzpicture}
\begin{groupplot}[
  group style={group size=2 by 1, horizontal sep=0.60cm},
  width=6.85cm, height=4.60cm,
  tick label style={font=\tiny},
  title style={font=\scriptsize\heiti},
  yticklabel style={font=\tiny},
  ylabel style={font=\tiny},
  grid=major, grid style={LightGray, line width=0.3pt},
  axis line style={InkGray, line width=0.5pt},
  enlarge x limits={abs=0.8},
  legend style={font=\scriptsize, at={(0.5,-0.20)}, anchor=north,
                legend columns=2, draw=none, fill=none, column sep=6pt},
]
\nextgroupplot[
  ybar, bar width=4pt,
  ymin=-0.10, ymax=1.08,
  xtick={0,1,2,3,4,5.7},
  xticklabels={2021,2022,2023,2024,2025,2026H1},
  ylabel={亿元},
]
\addplot[fill=AgriGreen!75, draw=AgriGreen, bar shift=-2.6pt] table[x=i, y=rev]{data/clean/profiles/fin_600883.dat};
\addplot[fill=SoftRed!60, draw=SoftRed, bar shift=2.6pt] table[x=i, y=ni]{data/clean/profiles/fin_600883.dat};
\addplot[fill=AgriGreen!30, draw=AgriGreen, pattern=north east lines, bar shift=-2.6pt] table[x=x, y=rev]{data/clean/profiles/finh_600883.dat};
\addplot[fill=SoftRed!20, draw=SoftRed, pattern=north east lines, bar shift=2.6pt] table[x=x, y=ni]{data/clean/profiles/finh_600883.dat};
\legend{营业收入（年/半年）, 归母净利润（年/半年）}
\nextgroupplot[
  ymin=-0.9, ymax=12.3,
  xtick={0,1,2,3,4,5.7},
  xticklabels={2021,2022,2023,2024,2025,2026H1},
  ylabel={\%},
]
\addplot[AgriGold, mark=*, mark size=1.3pt] table[x=i, y=roe]{data/clean/profiles/fin_600883.dat};
\addplot[OilBlue, mark=square*, mark size=1.3pt] table[x=i, y=debt]{data/clean/profiles/fin_600883.dat};
\addplot[only marks, AgriGold, mark=*, mark size=2.0pt] table[x=x, y=roe]{data/clean/profiles/finh_600883.dat};
\addplot[only marks, OilBlue, mark=square*, mark size=2.0pt] table[x=x, y=debt]{data/clean/profiles/finh_600883.dat};
\legend{ROE, 资产负债率}
\end{groupplot}
\end{tikzpicture}
\begin{tikzpicture}
\begin{axis}[
  name=finbot,
  width=13.75cm, height=3.10cm, scale only axis,
  ymin=-0.16, ymax=1.82,
  xtick={0,1,2,3,4,5.7},
  xticklabels={2021,2022,2023,2024,2025,2026H1},
  tick label style={font=\tiny},
  yticklabel style={font=\tiny},
  ylabel={亿元}, ylabel style={font=\tiny},
  ybar, bar width=4pt,
  grid=major, grid style={LightGray, line width=0.3pt},
  axis line style={InkGray, line width=0.5pt},
  enlarge x limits={abs=0.8},
  legend style={font=\scriptsize, at={(0.5,-0.26)}, anchor=north,
                legend columns=2, draw=none, fill=none, column sep=10pt},
]
\addplot[fill=AgriGreen!55, draw=AgriGreen, bar shift=-3.0pt] table[x=i, y=cash]{data/clean/profiles/fin_600883.dat};
\addplot[fill=SoftRed!45, draw=SoftRed, bar shift=3.0pt] table[x=i, y=ibd]{data/clean/profiles/fin_600883.dat};
\addplot[fill=AgriGreen!25, draw=AgriGreen, pattern=north east lines, bar shift=-3.0pt] table[x=x, y=cash]{data/clean/profiles/finh_600883.dat};
\addplot[fill=SoftRed!20, draw=SoftRed, pattern=north east lines, bar shift=3.0pt] table[x=x, y=ibd]{data/clean/profiles/finh_600883.dat};
\legend{货币资金（左轴）, 有息负债（左轴）}
\end{axis}
\begin{axis}[
  at={(finbot.south east)}, anchor=south east,
  width=13.75cm, height=3.10cm, scale only axis,
  axis y line*=right, axis x line*=none, xtick=\empty,
  ymin=-0.50, ymax=0.12,
  tick label style={font=\tiny},
  yticklabel style={font=\tiny}, scaled y ticks=false,
  legend style={font=\scriptsize, at={(0.5,-0.44)}, anchor=north,
                legend columns=1, draw=none, fill=none},
]
\addplot[OilBlue, mark=*, mark size=2.2pt, line width=1.3pt] table[x=i, y=ocf]{data/clean/profiles/fin_600883.dat};
\addplot[only marks, OilBlue, mark=*, mark size=3.2pt] table[x=x, y=ocf]{data/clean/profiles/finh_600883.dat};
\legend{经营活动现金流净额（右轴，亿元，蓝点）}
\end{axis}
\end{tikzpicture}

\captionof{figure}[博闻科技（600883）近年主要财务指标]{博闻科技（600883）近年主要财务指标（左上：营业收入与归母净利润；右上：ROE 与资产负债率；下：货币资金、有息负债为左轴柱状，经营活动现金流净额为其右轴的蓝点折线。
2021---2025 年为年报口径，2026H1 为 2026 年半年报/半年末，与其前一年报之间留出间隔、以斜纹或空心点区分；半年报数据未年化）}
\end{minipage}
\end{center}

\pfl{财务分析} 2026 年 6 月末资产负债率 6.6\%，较 2025 年末上升 0.7 个百分点（样本中位数 52.2\%，所处三级行业内按由低到高排第\zhnumber{1}位）。 2026 年 6 月末货币资金 1.43 亿元、短期债务（短期借款与一年内到期非流动负债）0.05 亿元、长期债务 0.07 亿元，有息负债合计 0.12 亿元，现金短债比 30.40 倍——覆盖充分。 净有息负债 -1.31 亿元，相当于其 2026 年上半年营业收入的 -502\%。 流动比率 5.89、速动比率 5.13，短期指标尚可。 2026 年上半年经营活动现金流净额 -0.05 亿元，净利润为正但经营现金流为负，利润的现金含量偏低。 综合看，现金短债比 30.40 倍，短期偿付压力可控；资产负债率低于样本中位数 45.6 个百分点；有息负债中长期债务占比更高，期限结构对短债滚动的压力较小。

\pfl{重大战略转型} 公司的战略动作可归为以下几类：并购与重组（1 项）：2024 年，持股 5\% 以上股东保山智源教育投资发展有限公司通过公开征集转让方式协议转让所持公司全部股份 1。 公告口径上，近三年（2023 年 9 月---2026 年 9 月）公司披露重大重组与并购类公告 0 条、对外投资与转型类公告 4 条、回购与股权激励类公告 0 条，合计公告 239 条。 第一大行业仍是“食用菌”，收入占比 73.5\%，与 2021 年年报相比没有方向性变化。 综合判断，报告期内公司有 4 条与战略相关的公告落地（重大重组与并购 0 条、对外投资与转型 4 条），资本运作与主业调整之间存在直接关联。

\begin{center}\small\setlength{\tabcolsep}{4pt}
\captionof{table}[博闻科技（600883）历史融资与资本运作]{博闻科技（600883）历史融资与资本运作（据公司公告与公开报道整理；金额为披露值，含首次公开发行、再融资、债券、银行借款与重整投资等）}
\begin{adjustbox}{max width=\textwidth}
\begin{tabular}{@{}p{1.45cm}p{2.55cm}p{4.05cm}p{4.95cm}@{}}
\toprule
时间 & 方式 & 规模 & 用途与说明 \\
\midrule
1995-12 & IPO（上海证券交易所挂牌） & 募集资金规模 & 挂牌时公司名为云南省保山水泥股份有限公司 \\
2025-03 & 自有资金出资设立咖啡业务子公司 & 以自有货币资金出资 2,\allowbreak{}000.00 万元人民币 & 在云南省怒江傈僳族自治州泸水市设立子公司，\allowbreak{}作为咖啡种植基地项目投资运营主体；年报同时披露怒江博闻咖啡有限公司实收资本 250.00 万元；截至 2025 年末 \\
2025-07-17 至 2025-08-08 & 减持所持新疆众和可转换公司债券（众和转债） & 减持 500，\allowbreak{}000 张，\allowbreak{}累计成交金额 6，\allowbreak{}491.87 万元，\allowbreak{}预计实现投资收益约 1 & 资产处置与资金回笼；本次减持完成后公司不再持有众和转债编号临 2025-023 \\
2026-05-22 & 2025 年度利润分配实施 & 按每股人民币 0.04 元（含税）、\allowbreak{}已发行股份 2.36 亿股计算，\allowbreak{}共计派发现金股利 8 & 股东回报；2025 年度股东大会批准；母公司 2025 年末可供分配利润 5.48 亿元 \\
\bottomrule
\end{tabular}
\end{adjustbox}
\end{center}

\pfl{历史融投资情况} 从现金流量表看，2025 年公司偿还债务支付的现金 0.00 亿元，筹资活动现金流净额 -0.12 亿元（2023 年为取得借款 0.03 亿元）。 2026 年上半年筹资活动现金流净额为 -0.07 亿元（取得借款 — 亿元、偿还债务 0.00 亿元，未年化），已转为净流出，处于净还债状态。 2025 年分配股利、利润或偿付利息支付的现金合计 0.11 亿元。 公开披露的融资记录共 4 笔，工具类型以 IPO、可转债为主；金额与用途见下表。 其中规模最大的两笔为：2025 年 7 月，可转债，减持 500，000 张，累计成交金额 6，491.87 万元，预计实现投资收益约 1；2025 年 3 月，自有资金出资设立咖啡，以自有货币资金出资 2,000.00 万元人民币。 股本方面，近三年披露股本变动 3 次；总股本自 2021 年末以来未发生变化（2.36 亿股）。 分红方面，近三年现金分红 4 次、合计每股派现 0.25 元，最近一次实施在 2025 年 12 月。 近三年“再融资”类公告 1 条，涉及可转债。 综合看，公司融资结构上股权与债务融资并举；2025 年筹资活动现金流净额 -0.12 亿元，融资节奏仍以维持存量债务周转为主，与前述经营与投资现金流的方向基本一致。

\pfl{未来潜在融资需求分析} 2026 年 6 月末货币资金 1.43 亿元，而短期债务 0.05 亿元（现金短债比 30.40 倍），2026 年上半年经营现金流净额 -0.05 亿元。账面货币资金可以覆盖短期债务，缺口不体现在静态偿付层面。 公司 2026 年 6 月末资产负债率 6.6\%、2026 年上半年 ROE 1.5\%（未年化），账面杠杆不高、盈利能力一般，短期内没有必须依靠外部融资才能维持的经营缺口；后续融资更可能指向技改、扩产或产业并购。 以现有货币资金与经营现金流衡量，公司不存在刚性补血的紧迫性，融资动作更多是主动型的。




